% Documento principal: estructura e inclusión de capítulos.
\documentclass[12pt,oneside]{report}
\usepackage[T1]{fontenc}
\usepackage[utf8]{inputenc}
\usepackage[spanish,es-nodecimaldot]{babel}
\usepackage{lmodern}

\usepackage{microtype}
\usepackage[a4paper,margin=2.8cm]{geometry}
\usepackage{setspace}
\onehalfspacing

\usepackage{amsmath,amssymb,amsfonts,mathtools}
\usepackage{bm}
\usepackage{siunitx}
\sisetup{detect-all}

\usepackage{graphicx}
\usepackage{booktabs}
\usepackage{tabularx}
\usepackage{longtable}
\usepackage{ragged2e}
\usepackage{caption}
\usepackage{float}
\usepackage{placeins}
\usepackage{adjustbox}
\usepackage{multirow}
\usepackage{array}
\newcolumntype{L}{>{\RaggedRight\arraybackslash}X}
\newcolumntype{C}{>{\Centering\arraybackslash}X}
\newcolumntype{P}[1]{>{\RaggedRight\arraybackslash}p{#1}}

\usepackage{enumitem}
\usepackage{csquotes}
\usepackage[table]{xcolor}
\usepackage[numbers,sort&compress]{natbib}
\usepackage[colorlinks=true,allcolors=blue]{hyperref}
\usepackage{xurl}
\urlstyle{tt}
\usepackage[spanish,nameinlink,noabbrev]{cleveref}
\crefname{chapter}{capítulo}{capítulos}
\Crefname{chapter}{Capítulo}{Capítulos}
\crefname{section}{sección}{secciones}
\Crefname{section}{Sección}{Secciones}
\crefname{subsection}{sección}{secciones}
\Crefname{subsection}{Sección}{Secciones}
\crefname{equation}{ecuación}{ecuaciones}
\Crefname{equation}{Ecuación}{Ecuaciones}
\crefname{table}{cuadro}{cuadros}
\Crefname{table}{Cuadro}{Cuadros}
\crefname{figure}{figura}{figuras}
\Crefname{figure}{Figura}{Figuras}
\usepackage{tikz}
\usepackage[titles]{tocloft}
% Espacio para números como 12.14.3 en el índice del manuscrito ampliado.
\setlength{\cftchapnumwidth}{2.3em}
\setlength{\cftsecindent}{2.3em}
\setlength{\cftsecnumwidth}{3.3em}
\setlength{\cftsubsecindent}{5.6em}
\setlength{\cftsubsecnumwidth}{4.2em}

\usetikzlibrary{arrows.meta,positioning,fit,shapes.geometric,babel}

\definecolor{azulacademico}{RGB}{45,73,104}
\definecolor{azulclaro}{RGB}{220,229,238}
\definecolor{griscronograma}{RGB}{225,225,225}

% ======================================================
%                   ACRÓNIMOS
% ======================================================
\usepackage[acronym,nogroupskip]{glossaries-extra}
\setabbreviationstyle[acronym]{long-short}
\makeglossaries
\newcommand{\ac}{\gls}
\newcommand{\acp}{\glspl}
\newcommand{\acs}{\glsentryshort}
\newcommand{\acf}{\glsfirst}

% RF / Comms / Sensing
\newacronym{RF}{RF}{radiofrecuencia}
\newacronym{WiFi}{Wi-Fi}{tecnología de redes inalámbricas basada en IEEE 802.11}
\newacronym{LoRa}{LoRa}{radio de largo alcance}
\newacronym{LoRaWAN}{LoRaWAN}{red de área amplia y largo alcance}
\newacronym{LPWAN}{LPWAN}{red de área amplia y bajo consumo}
\newacronym{OFDM}{OFDM}{multiplexación por división de frecuencias ortogonales}
\newacronym{CSI}{CSI}{información del estado del canal}
\newacronym{CFR}{CFR}{respuesta en frecuencia del canal}
\newacronym{CIR}{CIR}{respuesta al impulso del canal}
\newacronym{MIMO}{MIMO}{múltiples entradas y múltiples salidas}
\newacronym{SISO}{SISO}{una entrada y una salida}
\newacronym{RT}{RT}{trazado de rayos}
\newacronym{DRT}{DRT}{trazado de rayos diferenciable}
\newacronym{WDT}{WDT}{gemelo digital inalámbrico}
\newacronym{ISAC}{ISAC}{comunicación y sensado integrados}
\newacronym{SDR}{SDR}{radio definida por software}
\newacronym{IQ}{I/Q}{componentes en fase y cuadratura}
\newacronym{AoA}{AoA}{ángulo de llegada}
\newacronym{AoD}{AoD}{ángulo de salida}
\newacronym{ToF}{ToF}{tiempo de vuelo}
\newacronym{PDP}{PDP}{perfil de potencia en función del retardo}
\newacronym{RMSDS}{RMS-DS}{dispersión cuadrática media del retardo}
\newacronym{SNR}{SNR}{relación señal a ruido}
\newacronym{SSNR}{SSNR}{relación señal a ruido para sensado}
\newacronym{CFO}{CFO}{desajuste de frecuencia de portadora}
\newacronym{SFO}{SFO}{desajuste de frecuencia de muestreo}
\newacronym{STO}{STO}{desajuste temporal de muestreo}
\newacronym{CPO}{CPO}{desajuste de fase de portadora}
\newacronym{PEOC}{PEOC}{calibración que ignora el error de fase}
\newacronym{UPEC}{UPEC}{calibración con error de fase uniforme}
\newacronym{PEAC}{PEAC}{calibración que considera el error de fase}
\newacronym{PHY}{PHY}{capa física}
\newacronym{CSS}{CSS}{espectro ensanchado mediante chirridos}
\newacronym{FFT}{FFT}{transformada rápida de Fourier}
\newacronym{IFFT}{IFFT}{transformada rápida inversa de Fourier}
\newacronym{FIM}{FIM}{matriz de información de Fisher}
\newacronym{CRLB}{CRLB}{cota inferior de Cramér--Rao}
\newacronym{LOS}{LOS}{línea de vista}
\newacronym{NLOS}{NLOS}{sin línea de vista}
\newacronym{AWGN}{AWGN}{ruido gaussiano blanco aditivo}
\newacronym{LTI}{LTI}{lineal e invariante en el tiempo}
\newacronym{LTV}{LTV}{lineal y variante en el tiempo}
\newacronym{WSSUS}{WSSUS}{estacionario en sentido amplio con dispersión no correlacionada}
\newacronym{CP}{CP}{prefijo cíclico}
\newacronym{PSD}{PSD}{densidad espectral de potencia}

% ML / Representation Learning
\newacronym{DL}{DL}{aprendizaje profundo}
\newacronym{ML}{ML}{aprendizaje automático}
\newacronym{SSL}{SSL}{aprendizaje autosupervisado}
\newacronym{JEPA}{JEPA}{Joint-Embedding Predictive Architecture}
\newacronym{PINN}{PINN}{red neuronal informada por la física}
\newacronym{BSS}{BSS}{separación ciega de fuentes}
\newacronym{PCA}{PCA}{análisis de componentes principales}
\newacronym{ICA}{ICA}{análisis de componentes independientes}
\newacronym{NMF}{NMF}{factorización matricial no negativa}
\newacronym{NMSE}{NMSE}{error cuadrático medio normalizado}
\newacronym{MAE}{MAE}{error absoluto medio}
\newacronym{RMSE}{RMSE}{raíz del error cuadrático medio}
\newacronym{OOD}{OOD}{fuera de distribución}
\newacronym{SMPL}{SMPL}{Skinned Multi-Person Linear Model}

% Health / Human sensing
\newacronym{RR}{RR}{frecuencia respiratoria}
\newacronym{HR}{HR}{frecuencia cardiaca}
\newacronym{ECG}{ECG}{electrocardiograma}
\newacronym{PPG}{PPG}{fotopletismograma}
\newacronym{GT}{GT}{valor de referencia}

% ======================================================
%                COMANDOS ÚTILES
% ======================================================
\newcommand{\cpx}{\mathbb{C}}
\newcommand{\real}{\mathbb{R}}
\newcommand{\E}{\mathbb{E}}
\newcommand{\e}{\mathrm{e}}
\newcommand{\jj}{\mathrm{j}}
\newcommand{\calP}{\mathcal{P}}
\newcommand{\calO}{\mathcal{O}}
\newcommand{\calI}{\mathcal{I}}
\newcommand{\calL}{\mathcal{L}}
\newcommand{\calD}{\mathcal{D}}
\newcommand{\calR}{\mathcal{R}}
\newcommand{\vect}[1]{\bm{#1}}
\newcommand{\mat}[1]{\bm{#1}}
\DeclareMathOperator{\wrap}{wrap}
\DeclareMathOperator{\diag}{diag}
\DeclareMathOperator{\rank}{rank}
\DeclareMathOperator{\tr}{tr}

% ======================================================
%                    DOCUMENTO
% ======================================================

% Evita estirar verticalmente las páginas con tablas extensas.
\raggedbottom

\begin{document}
\hypersetup{pageanchor=false}

% ------------------------------------------------------
% Portada provisional
% ------------------------------------------------------
\begin{titlepage}
    \centering
    {\Large Tecnológico de Monterrey\par}
    \vspace{1.2cm}
    {\large Programa de Doctorado\par}
    \vfill
    {\Huge\bfseries Gemelo digital inalámbrico con estructura física y representación latente para el sensado humano mediante infraestructura de comunicaciones existente\par}
    \vspace{0.7cm}
    {\Large\itshape Modelado diferenciable, análisis de observabilidad y transferencia multimodal entre simulación y medición con Wi-Fi y LoRa\par}
    \vfill
    {\large Borrador de trabajo para tesis doctoral\par}
    \vspace{0.4cm}
    {\large Autor: [Nombre del doctorando]\\[0.2cm]
    Dirección de tesis: [Por completar]\par}
    \vfill
    {\large Guadalajara, México --- 2026\par}
\end{titlepage}


\hypersetup{pageanchor=true}
\pagenumbering{roman}
\chapter*{Resumen}
\addcontentsline{toc}{chapter}{Resumen}
Esta tesis propone un marco metodológico para el sensado humano mediante infraestructura inalámbrica existente. Una medición de canal no proporciona una observación directa de la presencia, la postura o la fisiología, sino el resultado de una cadena constituida por la geometría, la propagación electromagnética, la forma de onda, la instrumentación, el protocolo, el ruido y la interferencia. El trabajo plantea representar dicha cadena mediante un gemelo digital inalámbrico diferenciable y físicamente estructurado, calibrarlo con mediciones reales e invertirlo con una cuantificación explícita de la incertidumbre y la observabilidad. Wi-Fi y LoRa se estudian como observadores complementarios de un mismo estado físico. La validación progresa desde casos analíticos y un gemelo estático hasta el movimiento controlado, la respiración, la presencia de varios usuarios, la transferencia de la simulación a condiciones reales y el diseño de redes distribuidas. A lo largo del manuscrito se distinguen las contribuciones propuestas de los resultados que aún deben demostrarse experimentalmente.

Los experimentos documentados abarcan S1/B0--B2 y S4.1--S4.6d.2. Se utiliza un banco de 10\,000 posiciones de DICHASUS dc01, dividido inicialmente en 5000 posiciones de entrenamiento, 100 de validación y 4900 de prueba. Los materiales neuronales B2 reducen el MAE de potencia de 6.29 a 2.22 dB y el de dispersión RMS de retardos de 43.21 a 16.50 ns respecto de B0, sin mejorar la fase bruta. La descomposición posterior distingue una pendiente espectral efectiva, un desfase constante por enlace y un residual con estructura espectral. Sus asociaciones con descriptores energéticos se examinan mediante análisis entre posiciones, dentro de posición y remuestreo agrupado. La predicción del retardo se evalúa primero bajo interpolación densa y después mediante cinco bloques espaciales con exclusión de 0.5 m. En esta última evaluación, el predictor que combina posición y descriptores del trazado alcanza un MAE medio de retardo de 18.69 ns, frente a 27.56 ns al usar sólo posición. Aplicado al canal, y manteniendo ajuste del CPO con la medición evaluada, obtiene una mediana de error angular de 64.08 grados y coherencia de 0.448. S4.6 caracteriza 640\,000 fases residuales y recupera el tiempo de 10\,000 capturas. La resultante global es 0.001088, mientras que la concentración mediana dentro del arreglo es 0.4216. Las predicciones por vecinos espaciales y temporales de la dirección común producen errores medianos próximos a 90 grados, sin ventaja clara frente al control aleatorio. Estos resultados apoyan una distinción operacional entre retardo parcialmente predecible y fase común que requiere un tratamiento de referencia explícito; no identifican su causa física ni demuestran impredecibilidad universal. La evaluación de invariantes espaciales y espectrales S4.7a, la incertidumbre por trayectoria S4.7b, las representaciones alternativas y la fidelidad diferencial constituyen los siguientes bloques metodológicos.

La propuesta se articula ahora mediante dos ramas complementarias: una robusta frente
a transformaciones instrumentales identificadas y otra coherente que preserve respuesta
al movimiento. Se propone contrastar geometría manual, RGB-D, LiDAR, regularización RF
y calibración mediante G0--G4, separando precisión global, registro de antenas y referencia
local de desplazamiento. La campaña E0--E5 progresará desde estabilidad instrumental,
reflector mecánico y fantoma respiratorio hasta personas y transferencia. Referencia,
repetibilidad y sensibilidad se validarán antes de interpretar el problema inverso humano.
Estos bloques actualizan la metodología y la propuesta de trabajo; sus resultados aún
no forman parte de la evidencia experimental documentada.

\chapter*{Abstract}
\addcontentsline{toc}{chapter}{Abstract}
This dissertation proposes a methodological framework for human sensing through existing wireless communication infrastructure. A channel measurement does not directly observe presence, pose, or physiology; it is the output of a chain involving geometry, electromagnetic propagation, waveform, hardware, protocol, noise, and interference. The proposed approach models this chain with a physically structured differentiable wireless digital twin, calibrates it against real measurements, and inverts it while explicitly quantifying uncertainty and observability. Wi-Fi and LoRa are treated as different observers of the same physical state rather than as unrelated tasks. Validation progresses from analytical cases and a static twin to controlled motion, respiration, multiple users, simulation-to-reality transfer, and distributed-network design. Throughout the manuscript, proposed contributions are separated from empirical claims that remain to be validated.

The documented experiments span S1/B0--B2 and S4.1--S4.6d.2. A DICHASUS dc01 benchmark contains 10,000 positions, initially split into 5,000 training, 100 validation, and 4,900 test positions. Relative to nominal ray tracing B0, neural materials B2 reduce power MAE from 6.29 to 2.22 dB and RMS delay-spread MAE from 43.21 to 16.50 ns, without improving raw phase agreement. Subsequent analyses separate an effective spectral delay, a per-link constant phase offset, and a spectrally structured residual. Associations with path-energy descriptors are examined across and within positions using grouped resampling. Delay prediction is evaluated under dense interpolation and five spatial blocks with a 0.5 m exclusion margin. In the blocked evaluation, position and ray-tracing descriptors yield a mean delay MAE of 18.69 ns, compared with 27.56 ns using position alone. Applied to the channel while still fitting the constant phase offset from the evaluated measurement, the combined predictor yields a median angular error of 64.08 degrees and coherence of 0.448. S4.6 characterizes 640,000 residual phases and recovers timestamps for 10,000 captures. The global resultant length is 0.001088, whereas the median within-array concentration is 0.4216. Spatial and temporal nearest-neighbor predictions of the common direction yield median errors close to 90 degrees, without a clear advantage over random sampling. These findings support an operational distinction between partially predictable delay and common phase requiring explicit reference handling; they neither identify a unique physical cause nor establish universal unpredictability. Spatial and spectral invariant evaluation in S4.7a, pathwise uncertainty in S4.7b, alternative representations, and differential fidelity define the next methodological stages.

The updated proposal uses two complementary branches: one robust to identified acquisition
transformations and one coherent branch preserving motion sensitivity. G0--G4 will compare
manual geometry, RGB-D, LiDAR, RF regularization, and calibration, separating global
accuracy, antenna registration, and local displacement metrology. Scenarios E0--E5 will
progress from instrumental stability, a mechanical reflector, and a respiratory phantom
to human measurements and transfer. Reference quality, repeatability, and sensitivity
will be assessed before interpreting the human inverse problem. These stages update the
methodology and work proposal; their outcomes are not yet part of the documented evidence.

\clearpage
\tableofcontents
\clearpage
\glsaddall[types=\acronymtype]
\printglossary[type=\acronymtype,title=Acrónimos]
\clearpage
\chapter*{Convenciones de notación}
\addcontentsline{toc}{chapter}{Convenciones de notación}
\begingroup
\small
\setstretch{1.15}
\begin{longtable}{P{3.6cm}P{10.0cm}}
\toprule
\textbf{Símbolo} & \textbf{Uso reservado en el manuscrito}\\
\midrule
\endfirsthead
\toprule
\textbf{Símbolo} & \textbf{Uso reservado en el manuscrito}\\
\midrule
\endhead
$m$ & Modalidad o radio; principalmente Wi-Fi o LoRa.\\
$t,\tau,f$ & Tiempo físico, retardo de propagación y frecuencia, respectivamente.\\
$n,k,\ell$ & Índice de captura o paquete, subportadora y trayectoria.\\
$N_p,N_{\rm link}$ & Número de trayectorias y número de enlaces espaciales.\\
$G_E,G_H,G_F$ & Geometría/materiales del entorno, geometría humana y deformación fisiológica.\\
$\Gamma(t)$ & Soporte geométrico común de interacciones posibles en la escena.\\
$\calP_m(t)$ & Conjunto de trayectorias electromagnéticas efectivas para la modalidad $m$.\\
$h_m(t,\tau)$ & Respuesta al impulso ideal de la modalidad $m$.\\
$H_m(t,f)$ & Respuesta en frecuencia ideal; $\widehat H_m$ denota su estimación.\\
$Y_m$ & Dato finalmente disponible después de la forma de onda, la instrumentación, el protocolo y el muestreo.\\
$\vect s_E,\vect s_H,\vect s_F$ & Estados físicos del entorno, macroestado humano y fisiología.\\
$\vect s_R,\vect s_N$ & Configuración de radio y perturbaciones instrumentales.\\
$Z_E,Z_\phi,Z_R$;\newline $Z_H,Z_F,Z_N$ & Bloques de representación latente propuestos; no son sinónimos del estado físico ni presuponen un codificador ya validado.\\
$\phi_0^\star,\Delta\tau_{\rm eff}$ & Desfase residual efectivo y retardo relativos a medición y modelo; no se identifican automáticamente con CPO/STO físicos.\\
$\beta_{u,a},\psi_{i,u}$ & Dirección marginal por antena y dirección común al arreglo en S4.6.\\
$G_u(k),K_u(k,l)$ & Productos espaciales y entre frecuencias utilizados en el protocolo invariante S4.7a.\\
$i,u,a,b$ & En los avances S1/S4: posición, arreglo, antena y modelo B0/B1/B2.\\
$E_{\phi,\rm deg},\rho$ & Error circular absoluto en grados y coherencia compleja normalizada.\\
$x_{\rm chest}(t)$ & Desplazamiento o forma de onda torácica; RR es una cantidad derivada.\\
$y_{\rm L}(t)$ & Señal IQ recibida de LoRa.\\
$\vect q$ & Vector de parámetros físicos de interés para inversión.\\
$\mat J_q$ & Jacobiano del observable respecto de $\vect q$.\\
$\mat F_q,\mat F_q^{\rm eq}$ & Matriz de información de Fisher y su forma equivalente tras eliminar perturbaciones.\\
$\mat\Sigma_w,\mat\Sigma_q$ & Covarianza del ruido y covarianza estimada/posterior del estado.\\
$\mathcal R_{\rm EM}^{(m)},\mathcal O_m$ & Simulador electromagnético y operador de observación.\\
$\calL,R$ & Función de pérdida y regularizador; $R_{a,k}$ denota específicamente el producto residual de canales en S4.\\
$\mathsf V_0$--$\mathsf V_7$ & Alternativas de observable, distancia o fusión para S4.7a; no son los regímenes de información $\mathcal R_i$.\\
$z_{\rm rob},z_{\rm coh}$ & Ramas robusta y coherente propuestas.\\
$T_{W\leftarrow s}$ & Transformación rígida desde el marco del sensor $s$ al marco global $W$.\\
$D_{\rm rep},D_{\rm fis},Q$ & Distancias de repetibilidad, respuesta física y cociente descriptivo; no sustituyen una prueba de identificabilidad.\\
G0--G4; Esc. E0--E5 & Condiciones de reconstrucción/calibración y escenarios físicos propuestos, respectivamente.\\

\bottomrule
\end{longtable}
\endgroup
\clearpage

\pagenumbering{arabic}
\chapter{Introducción}
\label{ch:introduccion}

\section{Del enlace de comunicación al instrumento de observación}

La ingeniería de comunicaciones inalámbricas estudia cómo transportar información de manera confiable a través de un medio que atenúa, dispersa, retarda y distorsiona la señal. Su teoría clásica conecta mecanismos de propagación, modelos de canal, formas de onda y arquitectura del receptor \citep{Molisch2022,TseViswanath2005,Goldsmith2005}. Esta tesis adopta esa misma cadena, pero cambia la variable de interés: el canal deja de ser solamente una perturbación que debe compensarse para recuperar bits y pasa a ser también una medición indirecta del entorno.

La reinterpretación es física antes que algorítmica. Las reflexiones, difracciones, dispersiones, obstrucciones y movimientos modifican la amplitud, el retardo, la fase, el desplazamiento Doppler y la dirección de llegada de las contribuciones de multitrayectoria. Por tanto, una observación del canal contiene información acerca de aquello que lo perturbó. Sin embargo, esa información es parcial: depende de la geometría transmisor--objetivo--receptor, la longitud de onda, el ancho de banda, la diversidad espacial, la tasa de muestreo, la instrumentación y el ruido. Un sistema de sensado no puede recuperar información que su operador de observación no haya preservado.

Esta perspectiva sustenta el campo del sensado inalámbrico. En Wi-Fi, la disponibilidad de \ac{CSI} ha permitido estudiar presencia, localización, actividades, postura y señales fisiológicas; aun así, la generalización entre habitaciones, sujetos y dispositivos continúa siendo un problema central \citep{Zhu2024Survey,Armenta2024Survey}. En LoRa, los estudios de sensado a través de paredes, con tráfico ambiente y en redes de gran escala sugieren una relación distinta entre alcance, densidad temporal e información de movimiento \citep{Zhang2020LoRa,SenLoRa2025,Xue2025LargeScaleLoRa}. La comparación de porcentajes obtenidos en conjuntos de datos independientes no permite determinar si esas diferencias provienen de la propagación, de la forma de onda, de la instrumentación o del protocolo.

\section{Vacío científico y tesis de trabajo}

El problema no se reduce a diseñar un clasificador sobre CSI. En una medición real aparecen mezcladas al menos cinco familias de variables: entorno, dinámica humana macroscópica, fisiología, configuración de radio y perturbaciones instrumentales. Un modelo puramente discriminativo puede minimizar su pérdida utilizando la firma de una habitación o de un dispositivo en vez de la variable humana pretendida. El buen desempeño dentro del dominio no demuestra que se haya identificado el mecanismo causal ni que el sistema pueda transferirse.

La pregunta rectora es:
\begin{quote}
\emph{¿Puede construirse una representación físico--latente del sistema de propagación inalámbrica que separe el entorno, la configuración de radio, la instrumentación, la dinámica humana y la fisiología, de modo que un mismo marco explique el canal directo, permita invertirlo para recuperar el estado humano y cuantifique cuándo dicho estado es realmente observable?}
\end{quote}

La hipótesis de trabajo sostiene que una representación jerárquica y físicamente estructurada, calibrada con mediciones y complementada con aprendizaje multimodal, conservará mejor las variaciones relevantes para el sensado y requerirá menos datos de adaptación que una representación monolítica de la señal. Esta proposición se someterá a prueba y no se asume como resultado.

\section{Modelo generativo y contrato de notación}

Sea el estado total
\begin{equation}
\vect s(t)=
\left[
\vect s_E,\vect s_H(t),\vect s_F(t),
\vect s_R,\vect s_N(t)
\right],
\label{eq:total_state}
\end{equation}
donde los subíndices $E$, $H$, $F$, $R$ y $N$ denotan, respectivamente, entorno, macroestado humano, fisiología, radio y factores instrumentales o no controlados. Se reserva $G_E$ para la geometría y los materiales del entorno, $G_H(t)$ para la geometría humana y $G_F(t)$ para la deformación fisiológica; $Z_\bullet$ denotará siempre una representación latente aprendida, no el estado físico mismo.

La escena dinámica determina un soporte geométrico de interacciones $\Gamma(t)$. Para cada modalidad $m$ y configuración de radio se obtiene un conjunto electromagnético de trayectorias
\begin{equation}
\calP_m(t)=
\mathcal R_{\rm EM}^{(m)}
\left(\Gamma(t);\vect s_R\right),
\label{eq:radio_dependent_paths}
\end{equation}
del cual se sintetiza el canal $H_m$ y finalmente la medición
\begin{equation}
Y_m(t)=
\mathcal O_m
\left(H_m[\calP_m(t)];\vect s_R,\vect s_N(t)\right)
+\vect{\eta}_m(t).
\label{eq:generative_observation}
\end{equation}
Esta distinción es deliberada: Wi-Fi y LoRa comparten mundo físico y movimiento, pero sus ganancias, fases, antenas y aun sus trayectorias efectivamente observables dependen de frecuencia y ancho de banda. El objeto común es $\Gamma(t)$; $\calP_m(t)$ es radio--dependiente.

La cadena completa de la tesis es
\begin{equation}
\begin{aligned}
\underbrace{G_E,G_H(t),G_F(t)}_{\text{mundo}}
&\rightarrow \Gamma(t)
\xrightarrow{\;\mathcal R_{\rm EM}^{(m)}\;}
\calP_m(t)\rightarrow H_m(t)\\
&\xrightarrow{\;\mathcal O_m\;}Y_m(t)
\xrightarrow{\;E_m\;}
\{Z_E,Z_\phi,Z_R,Z_H,Z_F,Z_N\}\\
&\rightarrow
\{\widehat{\vect q},\mat\Sigma_q,\mat F_q\}.
\end{aligned}
\label{eq:thesis_chain}
\end{equation}
Cada flecha establece una relación que se utilizará en el resto del manuscrito: el gemelo resuelve el modelo directo hasta $Y_m$; el codificador factoriza la observación; el problema inverso estima el estado; y la información de Fisher $\mat F_q$ expresa qué direcciones del estado pueden determinarse localmente.

\section{Posicionamiento metodológico y alcance}

El trabajo no corresponde, en sentido estricto, a una \ac{PINN}. Una red de este tipo aproxima la solución de una ecuación diferencial y penaliza su residuo \citep{Raissi2019PINN}. En esta investigación, la física interviene principalmente como estructura del modelo generativo: la geometría, la propagación, la síntesis del canal y la observación de radio forman un grafo computacional diferenciable. El aprendizaje se reserva para calibrar parámetros desconocidos, aproximar bloques costosos, construir representaciones y acelerar la solución del problema inverso. En adelante, esta formulación se denomina \emph{gemelo digital inalámbrico diferenciable y fundamentado en la física}.

El alcance primario comprende Wi-Fi y LoRa, presencia, movimiento macroscópico y respiración. La postura detallada, el ritmo cardiaco, la fisiología de varios usuarios y el despliegue exterior son extensiones progresivas sujetas a los criterios de observabilidad definidos posteriormente. La tesis no presupone que toda variable sea recuperable ni que una sola representación sea óptima para todas las tareas. Precisamente busca distinguir un fallo del estimador de una carencia física de información.

\section{Avance experimental y delimitación de la evidencia}

El aporte principal sobre la fase es una mejora efectiva mediante un retardo
parcialmente predecible. Se distinguen concordancia bruta, capacidad de alineación
con la medición y transferencia de la corrección a datos reservados. Con retardo
predicho y manteniendo ajuste del desfase constante por enlace, el error angular
mediano disminuye de 81.42 a 45.59 grados en interpolación densa y de 81.49 a
64.08 grados por bloques espaciales. La corrección completa de fase absoluta
permanece abierta, pero la utilidad predictiva del retardo queda documentada.
La \cref{sec:synthesis} y el \cref{ch:contributions} presentan respectivamente
el balance experimental y los aportes consolidados.


La primera etapa comprende el banco DICHASUS dc01, la calibración B0--B2 y la
serie S4.1--S4.6d.2. Tras comprobar que la mejora macroscópica no se traduce en
alineación angular bruta, se estudian retardo efectivo, CPO y residual espectral.
S4.3 relaciona el residual con descriptores energéticos mediante análisis
agrupado; S4.4 compara predictores y audita proximidad y separación espacial;
S4.5 cuantifica el efecto de predecir el retardo sobre el canal corregido.
S4.6 analiza 640\,000 fases residuales, recupera sus marcas temporales y compara
reglas de vecindad espacial y temporal para la dirección común al arreglo.
Su falta de ventaja clara motiva evaluar productos relativos e invariantes
en S4.7a, sin atribuir automáticamente esa fase al hardware.
La validación por bloques aporta evidencia regional dentro de dc01, condicionada
a un CPO ajustado con la medición evaluada. La fase completamente predicha,
la transferencia entre escenas y la fidelidad ante movimiento humano requieren
experimentos posteriores. El \cref{ch:avances} presenta la evidencia y sus límites.

\section{Orientación de la propuesta después del diagnóstico de fase}

El programa actualizado conserva una representación robusta de estructura relativa y
una rama coherente sensible al movimiento. La distinción requiere medir perturbaciones
instrumentales e intervenciones físicas: no toda fase común es irrelevante para la tarea.
El gemelo incorporará geometría métrica con incertidumbre, registro de antenas y una
referencia local de desplazamiento, en lugar de atribuir a la reconstrucción visual
la capacidad de determinar por sí sola la fase absoluta.

La campaña propuesta progresa desde E0, estabilidad instrumental, hasta E1, reflector
medido, y E2, fantoma respiratorio. E3 introduce una persona, E4 la separabilidad de
dos fuentes y E5 transferencia entre condiciones independientes. La comparación G0--G4
cuantificará el aporte de geometría y calibración. Esta secuencia vincula DICHASUS con
la investigación humana sin equiparar el desplazamiento del transmisor de dc01 con
movimiento torácico. Las hipótesis, el diseño y la propuesta administrativa se ajustan
a esos contrastes en los capítulos correspondientes.

\section{Ruta de lectura}

El manuscrito se desarrolla en cuatro niveles sucesivos. Primero se establece el lenguaje de comunicaciones necesario para describir propagación, canales variantes en el tiempo, multitrayectoria, coherencia, sistemas MIMO y estimación. Después se introduce el gemelo digital y se estudian las representaciones del entorno. El tercer nivel incorpora los operadores de observación de Wi-Fi y LoRa, la dinámica humana, la observabilidad y la transferencia entre simulación y medición. Finalmente, el marco conceptual se traduce en preguntas, hipótesis, líneas metodológicas, experimentos y criterios verificables de avance. De este modo, las cantidades definidas al inicio ---retardo, Doppler, fase, ancho de coherencia, jacobiano y covarianza de ruido--- reaparecen como magnitudes observables en el programa experimental. Esta progresión también determina el orden de los capítulos y evita presentar la simulación, el aprendizaje y la instrumentación como componentes aislados.

Después del diseño experimental se incorpora un capítulo de avances con la secuencia S1/B0--B2 y S4.1--S4.6d.2. Esta ubicación permite leer primero los criterios generales y después contrastarlos con la evidencia disponible, antes de revisar las contribuciones, el plan de trabajo y las conclusiones provisionales.

\part{Fundamentos y modelo físico}
\chapter{Fundamentos del canal inalámbrico: de propagación a observación}
\label{ch:marco_canal}

\section{Jerarquía de modelos y escalas}

Un modelo útil no reproduce todos los grados de libertad de la realidad; conserva aquellos que determinan la pregunta estudiada. En electromagnetismo, el nivel más completo parte de las ecuaciones de Maxwell y de condiciones de frontera sobre materiales. En sistemas de comunicación se adopta habitualmente una jerarquía de reducciones: campo electromagnético $\rightarrow$ trayectorias o respuesta de Green $\rightarrow$ canal equivalente en banda base $\rightarrow$ muestras del receptor. Molisch organiza precisamente la teoría inalámbrica alrededor de la interacción entre propagación, canal y diseño de sistema, mientras que Tse y Viswanath enfatizan cómo la estructura del canal condiciona diversidad, multiplexación y estimación \citep{Molisch2022,TseViswanath2005}.

En esta tesis se utilizará el modelo más sencillo que conserve la variable de sensado pertinente. La cobertura promedio puede estudiarse mediante la pérdida de trayecto; la respuesta compleja de banda ancha exige retardos y ganancias de multitrayectoria; la respiración exige, además, coherencia de fase y evolución temporal. Esta regla evita dos extremos: atribuir realidad física a un modelo estadístico que sólo describe promedios, o resolver un problema electromagnético completo cuando una aproximación local ofrece la precisión necesaria.

\section{De Maxwell a los mecanismos de propagación}

En un medio lineal, isotrópico y local, los campos satisfacen
\begin{align}
\nabla\times\vect E&=-\frac{\partial\vect B}{\partial t}, &
\nabla\times\vect H&=\vect J+\frac{\partial\vect D}{\partial t},\\
\nabla\cdot\vect D&=\rho, &
\nabla\cdot\vect B&=0,
\end{align}
con relaciones constitutivas $\vect D=\epsilon\vect E$, $\vect B=\mu\vect H$ y $\vect J=\sigma\vect E$. La permitividad compleja resume almacenamiento y pérdida; su valor depende de material y frecuencia. Las discontinuidades de $\epsilon$, $\mu$ y $\sigma$ producen reflexión y transmisión; aristas y objetos finitos introducen difracción y dispersión. En interiores, la señal recibida resulta de la superposición de estas interacciones, no de un único rayo geométrico \citep{Molisch2022,Rappaport2002}.

La longitud de onda
\begin{equation}
\lambda=\frac{c}{f_c}
\label{eq:wavelength}
\end{equation}
fija la escala de sensibilidad de fase. Un desplazamiento que parece pequeño en metros puede representar una fracción apreciable de $\lambda$ y modificar notablemente el canal complejo. Esta relación será reutilizada en la derivación biestática de la \cref{eq:bistatic_sensitivity} y en el experimento de fidelidad diferencial.

\section{Presupuesto de enlace y variación a gran escala}

En espacio libre y campo lejano, la ecuación de Friis expresa
\begin{equation}
P_r=P_tG_tG_r
\left(\frac{\lambda}{4\pi d}\right)^2,
\label{eq:friis}
\end{equation}
donde $d$ es distancia y $G_t,G_r$ son ganancias en las direcciones relevantes \citep{Friis1946}. En entornos reales se emplea con frecuencia un modelo log-distancia
\begin{equation}
PL(d)=PL(d_0)+10n\log_{10}\!\left(\frac{d}{d_0}\right)+X_\sigma,
\label{eq:logdistance}
\end{equation}
con exponente $n$ y sombreado lognormal $X_\sigma$. Estas ecuaciones describen tendencia y cobertura de comunicación; no describen por sí solas la respuesta a un desplazamiento torácico. Esta separación entre potencia media y sensibilidad local anticipa una tesis recurrente: buen RSSI o buena \ac{SNR} son condiciones de medición, pero no garantizan observabilidad de una variable humana.

\section{Representación equivalente en banda base}

Sea una señal pasabanda real
\begin{equation}
x_{\rm RF}(t)=\sqrt{2}\,\Re\left\{x(t)e^{j2\pi f_ct}\right\},
\label{eq:complex_envelope}
\end{equation}
donde $x(t)$ es su envolvente compleja. La representación en banda base conserva amplitud y fase relativas y permite trabajar con muestras \ac{IQ}. Para un canal lineal variante en el tiempo, la relación entrada--salida es
\begin{equation}
y(t)=\int h(t,\tau)x(t-\tau)\,d\tau+n(t),
\label{eq:ltv_channel}
\end{equation}
donde $h(t,\tau)$ es la respuesta al impulso dependiente del tiempo y $n(t)$ agrega ruido e interferencia. La caracterización de canales \ac{LTV} mediante retardo y Doppler proporciona el lenguaje formal para separar variación rápida y dispersión temporal \citep{Bello1963,Molisch2022}.

Si el canal cambia poco durante una ventana de observación, puede aproximarse como \ac{LTI}: $h(t,\tau)\approx h(\tau)$. Esta es una aproximación local, no una propiedad permanente del entorno. En el sensado, el movimiento humano es precisamente aquello que viola la invariancia temporal y genera información; por tanto, la duración de la ventana debe declararse siempre que se adopte el modelo LTI.

\section{Retardo, Doppler y coherencia}

El perfil de potencia--retardo se define como $P_h(\tau)=\E[|h(t,\tau)|^2]$. A partir de él se obtienen el retardo medio $\bar\tau$ y el \ac{RMSDS}
\begin{align}
\bar\tau&=\frac{\int \tau P_h(\tau)d\tau}{\int P_h(\tau)d\tau},\\
\tau_{\rm rms}&=
\left[
\frac{\int(\tau-\bar\tau)^2P_h(\tau)d\tau}
{\int P_h(\tau)d\tau}
\right]^{1/2}.
\label{eq:rms_delay}
\end{align}
El ancho de coherencia $B_c$ es inversamente proporcional a la dispersión de retardos, aunque el factor exacto depende del umbral de correlación adoptado. De forma dual, el espectro Doppler describe la distribución de variación temporal y el tiempo de coherencia $T_c$ es inversamente proporcional a su extensión. No se usarán $B_c$ o $T_c$ sin indicar la definición operativa, pues no son constantes universales \citep{Molisch2022,Goldsmith2005}.

Estas cantidades conectan directamente fundamentos y experimento. Si el ancho de señal $B$ supera de forma apreciable a $B_c$, el canal es selectivo en frecuencia y las subportadoras ofrecen diversidad; el tiempo de coherencia caracteriza correlación, pero no proporciona por sí solo un teorema de muestreo. La reconstrucción dinámica depende del espectro de la variable objetivo, los tiempos observados, los huecos de tráfico y el ruido. La selectividad frecuencial y la disponibilidad temporal se evalúan por separado.

\section{Modelos deterministas y estadísticos}

Los modelos estadísticos describen distribuciones de pérdida, desvanecimiento, retardo y Doppler sobre una clase de entornos. Son indispensables para el diseño robusto y para sintetizar variabilidad, pero no identifican necesariamente qué pared o persona generó una componente. Los modelos deterministas condicionados por la geometría, como el trazado de rayos, conservan esa correspondencia a costa de requerir una escena y propiedades materiales. Los modelos híbridos combinan una estructura geométrica con componentes difusas o parámetros aprendidos \citep{Molisch2022,SalehValenzuela1987}.

La tesis asigna funciones distintas a ambos paradigmas. El gemelo utiliza un modelo condicionado por la escena para realizar intervenciones y calcular derivadas; las distribuciones estadísticas describen los errores residuales, la instrumentación, la dispersión no modelada y las condiciones fuera de distribución. Un ajuste estadístico no se interpretará como una reconstrucción causal de la escena.

\section{Canal multitrayecto condicionado al entorno}

Considérese un transmisor en $\vect p_T$, un receptor en $\vect p_R$ y un entorno $G_E$. Para una modalidad $m$, el canal complejo puede aproximarse mediante $N_p^{(m)}$ contribuciones:
\begin{equation}
H_m(f,t)=\sum_{\ell=1}^{N_p^{(m)}(t)}
\alpha_{m,\ell}(f,t)
\exp\{-j2\pi f\tau_{m,\ell}(t)\},
\label{eq:multipath}
\end{equation}
con $\alpha_{m,\ell}$ el coeficiente complejo y $\tau_{m,\ell}=d_{m,\ell}/c$ el retardo asociado a longitud $d_{m,\ell}$. Cada elemento se representa como
\begin{equation}
\calP_{m,\ell}=
\left\{\alpha_{m,\ell},\tau_{m,\ell},
\theta^{\rm AoA}_{m,\ell},
\theta^{\rm AoD}_{m,\ell},
\nu_{m,\ell},
\chi_{m,\ell}\right\},
\label{eq:path_element}
\end{equation}
donde $\nu_{m,\ell}$ es el desplazamiento Doppler y $\chi_{m,\ell}$ codifica la secuencia de interacciones. El subíndice $m$ impide suponer que las ganancias y las trayectorias efectivas son idénticas entre bandas. La correspondencia entre tecnologías de radio se establece mediante el soporte geométrico $\Gamma$, no mediante la igualdad directa de $\calP_{\rm WiFi}$ y $\calP_{\rm LoRa}$.

La respuesta al impulso es
\begin{equation}
h_m(t,\tau)=\sum_{\ell=1}^{N_p^{(m)}(t)}
\alpha_{m,\ell}(t)\,\delta(\tau-\tau_{m,\ell}(t)),
\end{equation}
y supone que la ganancia de cada interacción es aproximadamente plana dentro del ancho de banda modelado. Si $\alpha_{m,\ell}(f,t)$ varía de forma apreciable con $f$, la contribución deja de ser un delta ponderado y debe representarse mediante una respuesta dispersiva por trayectoria.
y la relación entre \ac{CIR} y \ac{CFR} está dada por transformada de Fourier:
\begin{equation}
H_m(t,f)=\int h_m(t,\tau)e^{-j2\pi f\tau}\,d\tau.
\end{equation}
En un sistema de banda limitada, dos trayectorias con retardos próximos pueden no ser separables. La resolución temporal aproximada es
\begin{equation}
\Delta\tau\approx \frac{1}{B},
\qquad
\Delta d_{\rm path}\approx \frac{c}{B},
\label{eq:delay_resolution}
\end{equation}
donde $\Delta d_{\rm path}$ es diferencia de longitud de trayectoria; para rango monostático el factor es $c/(2B)$. Esta distinción explica por qué Wi-Fi convencional conserva información de micro-movimiento en la fase aun cuando no sea capaz de resolver físicamente cada reflector en retardo.

\section{Dimensión espacial y canal MIMO}

Con $N_t$ antenas transmisoras y $N_r$ receptoras, cada frecuencia posee una matriz $\mat H(f,t)\in\cpx^{N_r\times N_t}$. Sus entradas no son repeticiones independientes: comparten elementos dispersores y presentan una correlación espacial determinada por la geometría, la separación entre antenas, la polarización y los diagramas de radiación. En comunicaciones, esta estructura determina la diversidad y la multiplexación; en sensado, proporciona vistas biestáticas con jacobianos distintos respecto del estado humano \citep{TseViswanath2005,Molisch2022}.

La resolución angular no depende sólo del número de antenas, sino de la apertura efectiva y de la relación entre espaciado y longitud de onda. Por ello, el número de enlaces no se utilizará como sinónimo de número de observaciones independientes. Más adelante, la FIM conjunta medirá cuánta información adicional aporta realmente cada enlace.

\section{OFDM y CSI}

Para una subportadora $k$ y un instante discreto $n$,
\begin{equation}
Y_n[k]=H_n[k]X_n[k]+W_n[k],
\end{equation}
y, cuando $X_n[k]$ es conocido,
\begin{equation}
\widehat H_n[k]=\frac{Y_n[k]}{X_n[k]}.
\end{equation}
En \ac{MIMO}, $H_n[k]\in\cpx^{N_r\times N_t}$. Por tanto, una captura puede representarse como un tensor complejo
\begin{equation}
\mat H\in\cpx^{N_{\rm snap}\times N_{\rm link}\times K},
\end{equation}
donde $N_{\rm snap}$ es el número de capturas, $N_{\rm link}=N_tN_r$ representa los enlaces espaciales y $K$ el número de subportadoras. La ortogonalidad de \ac{OFDM} convierte una convolución con memoria en ganancias aproximadamente escalares por subportadora cuando el \ac{CP} excede el soporte significativo del canal \citep{Molisch2022,ProakisSalehi2008}. La CSI es una estimación de esas ganancias, no el canal físico libre de efectos instrumentales.

Las revisiones del sensado mediante Wi-Fi distinguen enfoques basados en modelos físicos, métodos basados en patrones y métodos de aprendizaje profundo. Asimismo, subrayan que la dificultad práctica se concentra en la variación entre dominios, la cobertura y la estabilidad de las mediciones \citep{Zhu2024Survey,Armenta2024Survey}. Esta taxonomía sugiere que un marco híbrido entre modelo y datos resulta más adecuado para el objetivo de esta tesis que una solución puramente discriminativa.

\section{Movimiento humano como perturbación del camino}

Sea un punto del cuerpo $\vect p_H(t)$ que perturba una trayectoria biestática Tx--humano--Rx. Su longitud es
\begin{equation}
L_H(t)=\|\vect p_T-\vect p_H(t)\|_2+
\|\vect p_H(t)-\vect p_R\|_2.
\end{equation}
La fase correspondiente es
\begin{equation}
\phi_H(t)=-\frac{2\pi}{\lambda}L_H(t).
\end{equation}
Se definen $\widehat{\vect u}_{TH}=(\vect p_H-\vect p_T)/\|\vect p_H-\vect p_T\|$
y $\widehat{\vect u}_{RH}=(\vect p_H-\vect p_R)/\|\vect p_H-\vect p_R\|$,
ambos dirigidos desde cada radio hacia el punto humano. Para un pequeño desplazamiento $x$ en dirección $\vect n$,
\begin{equation}
\frac{\partial \phi_H}{\partial x}
=-\frac{2\pi}{\lambda}
\left(\widehat{\vect u}_{TH}+\widehat{\vect u}_{RH}\right)^T\vect n.
\label{eq:bistatic_sensitivity}
\end{equation}
Esta ecuación es central: el mismo desplazamiento torácico puede ser fuertemente visible o casi invisible dependiendo de la geometría. Por ello, la potencia recibida o la \ac{SNR} de comunicación no son equivalentes a observabilidad de una variable fisiológica.

\section{Fresnel, interferencia constructiva y puntos ciegos}

La contribución dinámica no se observa aislada, sino sumada a un canal estático:
\begin{equation}
H(t)=H_s+H_d(t).
\end{equation}
Esta descomposición se define respecto de un estado basal y una ventana temporal; no es una separación física única. La amplitud resultante depende de la relación geométrica y de fase entre ambas componentes. Dos posiciones con niveles de potencia similares pueden presentar sensibilidades radicalmente diferentes a desplazamientos pequeños. Este fenómeno aparece repetidamente en la literatura de respiración basada en CSI y justifica el uso de razones complejas, diversidad de enlaces y métricas basadas en estructura de fase, en lugar de seleccionar canales exclusivamente por energía \citep{Mosleh2022BreatheSmart,Zhu2024Survey}.

\section{Distorsiones instrumentales}

Una observación COTS puede modelarse como
\begin{equation}
\widetilde H_n[k]
=g_n[k]H_n[k]\exp\left\{j\left[
\varphi_{0,n}+2\pi\Delta f_{\rm CFO}t_n
-2\pi f_k\Delta\tau_n
\right]\right\}+w_n[k],
\label{eq:hardware_model}
\end{equation}
donde $g_n[k]\geq0$ incluye la ganancia y la respuesta de la etapa de entrada del receptor, $\varphi_{0,n}$ es la fase común, $\Delta f_{\rm CFO}$ produce rotación temporal y $\Delta\tau_n=\tau_{\rm STO}+t_n\epsilon_{\rm SFO}$ produce una pendiente de fase dependiente de la subportadora que puede derivar entre paquetes. El modelo es deliberadamente parsimonioso; las versiones específicas de cada dispositivo podrán incluir ruido de fase, no linealidad y control automático de ganancia. Una metodología rigurosa debe distinguir la calibración de la medición de la calibración electromagnética: la primera busca un observable coherente; la segunda ajusta materiales, geometría e interacciones.

\subsection{Taxonomía operacional de la fase}
\label{sec:phase_taxonomy}

La fase observada no constituye una sola variable física. Una factorización mínima,
válida como modelo local de adquisición, es
\begin{equation}
\widehat H_{a,k,t}=
g_{a,k,t}\exp\!\left[j\psi_{a,t}-j2\pi\nu_k\delta\tau_t\right]
H_{a,k,t}^{\rm phys}+n_{a,k,t},
\label{eq:observed_phase_taxonomy}
\end{equation}
donde $H^{\rm phys}$ es la suma coherente de las trayectorias del entorno y del
objeto dinámico; $g$ reúne ganancia y respuesta compleja de cadena; $\psi$ incluye
la referencia de fase de la captura o de la cadena; $\delta\tau$ es un error de
referencia temporal en la convención adoptada; y $n$ representa ruido y discrepancia.
La expresión no afirma independencia ni identificabilidad desde una sola observación.

La metodología distinguirá cuatro orígenes que pueden producir patrones algebraicamente
semejantes:
\begin{enumerate}
\item \emph{fase de propagación física}, determinada por longitud de trayectoria,
interacciones, movimiento y frecuencia;
\item \emph{fase por discrepancia geométrica}, causada por errores de escena, registro,
centro de fase o posición;
\item \emph{fase de adquisición}, asociada con relojes, osciladores, cadenas RF,
reinicios, CFO/SFO/STO y referencia común;
\item \emph{fase de procesamiento}, introducida por compensación, interpolación,
selección de portadora, alineación o convención de banda base.
\end{enumerate}
Una pendiente espectral no identifica por sí sola STO, pues una discrepancia de
longitud produce la misma forma. Una rotación común tampoco es automáticamente
instrumental: un movimiento puede proyectarse sobre esa dirección. La identificación
exige metadatos, una referencia independiente o intervenciones que modifiquen un
factor manteniendo los demás controlados. Esta taxonomía conecta la calibración de
errores locales de fase de Ruah y colaboradores \citep{Ruah2024Phase} con los contrastes de invariancia y
fidelidad diferencial, sin interpretar sus variables latentes como geometría recuperada.


\section{Convenciones de frecuencia, materiales e interferencia coherente}
\label{sec:phase_conventions}

La síntesis del canal debe declarar si los coeficientes complejos contienen la fase de
propagación a la portadora. Con frecuencia física $f=f_{\rm ref}+\nu$ y coeficiente
$\widetilde\alpha_\ell(f)$ que excluye esa fase, se tiene
\begin{equation}
H(f)=\sum_\ell\widetilde\alpha_\ell(f)e^{-j2\pi f\tau_\ell}
=\sum_\ell\alpha_\ell(f;f_{\rm ref})e^{-j2\pi\nu\tau_\ell},\qquad
\alpha_\ell=\widetilde\alpha_\ell e^{-j2\pi f_{\rm ref}\tau_\ell}.
\label{eq:reference_frequency}
\end{equation}
Ambas expresiones son equivalentes. Aplicar $e^{-j2\pi f\tau_\ell}$ a un coeficiente
que ya contiene la fase de portadora la contaría dos veces. En los experimentos S4 se
usa una referencia común y $\nu_k=f_k-f_{\rm ref}$ para separar intercepto y pendiente.
En el modelo instrumental de la \cref{eq:hardware_model}, la frecuencia multiplicada
por el error de temporización se interpreta igualmente como frecuencia de banda base;
el término constante que produzca otra referencia se absorbe en el intercepto.

Para la convención temporal $e^{j\omega t}$, un dieléctrico conductor puede representarse
por $\epsilon_c=\epsilon_0\epsilon_r-j\sigma/\omega$. En incidencia plana entre medios
homogéneos, los coeficientes de Fresnel dependen de polarización, ángulo e impedancia;
por ejemplo, para polarización transversal eléctrica,
\begin{equation}
\Gamma_{\rm TE}=\frac{\eta_2\cos\theta_i-\eta_1\cos\theta_t}
{\eta_2\cos\theta_i+\eta_1\cos\theta_t},\qquad
\eta_m=\sqrt{\frac{j\omega\mu_m}{\sigma_m+j\omega\epsilon_m}}.
\end{equation}
La rama se elige compatible con un medio pasivo y las condiciones de radiación.
Esta dependencia explica por qué un material modifica amplitud y fase de cada interacción.
No obstante, los parámetros recuperados con un canal finito pueden ser efectivos si
compensan geometría o mecanismos ausentes. La jerarquía de campo, mecanismos de
propagación, estadísticas y canal de banda base sigue el enfoque de Molisch
\citep{Molisch2022}.

Con caminos coherentes, la potencia contiene términos cruzados:
\begin{equation}
|H(f)|^2=\sum_\ell|\alpha_\ell|^2+
2\Re\sum_{\ell<r}\alpha_\ell\alpha_r^*
 e^{-j2\pi\nu(\tau_\ell-\tau_r)}.
\label{eq:coherent_power}
\end{equation}
Por tanto, sumar energías individuales no equivale a calcular la potencia del canal.
Las fracciones de energía de S4.3 son descriptores del soporte multitrayecto. Su utilidad
predictiva debe evaluarse frente al canal complejo, sin interpretar la normalización
como una separación experimental de cada camino físico.

\section{Retardo absoluto, retardo relativo y dispersión temporal}
\label{sec:delay_gauge}

Una transformación $H'(\nu)=e^{j\phi_0}e^{-j2\pi\nu\delta\tau}H(\nu)$ conserva
la magnitud y traslada la respuesta temporal ideal: $h'(\tau)=e^{j\phi_0}h(\tau-\delta\tau)$.
Si se observa el perfil completo, cambia el retardo medio pero no su dispersión RMS.
Así, potencia y RMS-DS no identifican el origen de retardos ni una fase constante.
En una IFFT finita deben controlarse ventana, máscara y plegado circular antes de
aplicar esta invariancia a un estadístico estimado.

Un retardo relativo ajustado entre medición y simulación puede contener diferencias
de referencia, registro geométrico o interferencia entre caminos. Aunque las mediciones
ya tengan calibración de STO/CPO, ese parámetro puede seguir siendo distinto de cero.
S4.2 lo denomina retardo efectivo y no una estimación identificada del reloj del receptor.

El caso analítico de espacio libre verifica la escala y fase; una reflexión controlada
verifica Fresnel; la multitrayectoria verifica interferencia y resolución en retardo;
el movimiento verifica continuidad y Doppler. Son validaciones complementarias:
una concordancia de potencia en una escena extensa no reemplaza las comprobaciones
de fase y derivada en configuraciones controladas \citep{Molisch2022,ProakisSalehi2008}.

\section{Síntesis: cantidades que reaparecerán en la tesis}

Este capítulo fija las variables que sostienen el resto del manuscrito. La longitud de onda de la \cref{eq:wavelength} determina sensibilidad de fase; la respuesta LTV de la \cref{eq:ltv_channel} fundamenta la dinámica; $\tau_{\rm rms}$ se relaciona con ancho de coherencia, mientras el ancho de señal $B$ fija la resolución nominal $1/B$; el modelo por trayectorias de la \cref{eq:multipath} sirve de interfaz para el gemelo; y el modelo instrumental de la \cref{eq:hardware_model} separa mundo de medición. Los capítulos siguientes no redefinen estas cantidades: las especializan para calibración, Wi-Fi, LoRa, respiración e información de Fisher.


\section{Fase de referencia e invariantes espaciales y espectrales}
\label{sec:phase_invariants}

Una respuesta compleja incluye una referencia de fase. La diferencia entre
simulación y medición puede contener una transformación común al arreglo,
\begin{equation}
\vect h'_u(\nu_k)=e^{j\psi_u-j2\pi\nu_k\delta\tau_u}\vect h_u(\nu_k),
\qquad \vect h_u(\nu_k)\in\mathbb C^{A}.
\label{eq:array_phase_group}
\end{equation}
El término común no representa necesariamente un error puramente instrumental:
una modificación física también puede producir una rotación compartida. Su
tratamiento como perturbación exige declarar qué variable física se desea
preservar. La distinción es coherente con la separación entre propagación,
canal equivalente y referencia de recepción en comunicaciones
\citep{Molisch2022,Bello1963}.

El producto exterior espacial
\begin{equation}
\mat G_u(k)=\vect h_u(\nu_k)\vect h_u(\nu_k)^H
\label{eq:spatial_gram_invariant}
\end{equation}
es hermítico, semidefinido positivo y de rango uno para un vector no nulo.
Es un producto instantáneo; sólo tras una esperanza o promedio con supuestos
adecuados puede interpretarse como matriz de segundo momento, y no es una
covarianza centrada por definición. Sus diagonales son potencias y sus términos
fuera de diagonal contienen productos relativos entre antenas.

La transformación de la \cref{eq:array_phase_group} satisface
\begin{equation}
\mat G'_u(k)=\mat G_u(k)
\end{equation}
para cualquier $\psi_u$ y $\delta\tau_u$. De hecho, este observable es invariante
a una fase arbitraria en cada frecuencia. Por ello permite evaluar estructura
espacial relativa, pero elimina también el retardo común: comparar $G$ antes
y después de corregir ese retardo no es una ablación informativa. Diferencias
numéricas superiores a la tolerancia indicarían cambios de máscara, ponderación
o una corrección que no es realmente común al arreglo.

Para preservar fase relativa entre frecuencias se utiliza
\begin{equation}
\mat K_u(k,l)=\vect h_u(\nu_k)\vect h_u(\nu_l)^H,
\qquad
\mat K'_u(k,l)=e^{-j2\pi(\nu_k-\nu_l)\delta\tau_u}\mat K_u(k,l).
\label{eq:cross_frequency_invariant}
\end{equation}
$K$ elimina $\psi_u$ y conserva sensibilidad al retardo cuando $\nu_k\ne\nu_l$.
La interpretación requiere que $\psi_u$ sea constante sobre las frecuencias
comparadas. Los productos entre arreglos distintos conservan la fase relativa
$\psi_u-\psi_v$; no son invariantes a rotaciones independientes de ambos arreglos.

Una alternativa compacta apila antenas y frecuencias válidas en
$\vect x_u=\operatorname{vec}(H_u)$. Para un canal medido $\vect x$ y uno
predicho $\vect y$, la discrepancia respecto de una sola fase común es
\begin{equation}
d_{\mathrm U(1)}^2(\vect x,\vect y)
=\min_{\psi}\|\vect x-e^{j\psi}\vect y\|_2^2
=\|\vect x\|_2^2+\|\vect y\|_2^2-2|\vect x^H\vect y|.
\label{eq:phase_quotient_distance}
\end{equation}
Cuando el producto interior no es nulo, la rotación aplicada a $\vect y$ es
$\psi^\star=-\angle(\vect x^H\vect y)$. El valor mínimo sigue definido si el
producto se anula, aunque la fase óptima deje de ser única. Esta distancia
no ajusta fases independientes por antena ni por frecuencia, y no predice
la fase absoluta: evalúa una clase de equivalencia expresamente declarada.

Para $\vect x$ y $\vect y$ normalizados a norma uno,
$\|\vect x\vect x^H-\vect y\vect y^H\|_F^2
=2(1-|\vect x^H\vect y|^2)$.
Esta identidad permite evaluar el producto exterior conjunto sin construir
una matriz de tamaño $AK\times AK$, con $(AK)^2$ entradas. La normalización elimina también escala,
por lo que debe acompañarse de métricas de potencia; el error sin normalizar
conserva esa discrepancia. Todas las comparaciones usan un soporte y pesos
fijados antes de evaluar el método, con tratamiento explícito de energía nula.

\section{Escalas geométricas y sensibilidad de fase}
\label{sec:geometry_phase_scales}

Para una trayectoria de longitud total $d(\mathbf q)$ y coeficiente de interacción
$\alpha(\mathbf q,f)$, la contribución al canal es
$H_p=\alpha\exp[-j2\pi d/\lambda]$. Por tanto,
\begin{equation}
\nabla_{\mathbf q}H_p=e^{-j2\pi d/\lambda}
\left(\nabla_{\mathbf q}\alpha-j\frac{2\pi}{\lambda}\alpha\nabla_{\mathbf q}d\right),
\qquad
\delta\phi_p\simeq-\frac{2\pi}{\lambda}\nabla_{\mathbf q}d^{\mathsf T}\delta\mathbf q.
\label{eq:geometry_phase_jacobian}
\end{equation}
La derivada incluye cambios de amplitud y de fase de reflexión, además de longitud.
En un reflector puntual biestático,
$d=\|\mathbf q-\mathbf p_T\|+\|\mathbf q-\mathbf p_R\|$; su gradiente es la
suma de los vectores unitarios desde ambos terminales hacia el reflector. El factor
$4\pi/\lambda$ corresponde al desplazamiento radial monostático, no a cualquier
movimiento en un enlace Wi-Fi. La orientación y la geometría pueden anular la
sensibilidad de primer orden aun cuando la potencia recibida sea elevada.

Los valores del \cref{tab:wavelength_metrology} no se extraen de una tabla
experimental: se calculan con
\begin{equation}
\lambda=\frac{c}{f},\qquad
|\delta\phi_p|_{\rm deg}=360^{\circ}\frac{|\delta d|}{\lambda},
\qquad c=299\,792\,458\ \mathrm{m/s},
\label{eq:phase_length_conversion}
\end{equation}
que es la conversión directa entre longitud eléctrica y fase
\citep{Molisch2022,Rappaport2002}. La frecuencia de 3.438 GHz corresponde a dc01
\citep{Euchner2023DICHASUSdcxx}; 5 GHz representa una banda Wi-Fi de trabajo, y
28/60 GHz sirven como escalas mmWave relacionadas con la extensión experimental
propuesta \citep{Yao2026RFDTChannel,TIIWRL6432BOOST}. Los cambios de longitud de
10 mm y 1 mm son ejemplos metrológicos deliberados, no errores medidos del banco.

\begin{table}[htbp]
\centering
\caption{Sensibilidad geométrica aproximada calculada con la
\cref{eq:phase_length_conversion}. $|\delta d|$ es el cambio de longitud total de
trayectoria, no necesariamente el desplazamiento del objeto.}
\label{tab:wavelength_metrology}
\begin{tabular}{rrrr}
\toprule
Frecuencia (GHz) & $\lambda$ (mm) & $|\delta d|$ (mm) & $|\delta\phi_p|$ (grados)\\
\midrule
3.438 & 87.2 & 10 & 41.3\\
5 & 60.0 & 10 & 60.0\\
28 & 10.7 & 1 & 33.6\\
60 & 5.0 & 1 & 72.0\\
\bottomrule
\end{tabular}
\end{table}

A 60 GHz, un milímetro radial en configuración monostática produce aproximadamente
144 grados. Estos valores explican por qué una nube de puntos visualmente convincente
puede resultar insuficiente para predecir fase. Para incertidumbre geométrica pequeña
$\delta\mathbf q$ con covarianza $\Sigma_q$, manteniendo fija la fase del coeficiente
de interacción, la contribución de longitud de una trayectoria satisface
\begin{equation}
\sigma_{\phi_p}^2\simeq\left(\frac{2\pi}{\lambda}\right)^2
\nabla d^{\mathsf T}\Sigma_q\nabla d.
\label{eq:geometric_phase_uncertainty}
\end{equation}
La aproximación exige una familia de trayectorias estable. Cerca de oclusiones,
transiciones de visibilidad o nulos por interferencia, se utilizarán perturbaciones
geométricas Monte Carlo y métricas del canal complejo; la fase de una suma casi nula
es inestable y no hereda directamente la varianza de una trayectoria aislada.

\part{Gemelos digitales y representaciones físico--latentes}
\chapter{Gemelos digitales inalámbricos y trazado de rayos diferenciable}
\label{ch:wdt}

El capítulo anterior estableció la cadena que conduce desde la propagación hasta la
medición de canal. Sobre esa base, este capítulo introduce el gemelo digital como una
realización computacional de dicha cadena. El interés no reside únicamente en reproducir
el canal observado, sino en conservar la relación entre los parámetros físicos de la escena
y las variaciones de la señal que posteriormente se utilizarán para inferir movimiento y
fisiología.

\section{De la simulación de cobertura al gemelo inalámbrico}

Un simulador de propagación predice $H$ a partir de una escena; un gemelo digital persigue, además, correspondencia con un entorno físico concreto, actualización con mediciones y capacidad de responder a cambios. En el contrato de la \cref{eq:thesis_chain}, no basta con generar mapas plausibles: el modelo debe declarar qué escena, frecuencia, antena y observable reproduce. La formulación mínima es
\begin{equation}
\widehat H_m(f,t)=F_m
\left(\Gamma(t),\vect p_T,\vect p_R,f;
\vect s_R,\vect\theta_E\right),
\end{equation}
donde $\theta_E$ incluye parámetros materiales, dispersión, patrones de antena y otros parámetros físicos no conocidos exactamente.

El reto es que una reconstrucción geométrica visualmente correcta puede producir errores de fase considerables. A frecuencia portadora $f_c$, un error de longitud $\delta L$ produce
\begin{equation}
\delta\phi=-2\pi\frac{\delta L}{\lambda}.
\end{equation}
Por tanto, los errores centimétricos o incluso milimétricos pueden ser importantes para el sensado de movimientos pequeños. Esta observación motiva dos líneas complementarias: la calibración diferenciable de propiedades electromagnéticas y la calibración que considera errores locales de fase.

\section{Calibración como problema inverso}

Sea $\vect\theta$ el vector de parámetros del gemelo y $\mathcal D_{\rm cal}$ un conjunto de mediciones alineadas. Se define
\begin{equation}
\vect\theta^\star=\arg\min_{\vect\theta}\;
\sum_{i\in\mathcal D_{\rm cal}}
\calL\left(
H_{m,i}^{\rm real},
F_m(\Gamma_i,\vect p_{T,i},\vect p_{R,i},f;
\vect s_{R,i},\vect\theta)
\right)+\beta R(\vect\theta).
\label{eq:calibration_inverse}
\end{equation}
Para una pérdida real se vectorizan parte real e imaginaria, $\vect u=[\Re H;\Im H]$, y la diferenciabilidad permite obtener
\begin{equation}
\nabla_{\vect\theta}\calL
=\left(\frac{\partial\vect u}{\partial\vect\theta^T}\right)^T
\nabla_{\vect u}\calL.
\end{equation}
Trabajos recientes de calibración con \ac{DRT} muestran que las propiedades materiales, patrones de antena y parámetros de dispersión pueden optimizarse utilizando canales reales \citep{Hoydis2024DiffRT}; otros introducen errores locales de fase para modelar discrepancias geométricas pequeñas que una calibración de potencia no puede absorber \citep{Ruah2024Phase}. En esta tesis se consideran ambas familias como componentes consecutivos, no sustitutos.

\section{Separación de fuentes de error}

Conviene escribir
\begin{equation}
H_{\rm real}-H_{\rm gemelo}=
\Delta H_{\rm geom}+
\Delta H_{\rm mat}+
\Delta H_{\rm scat}+
\Delta H_{\rm ant}+
\Delta H_{\rm hw}+
\epsilon_{\rm model}.
\label{eq:domain_gap_decomp}
\end{equation}
Esta suma es una contabilidad conceptual o una linealización local, no una descomposición identificable a partir de $H$ por sí solo. Un modelo flexible todavía puede usar materiales para compensar coordenadas. Por ello la calibración será jerárquica: registro y sincronización; antenas y escala; materiales y scattering; y sólo al final residuales de fase/modelo. Priors, mediciones auxiliares e intervenciones separadas anclarán cada bloque.

\section{Jerarquía de la literatura y de la evidencia}
\label{sec:literature_evidence_hierarchy}

Los trabajos revisados no se tratarán como soluciones intercambiables. La tesis los
ordena por el operador que modelan, la información que reciben y la evidencia con que
se validan:
\begin{table}[htbp]
\centering\small
\caption{Función científica de la literatura en la cadena metodológica.}
\label{tab:literature_evidence_hierarchy}
\begin{tabularx}{\textwidth}{P{2.6cm}L L}
\toprule
Bloque & Referencias y uso en la tesis & Límite de inferencia\\
\midrule
Calibración estática & Sionna RT, calibración diferenciable y error local de fase
\citep{Hoydis2023SionnaRT,Hoydis2024DiffRT,Ruah2024Phase}. & Mejorar potencia o una
verosimilitud con fase latente no prueba geometría verdadera ni sensado dinámico.\\
Campo y objetos & NeRF$^2$, Learnable WDT, RFScape y RadTwin
\citep{Zhao2023NeRF2,Jiang2025LearnableWDT,Chen2025RFScape,Zhang2026RadTwin}. &
Interpolación de campo, edición de objetos y generalización geométrica son tareas
distintas y requieren entradas equivalentes para compararse.\\
Dinámica simulada & WaveVerse \citep{Zheng2026WaveVerse}. & La coherencia de una
simulación y sus contrastes externos no validan el Jacobiano de nuestro banco ni una
forma de onda respiratoria propia.\\
Inversión & RayLoc, RFDT e InverTwin
\citep{Han2026RayLoc,Chen2026RFDT,Chen2025InverTwin}. & La inversión se evaluará
después de validar el modelo directo, su referencia y sus discontinuidades relevantes.\\
Adquisición y robustez & FarSense, HiSAC, RadioRange y OpenLoRa
\citep{Zeng2019FarSense,Pegoraro2024HiSAC,Lyu2026RadioRange,Mishra2023OpenLoRa}. &
Cocientes, anclas o perturbaciones simuladas dependen de sus supuestos de hardware;
no acreditan coherencia del equipo propio.\\
Representación humana & BreatheSmart, WiMANS, WiMUSE, MURAL-Fi, MM-Fi y WiFi-JEPA
\citep{Mosleh2022BreatheSmart,Huang2024WiMANS,Yi2025WiMUSE,Li2026MURALFI,
Yang2023MMFi,Kim2026WiFiJEPA}. & Una tarea humana bien resuelta no valida por sí
misma la geometría, la fase absoluta ni el gemelo electromagnético.\\
\bottomrule
\end{tabularx}
\end{table}

Esta jerarquía determina el orden de ejecución. DICHASUS, Hoydis y Ruah fijan el
problema del modelo estático; los invariantes y las referencias instrumentales delimitan
el observable; las representaciones de escena se comparan después bajo regímenes de
información comunes; el reflector y el fantoma contrastan fidelidad diferencial; y
la inversión y el aprendizaje humano se activan cuando el observable ha demostrado
sensibilidad. La secuencia no resta valor a los métodos pospuestos: evita usar éxito
en una tarea como evidencia de otra.

\section{DICHASUS dc01 como banco de pruebas del gemelo estático}

La primera validación desarrollada utiliza DICHASUS dc01 y la escena correspondiente empleada en trabajos de calibración diferenciable \citep{Euchner2023DICHASUSdcxx,Hoydis2024DiffRT}. El banco reportado comprende $N=10000$ posiciones y conserva el objeto experimental
\begin{equation}
\mathcal D_{\rm dc01}
=\left\{
\vect p_i,\;H_{m,i}^{\rm real},\;\calP_{m,i}^{\rm RT},\;H_{m,i}^{\rm RT}
\right\}_{i=1}^{N}.
\label{eq:dc01_dataset}
\end{equation}
El soporte geométrico y el conjunto de trayectorias radio--dependiente se conservan como interfaz física, en lugar de almacenar únicamente el canal sintetizado. Esto permite modificar posteriormente materiales, frecuencia, observación y dinámica sin confundir el mundo común $\Gamma$ con una realización particular $\calP_m$.

La comparación B0--B2 y los diagnósticos de fase se documentan en el \cref{ch:avances}. DICHASUS es un banco de sondeo OFDM: su respuesta en frecuencia resulta adecuada para estudiar la propagación, mientras la validación del operador de observación de dispositivos Wi-Fi comerciales requiere pruebas específicas. Conservar la geometría tampoco autoriza reutilizar sin cambios las ganancias y trayectorias efectivas al pasar a LoRa.

\section{Fidelidad de campo y fidelidad diferencial}

Para comunicaciones suele ser suficiente comparar $H_{\rm sim}$ y $H_{\rm real}$. Para el sensado se requiere una condición adicional. Si $q$ es una variable física de interés,
\begin{equation}
J_q=\frac{\partial H}{\partial q},
\end{equation}
y se propone medir
\begin{equation}
E_J=
\frac{\|J_q^{\rm model}-J_q^{\rm ref}\|_F}
{\|J_q^{\rm ref}\|_F+\epsilon}.
\label{eq:jacobian_error}
\end{equation}
Un modelo con bajo \ac{NMSE} de canal, pero con derivadas incorrectas respecto del movimiento humano, puede ser inadecuado para el sensado. Por ello, la tesis distingue entre \emph{fidelidad del modelo directo} y \emph{fidelidad diferencial}.

$J_q^{\rm ref}$ se estimará con perturbaciones mecánicas bidireccionales, paso declarado, repeticiones y alineamiento de fase/ganancia definido antes de observar el resultado. Como el error relativo de la \cref{eq:jacobian_error} es inestable cuando $\|J_q^{\rm ref}\|$ es cercano a cero, se reportará también error absoluto con unidades, correlación compleja y error ponderado por la tarea. El término $\epsilon$ se fijará en las mismas unidades que la norma de referencia y no como constante adimensional arbitraria.

\section{Reconstrucción métrica y registro del gemelo}
\label{sec:geometric_metrology}

La geometría debe representarse con unidades, sistema de coordenadas e incertidumbre
verificables. Se distinguen tres objetivos de diseño: escala global de paredes,
puertas y mobiliario del orden de 1--2 cm; registro de antenas del orden de 2--5 mm;
y referencia local de desplazamiento del orden de 0.1--1 mm. Son rangos de trabajo
que deberán demostrarse mediante controles independientes, no especificaciones
atribuidas a cualquier LiDAR, cámara o actuador. Un experimento con pasos de
0.25 mm sólo es interpretable si su incertidumbre efectiva permite resolverlos.

Se fijará un marco global $W$ y se almacenará, para cada sensor $s$,
\begin{equation}
T_{W\leftarrow s}=\begin{bmatrix}R_{Ws}&\mathbf t_{Ws}\\0&1\end{bmatrix}\in SE(3),
\qquad \widetilde{\mathbf p}_W=T_{W\leftarrow s}\widetilde{\mathbf p}_s,
\label{eq:scene_extrinsics}
\end{equation}
con $R_{Ws}\in SO(3)$. Las transformaciones compuestas deberán respetar esa dirección.
Se registrarán extrínsecos de cámara, LiDAR, radar, transmisores, receptores y antenas;
el punto de montaje no se confundirá con el centro de fase, que puede depender de
frecuencia y dirección. Marcadores rígidos, distancias de control y mediciones repetidas
permitirán verificar el registro. Un residual pequeño de ICP mide concordancia entre
nubes, pero no garantiza exactitud absoluta. Se reservarán, cuando el recinto lo permita,
10--20 distancias de control que no se hayan utilizado para ajustar la reconstrucción.

La comparación propuesta distinguirá cinco condiciones:
\begin{description}
\item[G0.] Planta manual y primitivas geométricas con dimensiones verificadas.
\item[G1.] Reconstrucción RGB-D, con registro y simplificación documentados.
\item[G2.] Reconstrucción LiDAR con las mismas referencias globales.
\item[G3.] Geometría LiDAR regularizada para RF: superficies, espesores, huecos y
objetos relevantes corregidos; materiales nominales y presupuesto RT controlados.
\item[G4.] Geometría anterior con calibración RF de materiales y, si se justifica,
parámetros geométricos restringidos por su incertidumbre metrológica.
\end{description}
G4 cambia más factores que la geometría. Por ello se añadirá una ablación factorial
con materiales nominales/calibrados para separar el efecto del registro, la regularización
y el ajuste electromagnético. Todos los ajustes utilizarán sólo entrenamiento y
validación; los puntos reservados medirán generalización, no servirán para corregir
la escena a posteriori. Se mantendrán mecanismos RT, frecuencias y presupuesto de
muestreo comparables, o se informará explícitamente su coste.

RFDT-Channel ofrece una referencia reciente para construir escenas RF a partir de
RGB y LiDAR y simular canales interiores a 28 GHz \citep{Yao2026RFDTChannel}.
El gemelo de Wang y colaboradores \citep{Wang2026WiFiTwin} contrasta predicción de señal con RSSI real.
Ambos orientan el flujo geométrico, pero una validación de potencia o un canal simulado
no establecen exactitud de fase medida ni fidelidad de micromovimiento de este banco.
La comparación G0--G4 se evaluará sucesivamente en error geométrico, potencia,
retardos, canal relativo y respuesta diferencial. No se exigirá resolver primero
la fase absoluta de toda la habitación para iniciar una prueba local controlada.

Esta distinción cierra el paso entre el modelado electromagnético y la representación: una
aproximación útil debe conservar no sólo valores de canal, sino también las variaciones
que contienen información acerca del estado humano. El capítulo siguiente examina qué
estructuras explícitas, implícitas o híbridas pueden satisfacer ese requisito.

\chapter{Representaciones neuronales del entorno y del campo de radiofrecuencia}
\label{ch:latent_env}

Una vez definido el criterio de fidelidad del gemelo, es necesario decidir cómo representar
la escena y el campo de radiofrecuencia. Esta decisión afecta la capacidad de edición del
modelo, el costo computacional, la transferencia entre escenas y la cantidad de mediciones
requerida para la calibración. Por ello, las representaciones neuronales se estudian aquí
como aproximaciones estructuradas del modelo físico y no como sustitutos desligados de él.

\section{Por qué representar y no sólo simular}

El trazado de rayos explícito conserva la interpretabilidad, pero puede ser costoso y depende de una geometría detallada. Los modelos neuronales del campo buscan aproximar
\begin{equation}
(\vect p_T,\vect p_R,G_E)\mapsto H
\end{equation}
a partir de mediciones o representaciones geométricas compactas. La cuestión científica no es únicamente qué modelo reduce el error cuadrático, sino qué representación conserva la estructura necesaria para interpolar, extrapolar, editar escenas, adaptarse con pocas mediciones y soportar inversión.

\section{Regímenes de información}

Para comparar familias heterogéneas se fija primero una consulta $\vect u=(\vect p_T,\vect p_R,f,\vect s_R)$ y un conjunto de pares RF de calibración $\mathcal D_{\rm RF}=\{(\vect u_i,H_i)\}$. Los cuatro regímenes indican información adicional disponible en inferencia:
\begin{align}
\mathcal R_0 &: \{\vect u,\mathcal D_{\rm RF}\},\\
\mathcal R_1 &: \{\vect u,\mathcal D_{\rm RF},G_E^{\rm geom}\},\\
\mathcal R_2 &: \{\vect u,\mathcal D_{\rm RF},G_E^{\rm geom},\{O_j\}\},\\
\mathcal R_3 &: \{\vect u,\mathcal D_{\rm RF},G_E^{\rm geom},\{O_j\},
\vect\theta_{\rm mat},\vect\theta_{\rm ant}\}.
\end{align}
Aquí $G_E^{\rm geom}$ contiene sólo geometría, a diferencia del estado
completo $G_E$ que también incluye materiales. $O_j$ es un objeto segmentado con geometría y clase semántica. El canal reservado para evaluación es siempre la variable que debe predecirse y no una entrada de la consulta. El número de muestras, la región espacial y los metadatos se mantendrán constantes dentro de cada comparación; así se evita presentar como ganancia arquitectónica una ventaja debida a la información disponible.

\section{Campos implícitos}

Un campo neural continuo puede aproximarse como
\begin{equation}
\widehat H=F_\theta(\vect p_T,\vect p_R,f).
\end{equation}
La ventaja es una representación continua y compacta; el riesgo es que la geometría y los materiales queden absorbidos en los pesos, lo que dificulta la edición y la transferencia. NeRF$^2$ representa una línea importante en esta dirección y se propone en S5 como método de referencia para campos implícitos \citep{Zhao2023NeRF2}.

\section{Representaciones centradas en objetos e informadas por la geometría}

Una alternativa factoriza el entorno:
\begin{equation}
Z_E=\left\{z_{\rm global},z_{O_1},\ldots,z_{O_M}\right\},
\end{equation}
con
\begin{equation}
z_{O_i}=\left[z_{\rm geom}^{(i)},z_{\rm mat}^{(i)},z_{\rm int}^{(i)}\right].
\end{equation}
Learnable WDT y trabajos condicionados a la geometría como RadTwin constituyen referencias para esta familia \citep{Jiang2025LearnableWDT,Zhang2026RadTwin}. RFScape constituye otro antecedente directo de objetos neuronales acoplados
a trazado editable \citep{Chen2025RFScape}. La hipótesis es que una representación explícitamente estructurada puede necesitar más información geométrica inicial, pero generalizar mejor a cambios de objetos y soportar intervenciones físicas interpretables.

\section{El espacio de trayectorias como representación intermedia}

Se propone representar cada trayectoria mediante un descriptor
\begin{equation}
\vect p_{m,\ell}=
[\Re\alpha_{m,\ell},\Im\alpha_{m,\ell},\tau_{m,\ell},
\theta^{\rm AoA}_{m,\ell},\theta^{\rm AoD}_{m,\ell},
\nu_{m,\ell},\vect c_{m,\ell}],
\end{equation}
donde $\vect c_{m,\ell}$ codifica interacciones; los ángulos se representarán mediante seno/coseno para evitar discontinuidad en $2\pi$. El conjunto es permutacionalmente invariante, por lo que una arquitectura de tipo Deep Sets puede escribirse como
\begin{equation}
z_{\calP_m}=\rho\left(\sum_{\ell}\phi(\vect p_{m,\ell})\right).
\label{eq:deepset_paths}
\end{equation}
Esta forma admite un número variable de trayectorias y mantiene la correspondencia con el modelo físico; su invariancia a permutaciones sigue la construcción de Deep Sets \citep{Zaheer2017DeepSets}. Un transformador de conjuntos o un grafo pueden extenderla cuando las relaciones entre trayectorias sean relevantes. Para el alineamiento entre tecnologías de radio se aprenderá una representación común del soporte $\Gamma$, no una igualdad artificial entre los descriptores de $\calP_m$.

\section{Hipótesis de factorización latente}

Se propone separar
\begin{equation}
Z=\{Z_E,Z_\phi,Z_R,Z_H,Z_F,Z_N\},
\label{eq:factor_latent}
\end{equation}
con $Z_E$ asociado al entorno, $Z_\phi$ a la estructura e incertidumbre de fase, $Z_R$ a la configuración de radio, $Z_H$ al macroestado humano, $Z_F$ a la fisiología y $Z_N$ a perturbaciones residuales. La factorización no supone independencia absoluta; busca que la información relevante sea manipulable y que el codificador no confunda la identidad del dispositivo con el estado humano.

La inclusión de $Z_\phi$ está motivada por los diagnósticos del
\cref{ch:avances}. S4.4--S4.5 muestran predictibilidad parcial del retardo;
S4.6 encuentra una dirección común al arreglo cuya predicción mediante vecinos
no resulta útil bajo los controles realizados. En consecuencia se propone
$Z_\phi=[Z_\phi^{\rm pred},Z_\phi^{\rm nuisance},Z_\phi^{\rm residual}]$,
según la \cref{eq:phase_latent_roles}. Esta clasificación distingue cantidades
predictivas, parámetros perturbadores y discrepancia residual; no asigna causas
físicas únicas ni presupone independencia. La elección de un observable
invariante debe conservar sensibilidad a la tarea. Su fundamento matemático
se presenta en la \cref{sec:phase_invariants}, y su validación es el bloque
propuesto S4.7a. Un codificador aprendido y la incertidumbre por trayectoria
permanecen pendientes. La clasificación se revisará si una referencia de
adquisición o metadatos adicionales modifican la predictibilidad observada.

La identificabilidad de representaciones separadas requiere restricciones e
información adicional \citep{Locatello2019Identifiability}. Esta partición no es identificable
a partir de una única observación sin supuestos adicionales: distintas transformaciones invertibles de las variables latentes pueden explicar el mismo $Y_m$. En consecuencia, cada bloque tendrá criterios observables: intervenciones simuladas que modifican un factor a la vez, etiquetas parciales, pares de la misma escena observados con radios distintos, predicción cruzada y pruebas de invariancia fuera de distribución. La tesis evaluará esas propiedades y no interpretará automáticamente cada coordenada latente como una causa física.

\section{Dos ramas complementarias de representación}
\label{sec:two_representation_branches}

Se propone conservar simultáneamente una rama robusta $z_{\rm rob}=f_{\rm rob}(Y)$
y una rama coherente $z_{\rm coh}=f_{\rm coh}(Y,Y_{\rm ref},M)$, donde $M$ incluye
metadatos de adquisición y $Y_{\rm ref}$ una referencia independiente cuando exista.
La primera describe estructura relativa estable frente a un grupo instrumental
especificado; la segunda conserva variaciones temporales de fase y amplitud necesarias
para inferir movimiento. Su fusión se condicionará por la calidad de referencia y
la observabilidad, y se comparará con cada rama por separado.

El requisito deseado es invariancia a perturbaciones instrumentales e identificación
de la respuesta a intervenciones físicas. No puede garantizarse sólo mediante una
arquitectura: si una deriva instrumental y un movimiento generan la misma rotación
común, eliminarlos algebraicamente elimina ambas contribuciones. La separación exige
referencias, intervenciones controladas o información adicional. La equivariancia
física se expresará respecto de una transformación conocida del estado; no se
supondrá que toda traslación de un objeto produce una rotación global del canal.

\subsection{Familia operacional para S4.7a}
\label{sec:representation_v_family}

Se denominarán $\mathsf V_0$--$\mathsf V_7$ las alternativas del análisis; esta
notación evita confundirlas con los regímenes de información $\mathcal R_i$ y los
hitos administrativos R1--R5. Para $H\in\mathbb C^{A\times K}$, $\mathbf h_k=H_{:,k}$
y $\epsilon>0$ con las unidades del denominador, se definen:
\begin{align}
\mathsf V_0(H)&=H, &\mathsf V_1(H)&=|H|,\nonumber\\
[\mathsf V_2^{(\ell)}(H)]_{a,k}&=\frac{H_{a,k+\ell}H_{a,k}^{*}}
{|H_{a,k+\ell}||H_{a,k}|+\epsilon},\qquad \ell\in\mathcal L,
\label{eq:v_frequency_products}\\
[\mathsf V_3(H)]_{a,b,k}&=\frac{H_{a,k}H_{b,k}^{*}}
{|H_{a,k}||H_{b,k}|+\epsilon},\label{eq:v_antenna_products}\\
\mathsf V_5(H)_k&=\frac{\mathbf h_k\mathbf h_k^{H}}
{\|\mathbf h_k\|_2^2+\epsilon},\label{eq:v_spatial_gram}\\
[\mathsf V_6(H)]_{a,k}&=H_{a,k}\exp(-j\widehat\psi+j2\pi\nu_k\widehat\tau).
\label{eq:v_canonical_channel}
\end{align}
$\mathsf V_4$ designa la comparación en el cociente por $U(1)$, mediante
$d_{U(1)}^2(x,y)=\|x\|^2+\|y\|^2-2|x^Hy|$, para el canal apilado de un arreglo.
Es una distancia entre clases de equivalencia, no un codificador unario ni una
predicción del desfase óptimo. La rotación es única para todas sus antenas y frecuencias.

$\mathsf V_2$ conserva diferencias espectrales y elimina una fase común constante;
$\mathcal L$ comenzará con $\{1,4,16,64\}$ cuando el soporte espectral lo permita,
pero su selección y las máscaras de borde se fijarán con entrenamiento/validación.
$\mathsf V_3$ y $\mathsf V_5$ eliminan además retardo común, pero no desfases
relativos independientes entre cadenas. El producto instantáneo de
\cref{eq:v_spatial_gram} es una matriz de Gram, no una covarianza temporal estimada.
La normalización con $\epsilon$ mantiene invariancia exacta a fase, pero sólo
aproxima invariancia a escala. Se declararán máscaras de desvanecimiento y pares
seleccionados antes de evaluar; el número de características y su coste también
forman parte de la comparación.

$\mathsf V_6$ requiere que $\widehat\psi$ y $\widehat\tau$ se obtengan de una
referencia disponible causalmente, metadatos o un predictor congelado. En dc01,
la corrección de fase común ajustada frente al canal objetivo sigue siendo oráculo.
Calcular productos del CSI recibido como entrada de un sensor es admisible;
alinear una predicción RT con su objetivo de prueba y llamarla predicción autónoma
no lo es. Ambas operaciones deben distinguirse también al entrenar modelos latentes.

$\mathsf V_7$ designará la fusión de magnitud, productos espectrales, productos
espaciales y canal canónico disponibles. No se definirá como concatenación fija antes
del experimento: cada ablación declarará componentes, dimensión, normalización y acceso
a referencia. La comparación con $\mathsf V_1$--$\mathsf V_6$ determinará si la
complementariedad compensa el aumento de capacidad y de información auxiliar.

Las representaciones anteriores conectan escena, estructura relativa y referencia de
adquisición. El \cref{ch:learning_foundations} establece cómo aprenderlas y evaluarlas;
el capítulo de observación especifica después qué información conserva cada radio.
Esta separación permite controlar arquitectura, adquisición y fidelidad física como
factores distintos de la comparación.

\part{De la escena a la observación y la inferencia}
\chapter{Fundamentos estadísticos, aprendizaje y optimización del gemelo}
\label{ch:learning_foundations}

La calibración electromagnética y la predicción de sus discrepancias constituyen problemas
inversos relacionados, pero distintos. En el primero se ajustan parámetros del simulador;
en el segundo se aprende una función que aproxima un descriptor de error a partir de
información disponible en una consulta nueva. Este capítulo establece los fundamentos
necesarios para interpretar las regresiones de S4.4, los análisis agrupados de S4.3, la caracterización circular de S4.6 y la
extensión posterior a representaciones profundas. La formulación probabilística sigue
la distinción entre modelo, inferencia y decisión de Bishop, mientras el tratamiento de
redes profundas y regularización se apoya en Goodfellow et al.\ \citep{Bishop2006PRML,Goodfellow2016DL}.

\section{Aprendizaje supervisado y objetivo de generalización}

Sea $\vect z_i$ un descriptor disponible en inferencia y $y_i$ una variable objetivo.
Para S4.4, $\vect z_i$ contiene posición y/o descriptores del trazado nominal, mientras
$y_i$ es un retardo efectivo obtenido durante la construcción de las etiquetas a partir
del canal medido. Una función $g_\theta$ se estima mediante
\begin{equation}
\widehat\theta=\arg\min_\theta
\frac{1}{N_{\rm ent}}\sum_{i\in\mathcal I_{\rm ent}}
\ell\bigl(y_i,g_\theta(\vect z_i)\bigr)+\lambda\Omega(\theta).
\label{eq:empirical_risk}
\end{equation}
El riesgo de interés es, en cambio,
\begin{equation}
\mathcal R_Q(g)=\mathbb E_{(\vect z,y)\sim Q}
\left[\ell\bigl(y,g(\vect z)\bigr)\right],
\end{equation}
donde $Q$ representa la distribución de despliegue. Una partición aleatoria estima
principalmente un riesgo bajo condiciones semejantes a las observadas. Una región
reservada modifica la distribución de consultas y exige estudiar extrapolación espacial.
La separación entre ambos riesgos es central para interpretar S4.4d--f.

Construir las etiquetas con mediciones es compatible con aprendizaje supervisado.
La restricción consiste en impedir que el descriptor, la normalización o la selección
de hiperparámetros utilicen las mediciones de prueba. Cuando el descriptor proviene
de un simulador calibrado, sus parámetros también deben entrenarse sin las regiones
reservadas; alternativamente, debe declararse un protocolo con calibración previa común
y una afirmación de generalización condicionada a ella.

\section{Referencias de regresión y complejidad del predictor}

\subsection{Constante, vecino próximo y regresión regularizada}

Una constante ajustada por error absoluto utiliza la mediana de entrenamiento; por error
cuadrático utiliza la media. Esta elección debe distinguirse del uso arbitrario de cero.
El vecino próximo predice
\begin{equation}
\widehat y(\vect p)=y_{i^\star},\qquad
i^\star=\arg\min_{i\in\mathcal I_{\rm ent}}\|\vect p-\vect p_i\|_2.
\end{equation}
Es una referencia necesaria para determinar cuánto aporta un predictor frente a la
proximidad espacial. No estima un mecanismo electromagnético: interpola una etiqueta
en la geometría observada.

Para descriptores centrados y escalados con estadísticos de entrenamiento, la regresión
ridge resuelve
\begin{equation}
(\widehat b,\widehat{\vect w})=
\arg\min_{b,\vect w}\frac{1}{N_{\rm ent}}
\|\vect y-b\vect 1-\mat Z\vect w\|_2^2+\lambda\|\vect w\|_2^2.
\label{eq:ridge_model}
\end{equation}
Con datos centrados, $\widehat{\vect w}=(\mat Z^T\mat Z+
N_{\rm ent}\lambda\mat I)^{-1}\mat Z^T\vect y$ para $\lambda>0$.
El intercepto no se penaliza. La regularización estabiliza direcciones correlacionadas;
no corrige por sí sola una relación no lineal ni un cambio de dominio
\citep{Bishop2006PRML,Boyd2004Convex}.

\subsection{Bosques aleatorios y redes neuronales}

Un bosque de regresión combina árboles entrenados con aleatorización de muestras y
variables \citep{Breiman2001RandomForests}:
\begin{equation}
g_{\rm bosque}(\vect z)=\frac{1}{B}\sum_{b=1}^{B}T_b(\vect z).
\end{equation}
En las tablas experimentales se conserva la abreviatura RF del registro de ejecución
para \emph{Random Forest}; fuera de esas tablas RF designa radiofrecuencia. El método
puede representar interacciones no lineales entre descriptores sin exigir un modelo
paramétrico de su dependencia. En regresión convencional sus predicciones promedian
valores de las hojas, por lo que no debe suponerse extrapolación funcional ilimitada
fuera del soporte de entrenamiento. La importancia de una variable tampoco identifica
su efecto causal, especialmente cuando varios descriptores están correlacionados.

Una red multicapa compone transformaciones afines y no linealidades. Para materiales,
la salida puede parametrizarse mediante
\begin{equation}
\epsilon_r(\vect x)=1+\operatorname{softplus}(a_\psi(\vect x)),\qquad
\sigma(\vect x)=\operatorname{softplus}(b_\psi(\vect x)),
\label{eq:material_constraints}
\end{equation}
como una restricción suficiente para una clase de dieléctricos pasivos no magnéticos
considerados en el modelo. La cota $\epsilon_r\geq1$ es una elección de esa clase,
no una propiedad universal de todo medio electromagnético. Estas expresiones describen
una parametrización admisible; no atribuyen retrospectivamente esa implementación a B2.
Para una representación con frecuencia variable deben comprobarse además dispersión,
validez del modelo material y consistencia entre bandas.

La codificación posicional amplía las escalas espaciales representables, pero puede
facilitar la memorización. La capacidad, regularización y distribución de las mediciones
deben evaluarse conjuntamente. El número de parámetros, por sí solo, no determina la
fidelidad física ni la generalización \citep{Goodfellow2016DL}.

\section{Optimización diferenciable, restricciones e identificabilidad}

Para un objetivo real sobre canales complejos puede trabajarse con
$\vect u=[\Re H;\Im H]$. La regla de la cadena del \cref{ch:wdt} produce gradientes
respecto de parámetros reales sin tratar una pérdida no holomorfa como función analítica
compleja. La optimización del trazado se efectúa sobre un conjunto de caminos válido
en una región de parámetros; cambios de visibilidad pueden introducir discontinuidades.

Ridge con penalización positiva constituye un problema convexo; calibrar materiales,
geometría y redes neuronales conjuntamente no lo es en general. Por ello deben
distinguirse convergencia numérica, unicidad del mínimo e identificabilidad física
\citep{Boyd2004Convex,KaipioSomersalo2005}. Una pérdida pequeña puede coexistir con
parámetros diferentes que producen canales casi indistinguibles. Las múltiples
inicializaciones y la evaluación reservada caracterizan estabilidad empírica, pero no
demuestran unicidad global.

La pérdida de ganancia y dispersión de retardos de la calibración de referencia puede
formularse con un error relativo simétrico para cantidades no negativas,
\begin{equation}
\ell_{\rm sym}(a,b)=\frac{2|a-b|}{a+b+\varepsilon_a},\qquad
\mathcal L_{\rm cal}=\lambda_P\ell_{\rm sym}(P,\widehat P)
+\lambda_\tau\ell_{\rm sym}(\tau_{\rm RMS},\widehat\tau_{\rm RMS}).
\label{eq:calibration_macro_loss}
\end{equation}
Cada estabilizador debe tener las unidades de su denominador. Esta familia de objetivos
explica por qué mejorar magnitudes macroscópicas no impone concordancia de fase;
la forma y pesos exactos de una ejecución deben proceder de su configuración
\citep{Hoydis2024DiffRT}. Añadir términos complejos o diferenciales modifica el problema
de estimación y requiere una nueva comparación, con acceso a datos equivalente.

En optimización multiobjetivo, GradNorm ajusta pesos según magnitudes de gradiente y
ritmos de aprendizaje; PCGrad modifica componentes de gradientes conflictivos
\citep{Chen2018GradNorm,Yu2020PCGrad}. Ninguno garantiza una solución global, una
frontera de Pareto completa ni corrección del modelo físico. Su aplicación se condiciona
a observar un conflicto relevante entre objetivos previamente definidos.

\section{Estadística circular de los errores de fase}
\label{sec:circular_foundations}

Las diferencias angulares se definen en el círculo y no sobre una recta. Para residuos
$r_k=\wrap_{(-\pi,\pi]}(\phi_k-\widehat\phi_k)$, se consideran
\begin{equation}
\overline z=\frac{1}{K}\sum_k e^{jr_k},\qquad
\overline R=|\overline z|,\qquad V_{\rm circ}=1-\overline R.
\end{equation}
La dirección media es $\angle\overline z$ cuando la resultante permite definirla de
forma estable. Una media aritmética de ángulos próximos a $-\pi$ y $\pi$ puede ser
engañosa; las cantidades circulares evitan esa discontinuidad
\citep{Mardia2000Directional}.

La estructura espectral puede describirse mediante
\begin{equation}
C_\phi(L)=\left|\frac{1}{|\mathcal K_L|}
\sum_{k\in\mathcal K_L}e^{j(r_{k+L}-r_k)}\right|,
\label{eq:circular_lag}
\end{equation}
donde $\mathcal K_L$ contiene pares válidos separados por $L$ índices de frecuencia.
Esta resultante de incrementos no es una correlación de Pearson centrada: una fase
residual independiente pero concentrada también puede producir un valor positivo.
Para contrastar independencia se compara con un control de permutación que conserve
la distribución marginal y las máscaras. La concentración elevada a retardos cortos
indica regularidad local; no demuestra que el residual sea determinista ni identifica
su causa física. El ajuste previo de retardo y CPO puede introducir dependencia y
debe aplicarse igualmente al control.

La distribución de von Mises,
\begin{equation}
p(r\mid\mu,\kappa)=\frac{e^{\kappa\cos(r-\mu)}}{2\pi I_0(\kappa)},
\end{equation}
permite representar concentración circular. En una predicción de fase, las salidas
seno y coseno, o una distribución circular completa, evitan una pérdida euclídea
inadecuada cerca del corte angular. La concentración no se interpreta como incertidumbre
calibrada sin comprobar su cobertura sobre observaciones reservadas.

\section{Descriptores de trayectorias y dependencia entre niveles}
\label{sec:multilevel_foundations}

Para energías no negativas $e_\ell$ de caminos activos se define
$p_\ell=e_\ell/\sum_r e_r$ si la suma es positiva. La concentración de energía se
caracteriza con $D_1=\max p_\ell$, fracciones de los caminos dominantes,
$N_{\rm eff}=1/\sum p_\ell^2$ y entropía normalizada
$\mathcal H_p=-\sum p_\ell\log p_\ell/\log L$ para $L>1$.
Se usa $0\log0=0$; para un único camino se adopta entropía cero, y sin energía válida
el descriptor se marca como no definido. Estas probabilidades describen energías
individuales y no eliminan los términos de interferencia coherente del canal total.

Sean $x_{i,a}$ y $y_{i,a}$ descriptores de una posición $i$ y enlace $a$.
La asociación entre agregados $\overline x_i,\overline y_i$, la asociación de variables
centradas $x_{i,a}-\overline x_i$ y $y_{i,a}-\overline y_i$, y la distribución de
correlaciones calculadas por posición son estimandos distintos. Sus signos pueden
divergir sin contradicción. En S4.3 deben conservarse por separado las reducciones
utilizadas en cada tabla; agruparlas bajo un único coeficiente ocultaría la jerarquía.

El remuestreo agrupado selecciona posiciones con reemplazo y conserva conjuntamente
sus enlaces, volviendo a calcular el estadístico en cada réplica. Este procedimiento
respeta la dependencia interna del grupo \citep{Efron1993Bootstrap}; si posiciones
vecinas siguen correlacionadas, no elimina la dependencia entre grupos. La extensión
espacial requiere bloques o sesiones de remuestreo y un análisis de sensibilidad a
su tamaño. Un intervalo estrecho no transforma una asociación observacional en
identificación causal.

\section{Métricas, agregación y validación espacial}

Para errores escalares $e_i=\widehat y_i-y_i$ se distinguen
\begin{equation}
\mathrm{MAE}=N^{-1}\sum_i|e_i|,\quad
\mathrm{MedAE}=\operatorname{mediana}_i|e_i|,\quad
\mathrm{RMSE}=\sqrt{N^{-1}\sum_i e_i^2}.
\end{equation}
El coeficiente $R^2=1-\sum_i e_i^2/\sum_i(y_i-\overline y)^2$ puede ser negativo
y no está definido para una respuesta constante. Pearson mide asociación lineal y
Spearman asociación por rangos; ninguno garantiza acuerdo sin sesgo. La media de
medianas entre particiones difiere, en general, de la mediana de las predicciones
agrupadas. Esta diferencia explica reducciones distintas entre S4.4f y S4.5b.

En validación espacial, todos los enlaces de una posición pertenecen a la misma
partición. Los conjuntos de ajuste y prueba se separan mediante una región de exclusión
de anchura declarada. La selección de hiperparámetros utiliza sólo entrenamiento y
validación, preferiblemente con una separación interior compatible con la evaluación
exterior. Los cinco bloques de S4.4f caracterizan regiones de una escena; no constituyen
cinco edificios independientes ni demuestran transferencia entre dispositivos.

La predicción de un retardo auxiliar debe evaluarse también en la respuesta compleja
que modifica. Esa evaluación debe registrar qué parte de la corrección se predice
y cuál se estima con una medición disponible. Esta separación fundamenta el alcance
condicionado de S4.5 y prepara la posterior validación del operador de observación.


\section{Predictibilidad circular y alcance de los controles negativos}
\label{sec:circular_predictability}

Para una fase $\psi\in\mathbb S^1$ y variables disponibles $\vect z$, la pérdida
circular $\ell(\psi,g)=1-\cos(\psi-g)$ tiene riesgo condicional
\begin{equation}
\mathbb E[\ell(\psi,g)\mid\vect z]
=1-\Re\{e^{-jg}m(\vect z)\},\qquad
m(\vect z)=\mathbb E[e^{j\psi}\mid\vect z].
\end{equation}
Si $m(\vect z)\ne0$, el predictor óptimo para esta pérdida es
$g^\star(\vect z)=\angle m(\vect z)$ y el riesgo mínimo es $1-|m(\vect z)|$.
Si $m=0$, ninguna dirección puntual mejora ese riesgo condicional. Esta propiedad
no identifica automáticamente el predictor óptimo bajo distancia angular absoluta,
cuya solución corresponde a una mediana circular. Tampoco la resultante marginal
$|\mathbb E e^{j\psi}|$ determina la resultante condicional: fases predecibles desde
$\vect z$ pueden cancelarse al agregarlas \citep{Mardia2000Directional,Bishop2006PRML}.

Bajo el control específico $\Delta\psi\sim\mathcal U(-\pi,\pi]$, se cumple
\begin{equation}
\mathbb E|\Delta\psi|=\operatorname{Med}|\Delta\psi|=\frac{\pi}{2},\qquad
\mathbb E\cos\Delta\psi=0,\qquad
P(|\Delta\psi|\leq\alpha)=\frac{\alpha}{\pi},\quad 0\leq\alpha\leq\pi.
\label{eq:uniform_phase_control}
\end{equation}
Así, 90 grados, $1/6$ para 30 grados y $1/4$ para 45 grados son referencias
analíticas útiles para S4.6. Su proximidad empírica no prueba que el error sea
uniforme o independiente. El fracaso del vecino próximo descarta una ventaja
práctica clara de esa regla en el soporte evaluado; no establece que
$m(\vect z)=0$ para cualquier conjunto de descriptores.

La concentración finita también requiere cautela. Para $A$ fases independientes
uniformes, $\mathbb E\overline R^2=1/A$, aunque la resultante poblacional sea
nula. Por tanto, una dirección común obtenida de $A=32$ antenas puede aparecer
por ajuste a fluctuaciones. La comparación pertinente repite el estimador de
$\beta$, el centrado por arreglo y las máscaras sobre controles que conserven
las marginales, o estima la dirección con unas antenas y evalúa las restantes.
Las antenas reales no se suponen independientes por esta identidad de referencia.

\section{Estructura temporal con muestreo irregular y trazabilidad}
\label{sec:temporal_circular_protocol}

Para marcas temporales físicas $t_i$, un intervalo $B$ define los pares válidos
$\mathcal I_B=\{(i,j):i<j,\ |t_j-t_i|\in B\}$ dentro de un mismo arreglo.
Se separan tres estadísticos:
\begin{align}
D_\psi(B)&=\operatorname{Med}_{(i,j)\in\mathcal I_B}
\left|\wrap(\psi_j-\psi_i)\right|,\\
C_\psi(B)&=\frac{1}{|\mathcal I_B|}\sum_{(i,j)\in\mathcal I_B}
\cos(\psi_j-\psi_i),\\
R_\psi(B)&=\left|\frac{1}{|\mathcal I_B|}\sum_{(i,j)\in\mathcal I_B}
e^{j(\psi_j-\psi_i)}\right|.
\label{eq:temporal_phase_structure}
\end{align}
Una rotación determinista no nula puede producir $C_\psi$ pequeño y $R_\psi$
elevado; ambos indicadores responden a preguntas diferentes. Los intervalos
sin pares válidos no se rellenan con cero. Cada curva debe especificar número
de pares, separación efectiva, regla de selección y posibles restricciones
espaciales; usar todos los pares o una muestra de ellos define ponderaciones
diferentes. En S4.6d.2 se conserva la tabla reportada, sin atribuirle conteos
por intervalo que el registro suministrado no especifica.

Los pares que comparten una captura son dependientes. Los intervalos y las
comparaciones entre vecinos temporales, espaciales y aleatorios requieren
remuestreo por capturas, bloques o sesiones y réplicas del control aleatorio.
Se debe informar incertidumbre de la diferencia pareada; proximidad de medianas
sin un margen prefijado no constituye equivalencia estadística. La elección de
bloques temporales debe considerar el desplazamiento espacial del transmisor:
una campaña móvil no identifica por sí sola un proceso de deriva a posición fija.

Recuperar una marca temporal mediante igualdad de coordenadas es válido para
el subconjunto auditado sólo si la clave es única y la transformación de
coordenadas está documentada. En campañas con posiciones repetidas se necesita
un identificador original de captura. Deben preservarse archivo, índice de
origen, marca temporal, partición y versión de calibración; reproducir únicamente
una semilla de permutación no reconstruye necesariamente la secuencia histórica.

\section{Aprendizaje condicionado por repetibilidad y sensibilidad física}
\label{sec:repeatability_sensitivity_learning}

La función de aprendizaje se elegirá después de caracterizar la adquisición. Para una
representación $\mathcal V$, se compararán pares con estado físico fijo y distinta
realización instrumental, y pares con una perturbación física conocida:
\begin{align}
D_{\rm rep}&=d\bigl(\mathcal V(Y(\mathbf q,\eta_1)),
\mathcal V(Y(\mathbf q,\eta_2))\bigr),\nonumber\\
D_{\rm fis}(\delta\mathbf q)&=d\bigl(\mathcal V(Y(\mathbf q+\delta\mathbf q,\eta)),
\mathcal V(Y(\mathbf q,\eta))\bigr),\label{eq:repeat_physical_distances}\\
Q(\delta\mathbf q)&=\frac{\operatorname{med}D_{\rm fis}(\delta\mathbf q)}
{\operatorname{med}D_{\rm rep}+\epsilon_d}.
\label{eq:physical_repeatability_ratio}
\end{align}
Las condiciones instrumentales idénticas de la segunda expresión son un ideal
experimental; su discrepancia residual se medirá con referencias y repeticiones.
$Q$ es descriptivo y depende de la escala y de $\epsilon_d$. Un codificador constante
colapsa ambas distancias; por ello no se optimizará únicamente este cociente.
Se exigirán respuesta física no degenerada, incertidumbre y capacidad de discriminar
perturbaciones frente a pares nulos. Se informarán curvas de detección a falsa alarma
fijada con validación, errores de desplazamiento y resultados por orientación.

Las comparaciones iniciales incluirán magnitud, canal complejo, productos relativos,
autocodificador complejo, reconstrucción enmascarada y aprendizaje contrastivo temporal.
La rama latente física añadirá supervisión de desplazamiento, respuesta diferencial
y consistencia entre modalidades cuando estas referencias existan. Las pérdidas de
reconstrucción y de tarea se evaluarán por separado: reducir una no implica reducir
la otra. Las regularizaciones de varianza y las pruebas de recuperación del estado
permitirán detectar colapso, sin atribuir significado causal a coordenadas latentes
por inspección visual.

WiFi-JEPA proporciona una comparación específica de aprendizaje autosupervisado de
CSI para postura 3D \citep{Kim2026WiFiJEPA}; su objetivo publicado no demuestra
preservación de fase respiratoria. Se incorporará una vez validadas las entradas,
con presupuesto de datos y cómputo comparable y sin incluir sesiones de prueba en el
preentrenamiento. GradNorm o PCGrad se utilizarán sólo si el diagnóstico de gradientes
muestra desequilibrio o conflicto relevante; su incorporación no sustituirá una
hipótesis física. Las restricciones de Maxwell se limitarán inicialmente a casos
canónicos verificables, antes de considerar una PINN de habitación completa.

\chapter{Observación de una misma escena mediante Wi-Fi y LoRa}
\label{ch:radio_obs}

El soporte geométrico y las representaciones del capítulo anterior son comunes a la escena,
pero la medición final depende de la frecuencia, la forma de onda, el receptor y el
protocolo. El presente capítulo separa ambos niveles para evitar atribuir a la propagación
una pérdida de información causada por la instrumentación o, a la inversa, interpretar
como equivalente información obtenida mediante tecnologías distintas.

\section{Principio de separación entre propagación y radio}

La \cref{eq:generative_observation} distingue tres objetos que no deben intercambiarse: el canal ideal $H_m$, su estimación $\widehat H_m$ y la observación almacenada $Y_m$. El simulador produce $\calP_m$ condicionado por la escena, la frecuencia y las antenas; el receptor determina qué transformación y muestreo de $H_m$ quedan disponibles. Por ello se escribe
\begin{equation}
Y_m=\mathcal O_m(H_m[\calP_m];\vect s_R,\vect s_N),
\end{equation}
con $m\in\{\mathrm{WiFi},\mathrm{LoRa}\}$. Esta separación permite comparaciones causales: se fijan $\Gamma(t)$ y el movimiento humano, y se modifican la frecuencia, el ancho de banda, la forma de onda y la etapa de entrada del receptor. No se exige $\calP_{\rm WiFi}=\calP_{\rm LoRa}$; se exige que ambas realizaciones provengan del mismo estado físico controlado.

\section{Modelo de observación Wi-Fi}

Para subportadoras $f_k$,
\begin{equation}
H_k(t)=\sum_{\ell}\alpha_\ell(t)e^{-j2\pi\nu_k\tau_\ell(t)}.
\end{equation}
La observación real se obtiene después de considerar la instrumentación y el protocolo:
\begin{equation}
\widetilde H_k(t)=g_k(t)H_k(t)e^{j\psi_k(t)}+n_k(t).
\end{equation}
Las variables $g_k$ y $\psi_k$ se modelan explícitamente durante la transferencia de la simulación a las mediciones reales. UWB-Fi demuestra que el muestreo disperso en frecuencia puede aumentar el ancho de banda efectivo, pero también que cada cambio de canal introduce desajustes aleatorios de fase que deben tratarse. Su enfoque combina la teoría de recuperación dispersa con un modelo aprendido guiado por trazado de rayos \citep{Li2024UWBFI}.

\section{Modelo de observación LoRa}

Sea $x_{\rm CSS}(t)$ la envolvente compleja de un chirrido LoRa. Se adopta un
coeficiente $\alpha_{{\rm L},\ell}(t)$ que ya contiene la fase de portadora
correspondiente a la longitud instantánea del camino. El canal multitrayecto de la \cref{eq:ltv_channel} produce
\begin{equation}
y_{\rm L}(t)=\sum_{\ell}\alpha_{{\rm L},\ell}(t)
 x_{\rm CSS}(t-\tau_{{\rm L},\ell}(t))+w_{\rm L}(t).
\label{eq:lora_channel}
\end{equation}
Con esta convención, el Doppler de propagación está contenido en la evolución
de fase de $\alpha_{{\rm L},\ell}(t)$; no se multiplica de nuevo por una
rotación independiente que represente el mismo movimiento. La aproximación
de banda estrecha permite ignorar dilataciones temporales cuando son
despreciables para el ancho de banda y velocidades considerados.

Después de retirar el chirrido,
\begin{equation}
d_{\rm L}(t)=y_{\rm L}(t)c^*(t),
\end{equation}
y posteriormente
\begin{equation}
D_{\rm L}[k]=\operatorname{FFT}\{d_{\rm L}(t)\}.
\end{equation}
La eliminación del chirrido convierte idealmente su pendiente en un tono cuya posición se detecta mediante la FFT; el CFO, el error temporal, la multitrayectoria y la interferencia alteran ese patrón \citep{Ghanaatian2019LoRaRx,Maleki2023CSS}. SenLoRa utiliza símbolos de alta confianza para incrementar la información de sensado disponible en tráfico ambiente y muestra que la disponibilidad temporal es tan importante como la sensibilidad física del canal \citep{SenLoRa2025}.

\section{Jerarquía de representaciones LoRa}

Se conservará la cadena de representaciones
\begin{equation}
IQ\rightarrow \text{dechirped IQ}\rightarrow
\text{intervalos de FFT}\rightarrow
(A,\phi)\rightarrow
(\mathrm{RSSI},\mathrm{SNR}).
\end{equation}
La eliminación del chirrido y una FFT compleja completa son transformaciones invertibles, y el par magnitud--fase conserva un número complejo salvo por la cuantización y las convenciones en magnitud nula. La pérdida aparece al seleccionar intervalos espectrales, descartar la fase, promediar, cuantizar o reducir todo el paquete a RSSI/SNR. La tesis documentará cada operador y evaluará qué estadístico es suficiente para la presencia, el movimiento macroscópico y la respiración.

La jerarquía no afirma que todo estimador basado en I/Q superará siempre a uno basado en RSSI. Afirma algo más preciso: una transformación determinista no invertible no puede crear información acerca del estado bajo un modelo probabilístico fijo. Un resumen puede, sin embargo, mejorar el desempeño con un número finito de muestras al descartar perturbaciones o regularizar el aprendizaje. Por ello se compararán tanto la información disponible como la robustez del estimador, siguiendo la distinción entre estadístico, estimador y cota de información \citep{Kay1993,CoverThomas2006}.

\section{LoRa como plataforma de alcance y sensado distribuido}

El estudio de sensado de largo alcance a través de paredes muestra la detección de desplazamientos pequeños mediante la razón entre las señales de dos antenas. En condiciones específicas, informa sensibilidad a movimientos de escala milimétrica y una dependencia clara entre el alcance y la magnitud del desplazamiento \citep{Zhang2020LoRa}. Más recientemente, el sensado de gran escala en edificios modela la señal como
\begin{equation}
s(t)=H_s+a(t)e^{-j2\pi f_cd(t)/c}
\end{equation}
y estudia la modificación de la cobertura efectiva de sensado mediante la configuración del transceptor \citep{Xue2025LargeScaleLoRa}. Esto sugiere que el problema de despliegue puede expresarse como una optimización de la observabilidad y no sólo de la cobertura de comunicación.

Los modelos de observación definidos hasta este punto describen cómo se registra un canal
dado. El siguiente paso consiste en introducir la evolución temporal de la escena y en
determinar cómo el movimiento macroscópico y las deformaciones fisiológicas modifican
las trayectorias que llegan a dichos observadores.

\section{Disponibilidad instrumental y conservación de información}

La recepción I/Q empleada en plataformas SDR debe distinguirse de los indicadores
por paquete disponibles en un dispositivo de comunicaciones. Los estudios LoRa con
dos antenas y el marco OpenLoRa justifican verificar sincronización y procesamiento,
pero no garantizan que un gateway convencional entregue el mismo observable
\citep{Zhang2020LoRa,Mishra2023OpenLoRa}. La validación especificará dispositivo,
frecuencia de muestreo, referencia entre antenas y continuidad entre paquetes.

Para $y_g=g\mu(q)+gw+v$, con ruido previo y posterior de covarianzas
$\Sigma_w$ y $\Sigma_v$, la información de la media contiene
\begin{equation}
2g^2\Re\{J_q^H(g^2\Sigma_w+\Sigma_v)^{-1}J_q\}.
\end{equation}
Si $v=0$ y el escalamiento es ideal e invertible, la ganancia se cancela.
La utilidad práctica de ajustarla depende de ruido añadido, cuantización,
rango dinámico o saturación. En el diseño de red no se mantendrá artificialmente
fija la covarianza mientras se amplifica señal y ruido conjuntamente.

\section{Referencia coherente y control de adquisición}
\label{sec:coherent_acquisition}

Para dos cadenas simultáneas con perturbación común $c(t,k)$,
$Y_s=c\,g_sH_s+n_s$ y $Y_r=c\,g_rH_r+n_r$. Un cociente regularizado útil es
\begin{equation}
\widetilde H_s=H_{r,\rm cal}\frac{Y_sY_r^{*}}{|Y_r|^2+\epsilon_r},
\label{eq:cabled_reference_ratio}
\end{equation}
donde $H_{r,\rm cal}$ representa la respuesta de referencia calibrada en la convención
adoptada, y la respuesta relativa $g_s/g_r$ deberá calibrarse. Sin ruido y para
$\epsilon_r\to0$, se cancela $c$ si es realmente común y $H_r\ne0$; con
regularización, el cociente conserva sesgo de amplitud. Una referencia por cable
requiere caracterizar pérdidas, fase y estabilidad de cables, divisores y cadenas.
Cambios de temperatura o movimiento del cable se incluyen en el presupuesto de error.

Compartir 10 MHz o PPS facilita sincronización de frecuencia o tiempo, pero no
garantiza por sí solo la fase relativa de osciladores locales después de un reinicio.
Una captura secuencial de referencia tampoco equivale a dos receptores coherentes
simultáneos. El ADALM-Pluto en su configuración documentada dispone de un transmisor
y un receptor \citep{AnalogDevicesPluto}; no se le atribuirá un segundo canal de
recepción simultánea sin una implementación adicional verificada. La arquitectura
se elegirá tras comprobar acceso a señal, relojes, disparos y repetibilidad.

FarSense utiliza el cociente entre antenas para sensado respiratorio
\citep{Zeng2019FarSense}. En este banco se ensayará el producto regularizado
$Y_aY_b^*/(|Y_b|^2+\epsilon_r)$ con máscaras fijadas previamente. Suprime una
rotación común, pero puede amplificar ruido cerca de desvanecimientos y eliminar
variación física compartida. HiSAC aporta una referencia para coherencia entre bandas
mediante una trayectoria de anclaje \citep{Pegoraro2024HiSAC}; la aplicabilidad del
anclaje a nuestra adquisición deberá comprobarse, no suponerse.

Por paquete Wi-Fi se conservarán CSI disponible, RSSI, ganancia/AGC, canal, ancho de
banda, configuración de antenas, MCS, identificador de secuencia y marcas temporales;
CFO y SFO se registrarán cuando el dispositivo los exponga. En LoRa se distinguirán
muestras I/Q, salida compleja completa de dechirp/FFT y telemetría del módem.
Una pasarela LoRaWAN no proporciona necesariamente I/Q ni CFR. El descarte de fase,
bins o muestras es una decisión de observación que debe documentarse; dechirp y FFT
completos no constituyen por sí solos una pérdida irreversible de información.

Para sincronización multimodal se estimará un mapa de reloj
$t_W=a_s t_s+b_s$ y su incertidumbre; saltos o deriva no afín obligarán a segmentar
la sesión. Índices de cuadro iguales no prueban simultaneidad. Antes y después de
cada sesión se repetirán controles de referencia y escena vacía, y se registrarán
reinicios, reconexiones, temperatura y versiones de firmware. RadioRange ofrece
modelos de perturbaciones para simulación de distancias UWB, Wi-Fi y 5G
\citep{Lyu2026RadioRange}; sus resultados de ranging no sustituyen la caracterización
de deriva de fase ni la referencia respiratoria de este experimento.

\chapter{Dinámica humana, fisiología y coherencia temporal}
\label{ch:human_dynamics}

La separación entre propagación y observación permite incorporar ahora la dinámica sin
confundirla con los efectos propios de la radio. Este capítulo distingue el movimiento
corporal macroscópico de las deformaciones fisiológicas y formula las condiciones de
coherencia temporal que un gemelo debe preservar para representar ambos fenómenos.

\section{Separación entre macroestado y microestado}

Se define
\begin{equation}
\vect s_H(t)=\left[\vect p_H(t),\vect r_H(t),\vect q_{\rm pose}(t),\vect v_H(t)\right]
\end{equation}
y
\begin{equation}
\vect s_F(t)=\left[x_{\rm chest}(t),\dot x_{\rm chest}(t),\ldots\right].
\end{equation}
El primero contiene movimiento de centímetros a metros; el segundo desplazamientos de milímetros. Son estados físicos; $Z_H$ y $Z_F$ se reservaron en la \cref{eq:factor_latent} para sus representaciones aprendidas. Mezclar ambas escalas desde el principio produce un problema numéricamente mal condicionado y metodológicamente ambiguo.

\section{Coherencia temporal de trayectorias}

Para una trayectoria móvil,
\begin{equation}
\phi_\ell(t)=-\frac{2\pi f_c}{c}L_\ell(t),
\qquad
\dot\phi_\ell(t)=2\pi\nu_\ell(t).
\end{equation}
Un simulador dinámico debe preservar la continuidad de fase y la correspondencia temporal entre trayectorias. Recalcular cada captura sin mantener la identidad de las trayectorias puede introducir saltos artificiales que destruyan precisamente la componente de micromovimiento que se pretende estudiar. WaveVerse se toma como referencia reciente de simulación dinámica con coherencia de fase \citep{Zheng2026WaveVerse}; en la implementación propia se estudiará una solución híbrida entre el nuevo trazado de trayectorias y su actualización local.

\section{Aproximación local diferenciable}

Para desplazamientos pequeños,
\begin{equation}
H(\vect q+\Delta\vect q)
\approx H(\vect q)+
\mat J_H(\vect q)\Delta\vect q+
\vect R_2(\Delta\vect q),
\label{eq:taylor_dynamic}
\end{equation}
donde $\vect R_2$ agrupa términos de segundo orden calculados componente a componente para el canal vectorial. Si se necesitara una expansión cuadrática explícita, cada componente compleja de $H$ tendría su propia Hessiana; no existe una única matriz Hessiana que produzca todo el tensor de canal.
En respiración, una primera aproximación es
\begin{equation}
\Delta H(t)\approx
\frac{\partial H}{\partial x_{\rm chest}}
 x_{\rm chest}(t).
\end{equation}
Esta aproximación se utilizará sólo cuando el conjunto de trayectorias sea estable. Para cambios grandes de postura, orientación o posición se realizará re-tracing.

La validez no se decidirá únicamente por la magnitud de $\Delta\vect q$. La aproximación también falla cerca de oclusiones, sombras, puntos de reflexión degenerados o eventos de aparición y desaparición de trayectorias. El reflector mecánico de la línea metodológica L3 permitirá medir el residuo de primer y segundo orden antes de aplicar la aproximación al tórax.

\section{Respiración como problema de recuperación de la forma de onda}

BreatheSmart estudia mediante un maniquí y clasificación supervisada cómo cambios pequeños en amplitud y fase de CSI contienen información de movimiento respiratorio y proporciona un entorno controlado para estudiar límites y factores de falla \citep{Mosleh2022BreatheSmart}. Wi-Actigraph amplía el contexto hacia sueño y perturbaciones respiratorias/motoras en un entorno realista \citep{Alzaabi2025WiActigraph}. En esta tesis, el objetivo primario no será clasificar BPM directamente, sino recuperar primero
\begin{equation}
Y(t)\rightarrow \widehat x_{\rm chest}(t)
\end{equation}
y después obtener métricas derivadas, por ejemplo $\mathrm{RR}=60f_{\rm resp}$ en respiraciones por minuto cuando una frecuencia instantánea sea una descripción válida. Esta elección conserva estructura temporal y permite evaluar pausas, irregularidad y variabilidad sin reducir toda la señal a un pico espectral.

\section{Movimiento macroscópico como fuente no lineal de interferencia}

MoVi-Fi, implementado con radares IR-UWB y FMCW, sostiene que la componente de signos vitales no desaparece necesariamente durante el movimiento, sino que se mezcla de manera no lineal con reflexiones de mayor energía, y emplea aprendizaje contrastivo para recuperar las formas de onda \citep{Chen2022MoViFi}. Su adaptación a CSI Wi-Fi constituye una hipótesis adicional. Este resultado motiva modelar
\begin{equation}
Y_m(t)=F_m(\vect s_E,\vect s_H(t),\vect s_F(t),
\vect s_R,\vect s_N(t))
\end{equation}
en lugar de asumir una suma aditiva independiente entre movimiento y fisiología.

\section{Varios usuarios y separabilidad física}

WiMANS contiene muestras con cero a cinco personas, CSI en dos bandas y video sincronizado, y muestra que el paso de uno a varios usuarios continúa siendo difícil \citep{Huang2024WiMANS}. WiMUSE formula la respiración de varias personas como un problema de \ac{BSS} de formas de onda y señala que los métodos puramente espectrales fallan cuando los sujetos presentan frecuencias cercanas o patrones no estacionarios \citep{Yi2025WiMUSE}. MURAL-Fi muestra, a su vez, que añadir diversidad espacial (AoA, ToF, AoD y Doppler) mejora la separación de regiones torácicas y la robustez frente a movimientos involuntarios \citep{Li2026MURALFI}. Estos trabajos respaldan una conclusión central: la complejidad del estimador no puede compensar de manera indefinida la ausencia de observabilidad física.

\section{Correspondencia entre variable física y referencia experimental}

La señal de una banda de expansión respiratoria no se identifica automáticamente
con desplazamiento local en milímetros. El estimando puede ser forma de onda
normalizada, desplazamiento calibrado o tasa; su referencia y métricas deben
corresponder. ViMo utiliza radio comercial de 60 GHz y una respuesta temporal
con recursos de separación distintos del CSI sub-7 GHz; sus resultados se
consideran un antecedente complementario, no un rendimiento directamente
transferible \citep{Wang2020ViMo}.

Se incluirán controles de sala vacía, actuador inmóvil y perturbaciones no
respiratorias. Una correlación elevada tras normalizar o deformar el eje temporal
no garantiza amplitud ni tiempo de eventos correctos. La evaluación informará
error temporal sin reajuste de prueba, tasa cuando esté definida, falsas detecciones
y riesgo frente a cobertura de las ventanas aceptadas.

\section{De la movilidad del terminal al movimiento de un dispersor}
\label{sec:motion_intervention}

El canal de dc01 cambia con la posición del transmisor en un entorno estático.
La derivada respecto de esa posición no equivale a la derivada respecto del tórax
u otro objeto con transmisor y receptor fijos. En el segundo caso pueden coexistir
trayectorias estáticas y dinámicas, cambios de visibilidad y variaciones del coeficiente
de dispersión. La intervención mecánica propuesta separará estos mecanismos antes
de interpretar señales humanas.

Para un estado de referencia $\mathbf q_0$ y perturbaciones pequeñas, se ensayará
\begin{equation}
\Delta H\simeq J_{\rm macro}\delta\mathbf q_{\rm macro}
+J_{\rm resp}\delta\mathbf q_{\rm resp}+r,
\label{eq:local_motion_decomposition}
\end{equation}
con $r$ que incluye términos no lineales y discrepancia. Es una expansión local;
no una separación exacta de contribuciones independientes en multitrayectoria.
La resta de canales requerirá una referencia común estable entre instantes.
Compensar libremente la fase de cada cuadro con el propio objetivo puede suprimir
la señal que se pretende medir. Por ello se evaluarán en paralelo el canal referido
y la respuesta de $\mathsf V_i$ al mismo desplazamiento.

Un fantoma mecánico permite variar amplitud, frecuencia, orientación y movimiento
lento con una referencia independiente. Una forma de onda periódica recuperada en
ese banco no valida respiración humana, y una pausa mecánica no valida detección
clínica de apnea. La transición a personas añade deformación no rígida, postura,
ropa, oclusión y movimiento espontáneo. Se analizarán primero esos factores y la
forma de onda de referencia antes de entrenar clasificadores de actividad o salud.

La conclusión anterior conduce de manera natural al problema de observabilidad. Antes de
seleccionar un estimador conviene determinar si la geometría, las tecnologías de radio y la
cadencia de medición contienen información suficiente para distinguir las variables de
interés. El capítulo siguiente formaliza esta pregunta mediante jacobianos e información
de Fisher.

\chapter{Observabilidad, información de Fisher e incertidumbre}
\label{ch:observability}

Los capítulos precedentes construyeron el modelo directo, desde la escena hasta la
medición. Este capítulo estudia su dirección inversa: bajo qué condiciones una variación
del estado humano produce cambios distinguibles en los observables. La respuesta se
formula localmente mediante sensibilidad, información estadística e incertidumbre, y
servirá después para evaluar tanto los algoritmos como el diseño de la red.

\section{Problema inverso, identificabilidad y observabilidad local}

El modelo directo lleva un estado $\vect q$ a una distribución de observaciones $p(\vect y\mid\vect q)$. La inversión pregunta qué estados son compatibles con una medición. La \emph{identificabilidad global} exige que estados distintos no induzcan la misma distribución; la \emph{identificabilidad local} estudia esa unicidad en una vecindad; y la \emph{estimabilidad práctica} incorpora ruido, número finito de muestras y error del modelo \citep{KaipioSomersalo2005,Kay1993}. En este manuscrito, \emph{observabilidad} denota información estadística local sobre variables físicas, no el criterio clásico de observabilidad de un sistema dinámico lineal.

Sea
\begin{equation}
\vect y=\vect\mu(\vect q)+\vect w,
\qquad
\vect w\sim\mathcal{CN}(\vect 0,\mat\Sigma_w).
\label{eq:local_observation}
\end{equation}
La matriz Jacobiana
\begin{equation}
\mat J_q=\frac{\partial\vect\mu}{\partial\vect q^T}
\label{eq:state_jacobian}
\end{equation}
describe sensibilidad, pero su norma no considera ruido ni correlación. Para parámetros reales, ruido complejo propio, covarianza conocida e independiente de $\vect q$, la \ac{FIM} es
\begin{equation}
\mat F_q(\vect q)
=2\Re\left\{
\mat J_q^H\mat\Sigma_w^{-1}\mat J_q
\right\}.
\label{eq:fim}
\end{equation}
Si la covarianza depende del estado aparecen términos adicionales; si el ruido no es gaussiano, la FIM se deriva de la función de puntuación y no de la \cref{eq:fim}. Bajo condiciones regulares, para parámetros identificables y estimadores localmente
insesgados, la cota de Cramér--Rao es
$\operatorname{Cov}(\widehat{\vect q})\succeq\mat F_q^{-1}$
\citep{Kay1993}. Si la información es singular, se declara un estimando local
$U^T\vect q$ en un subespacio identificable de base ortonormal $U$ y se aplica
la cota correspondiente a esa parametrización regular. La pseudoinversa no
asigna incertidumbre nula a una dirección invisible y no garantiza la existencia
de un estimador insesgado del vector completo. La cota tampoco equivale a un
intervalo posterior ni incorpora automáticamente discrepancia del gemelo.

\section{Perturbaciones e información equivalente}

Sea $\vect\zeta$ un vector de perturbaciones que contiene, por ejemplo, la fase común, la ganancia o la incertidumbre geométrica. La FIM conjunta se particiona como
\begin{equation}
\mat F(\vect q,\vect\zeta)=
\begin{bmatrix}
\mat F_{qq}&\mat F_{q\zeta}\\
\mat F_{\zeta q}&\mat F_{\zeta\zeta}
\end{bmatrix}.
\end{equation}
Cuando $\vect\zeta$ debe estimarse junto con $\vect q$, la información equivalente para el estado de interés es el complemento de Schur
\begin{equation}
\mat F_q^{\rm eq}=
\mat F_{qq}-
\mat F_{q\zeta}\mat F_{\zeta\zeta}^{\dagger}
\mat F_{\zeta q}.
\label{eq:equivalent_fim}
\end{equation}
Esta expresión formaliza por qué un canal sensible al movimiento puede ser poco informativo si el mismo cambio puede explicarse mediante CFO, fase desconocida o movimiento de otro sujeto.

\section{Escala, rango y criterios espectrales}

Los autovalores de una FIM dependen de las unidades de $\vect q$. Antes de comparar posición en metros, orientación en radianes y desplazamiento torácico en milímetros, se define
\begin{equation}
\bar{\vect q}=\mat D_q^{-1}(\vect q-\vect q_0),
\qquad
\mat F_{\bar q}=\mat D_q^T\mat F_q^{\rm eq}\mat D_q,
\label{eq:fim_normalization}
\end{equation}
con escalas físicas declaradas en $\mat D_q$. Para un conjunto fijo de variables
y $\mat F_{\bar q}\succ0$, se considerarán
\begin{align}
\lambda_{\min}(\mat F_{\bar q})
&\quad\text{(E-optimality)},\\
\log\det(\mat F_{\bar q})
&\quad\text{(D-optimality)},\\
\tr(\mat F_{\bar q}^{-1})
&\quad\text{(A-optimality)}.
\end{align}
Los criterios se comparan sobre el mismo subespacio objetivo; un diseño singular
se declara no admisible para recuperar todas sus variables. Añadir una matriz
de información previa define un criterio regularizado diferente y exige declarar
ese prior. En particular, minimizar la traza de una pseudoinversa sin controlar
el rango puede favorecer una configuración que pierde una variable.

El rango identifica el subespacio localmente visible; un valor propio pequeño indica una dirección mal condicionada, no necesariamente una variable individual invisible. Estos criterios provienen del diseño óptimo de experimentos y se reutilizarán para determinar la ubicación de los nodos \citep{Pukelsheim2006}.

\section{Respiración como ejemplo conductor}

Para $q=x_{\rm chest}$, y bajo las condiciones de la \cref{eq:fim}, el índice escalar
\begin{equation}
o_{\rm resp}(\vect x)=
2\Re\left\{
\left(\frac{\partial\vect\mu}{\partial x_{\rm chest}}\right)^H
\mat\Sigma_w^{-1}
\left(\frac{\partial\vect\mu}{\partial x_{\rm chest}}\right)
\right\}
\label{eq:resp_observability}
\end{equation}
define un mapa local de información del desplazamiento torácico con las
perturbaciones restantes conocidas. Si son desconocidas se utiliza la
información equivalente de la \cref{eq:equivalent_fim}. Su derivada hereda la geometría biestática de la \cref{eq:bistatic_sensitivity}; por ello el mapa puede diferir radicalmente de RSSI. La tesis validará si predice zonas favorables y puntos ciegos mejor que métricas de comunicación.

Presencia requiere un tratamiento diferente porque es una hipótesis discreta. Para $\mathcal H_0$ (ausencia) y $\mathcal H_1$ (presencia) se utilizarán razón de verosimilitudes, curvas ROC y separaciones distribucionales como KL o Chernoff, no una FIM ordinaria salvo que se introduzca y justifique una relajación continua.

\section{Redes distribuidas y dependencia estadística}

Si $N_{\rm link}$ enlaces son condicionalmente independientes dado $\vect q$, la información se suma:
\begin{equation}
\mat F_{\rm net}(\vect q)
=\sum_{l=1}^{N_{\rm link}}\mat F_{q,l}(\vect q).
\label{eq:fim_network}
\end{equation}
Con osciladores, interferencia o errores ambientales compartidos, la hipótesis falla y debe utilizarse la covarianza conjunta en la \cref{eq:fim}. Esta precaución distingue diversidad nominal de información realmente independiente.
Si los enlaces comparten una incógnita instrumental, se construye primero
la información conjunta y después se elimina esa incógnita; sumar
complementos de Schur independientes describe en general otro modelo.

\section{Observabilidad física y disponibilidad temporal}

En LoRa, la información puede existir físicamente pero no muestrearse con densidad suficiente. Para un estado aproximadamente constante y muestras condicionalmente independientes en $\mathcal T_{\rm obs}$,
\begin{equation}
\mat F_{\rm eff}=
\sum_{n\in\mathcal T_{\rm obs}}\mat F_{q,n}.
\end{equation}
Para respiración dinámica se empleará, en cambio, un modelo de estado y una FIM Bayesiana o posterior que incorpore correlación temporal. SenLoRa introduce métricas de disponibilidad en tráfico ambiente y estrategias para aumentar los símbolos útiles sin abandonar comunicación \citep{SenLoRa2025}. La tesis separará así \emph{información física por muestra}, \emph{cadencia disponible} e \emph{información acumulada}.

\section{Incertidumbre y abstención}

Se distinguirán el ruido aleatorio, la incertidumbre epistémica del estimador y la discrepancia del gemelo. La FIM caracteriza un límite local bajo un modelo supuesto; no cubre por sí sola el error de escena ni el cambio de dominio. El sistema informará una estimación, un intervalo calibrado y un índice de observabilidad, y podrá abstenerse cuando la evidencia sea insuficiente. La calibración se verificará mediante cobertura empírica y curvas de riesgo frente a cobertura, especialmente antes de interpretar una salida fisiológica.

\section{Información de la forma de onda y de la frecuencia respiratoria}

Para $x(t)=A\sin(2\pi f_rt+\varphi)$, la sensibilidad respecto de la tasa se
obtiene combinando la derivada del canal con
\begin{equation}
\frac{\partial x(t)}{\partial f_r}=2\pi At\cos(2\pi f_rt+\varphi).
\end{equation}
La información de $f_r$ depende de amplitud, duración, tiempos observados y
confusión con $A$, $\varphi$ y perturbaciones instrumentales. Si $A=0$, la
frecuencia no es identificable en este modelo, aunque $\partial H/\partial x$
sea grande. Por ello un mapa de sensibilidad del canal no basta para asegurar
una estimación de respiraciones por minuto.

En la validación externa, los índices predichos por el gemelo se contrastarán
con errores y fallos de estimadores sobre mediciones reservadas. Se distinguirá
el índice calculado con estado verdadero, utilizado como referencia explicativa,
del disponible con posición u orientación inciertas en inferencia.

\section{Efecto de eliminar perturbaciones sobre la sensibilidad útil}

Tras blanquear y apilar partes real e imaginaria de los jacobianos, sean
$A_q=\sqrt{2}[\Re(\Sigma_w^{-1/2}J_q)^T,
\Im(\Sigma_w^{-1/2}J_q)^T]^T$ y $A_\zeta$ su equivalente para las
perturbaciones. En el modelo local de covarianza fija,
\begin{equation}
F_q^{\rm eq}=A_q^T(I-A_\zeta A_\zeta^\dagger)A_q.
\end{equation}
La información útil corresponde a la sensibilidad que no puede explicarse
mediante esas perturbaciones. En particular, una fase instrumental arbitraria
por instante puede confundirse con el desplazamiento de un único camino.
La calibración basal o una referencia independiente permiten estudiar otra
condición de identificabilidad. Esta formulación fundamenta la prueba S6 de
preservación de movimiento después del procesamiento instrumental.


\section{Invariancia de fase y pérdida de sensibilidad física}
\label{sec:invariance_sensing_limit}

Para el observable espacial de la \cref{eq:spatial_gram_invariant}, una
perturbación pequeña satisface
\begin{equation}
\delta\mat G=(\delta\vect h)\vect h^H+\vect h(\delta\vect h)^H
+\mathcal O(\|\delta\vect h\|^2).
\end{equation}
Si $\delta\vect h=j\alpha\vect h$ con $\alpha\in\mathbb R$, los términos
lineales se cancelan. Para un cambio finito puramente rotacional la invariancia
es exacta. Por tanto, una representación que elimina una fase instrumental
común también puede eliminar una variación física que actúe en esa dirección.
Esto ocurre, por ejemplo, en la idealización de un único camino cuyo movimiento
sólo modifica su fase común y mantiene amplitud y vector espacial.

S4.7a contrastará fidelidad estática en el espacio invariante; S6 deberá evaluar
además $\partial\mat G/\partial x_H$ o la derivada del observable elegido,
comparándola con una referencia mecánica. La información de Fisher del observable
transformado debe derivarse de su distribución: los productos de canales
ruidosos no conservan automáticamente el modelo gaussiano aditivo ni su
covarianza original. Una buena concordancia de $G$ no garantiza que permanezca
información suficiente para recuperar respiración. Esta limitación conecta
la evaluación de invariantes con la proyección de perturbaciones y justifica
mantener una referencia de fase estable cuando la tarea física la requiera.

\section{Contraste experimental de la información preservada}
\label{sec:observable_jacobian_test}

La observabilidad se contrastará después de aplicar el operador realmente utilizado.
Para $\mathcal V$ no holomorfa, como magnitud o productos conjugados, se apilarán
partes reales e imaginarias del canal en $\mathbf y_{\mathbb R}$. Entonces
\begin{equation}
J_{\mathcal V,q}=\frac{\partial\mathcal V}{\partial\mathbf y_{\mathbb R}}
\frac{\partial\mathbf y_{\mathbb R}}{\partial\mathbf q},\qquad
\widehat J_{\mathcal V,q}^{(r)}=
\frac{\mathcal V(Y(\mathbf q+h\mathbf e_r))-
\mathcal V(Y(\mathbf q-h\mathbf e_r))}{2h}.
\label{eq:observable_real_jacobian}
\end{equation}
Se probarán varios pasos $h$ resolubles por la referencia mecánica para separar
error de truncamiento, histéresis y ruido amplificado por la diferencia finita.
La distancia cociente $\mathsf V_4$ exige analizar su geometría local o un representante
con referencia declarada; no se tratará como un vector de características arbitrario.

Se compararán dirección, norma y valores singulares del Jacobiano medido y simulado,
además de error relativo con un piso de denominador declarado. Si la sensibilidad
medida no se distingue del ruido, un error relativo pequeño no certifica un gemelo
útil. La FIM utilizará la covarianza del observable transformado y perturbaciones
instrumentales residuales; no se trasladará sin cambios el supuesto gaussiano del
CSI a productos o cocientes. Se confrontará su predicción con errores empíricos,
cobertura de intervalos y tasa de abstención en datos reservados.

Para dos personas, se estudiará el Jacobiano conjunto respecto de ambas fuentes y
sus direcciones casi colineales. Una tasa respiratoria correcta del componente
más fuerte no acredita separación de usuarios. Deben evaluarse asociación, errores
por persona y condiciones no identificables, antes de optimizar una red según esa
información. Los umbrales de confianza y de abstención se fijarán con validación.

La observabilidad establece qué puede recuperarse; no determina por sí sola cómo aprender
una representación robusta a partir de datos heterogéneos. El capítulo siguiente completa
el marco conceptual con aprendizaje autosupervisado, supervisión multimodal y adaptación
desde la simulación hacia las mediciones reales.

\chapter{Aprendizaje de representaciones y transferencia a mediciones reales}
\label{ch:representation_learning}

El aprendizaje se incorpora después de fijar el modelo físico y los criterios de
observabilidad. Este orden permite asignarle una función precisa: aproximar componentes
costosos, aprovechar datos sin etiquetar y adaptar la representación, sin exigir que una
red neuronal descubra por sí sola toda la estructura del problema.

\section{Motivación para aprendizaje autosupervisado}

Los conjuntos de datos reales de radiofrecuencia son costosos y abarcan una gran diversidad de geometrías. La simulación permite controlar factores y generar intervenciones, pero introduce discrepancias respecto de las mediciones. Una estrategia natural consiste en preentrenar representaciones mediante tareas que no requieran etiquetas humanas exhaustivas y, posteriormente, alinear la simulación, las mediciones reales de radiofrecuencia y las modalidades de referencia.

\section{Representación compartida y privada}

Para modalidad $m$ se propone
\begin{equation}
z_m=\left[z_{\rm comun},z_m^{\rm propio}\right].
\end{equation}
La parte compartida contiene variables físicas comunes (posición, postura, movimiento y respiración), mientras que la parte específica puede conservar peculiaridades de cada modalidad. No se impone una igualdad completa entre $z_{\rm RF}$ y $z_{\rm vision}$.

\section{Destilación multimodal}

Durante entrenamiento se dispone de una modalidad directa $V$ y una indirecta $Y$:
\begin{equation}
z_T=E_T(V),\qquad z_S=E_S(Y).
\end{equation}
Una pérdida de alineamiento es
\begin{equation}
\calL_{\rm align}=
\|P_S z_S^{\rm comun}-P_Tz_T^{\rm comun}\|_2^2.
\end{equation}
El trabajo de MERL, centrado en descripciones de actividades y no en
reconstrucción fisiológica, muestra que las modalidades directas e indirectas pueden combinarse durante el entrenamiento y que el conocimiento multimodal puede transferirse a un modelo que opera con información indirecta \citep{Hori2024Fusion}. En esta tesis, la cámara, el sensor de profundidad y el radar serán instrumentos de supervisión y validación, no necesariamente sensores del sistema desplegado.

\section{Objetivos predictivos y enmascaramiento estructurado}

Una arquitectura inspirada en I-JEPA \citep{Assran2023IJEPA} puede predecir representaciones ocultas en lugar de reconstruir todas las muestras. Las máscaras pueden operar sobre el tiempo, la frecuencia o el enlace:
\begin{align}
\mathcal M_T &: \text{ventanas temporales},\\
\mathcal M_F &: \text{subportadoras/canales},\\
\mathcal M_L &: \text{enlaces transmisor--receptor}.
\end{align}
El objetivo es inducir al codificador a capturar la estructura redundante del fenómeno y no únicamente artefactos locales.

\section{Transferencia de la simulación a las mediciones reales}

Se consideran tres dominios:
\begin{equation}
\mathcal D_{\rm sim},\qquad
\mathcal D_{\rm public},\qquad
\mathcal D_{\rm own}.
\end{equation}
El entrenamiento se organiza en tres fases: preentrenamiento físico a gran escala, adaptación con conjuntos de datos públicos y calibración con pocas muestras de adquisición propia. La eficiencia de datos se expresa mediante la curva de error frente a
mediciones reales y el número $N_{\rm objetivo}$ necesario para alcanzar
un umbral de tarea fijado en validación. Si se conserva $N_{90}$, se define
como el número que logra el 90\% de la mejora entre una referencia y una
saturación estimada, declarando ambas y el sentido de la métrica. Se cuentan
por separado posiciones, sesiones, participantes y datos auxiliares de adaptación.

\section{Separación de las perturbaciones instrumentales}

Se buscará reducir dependencia entre $Z_H$ y $Z_N$. Una regularización posible usa HSIC:
\begin{equation}
\calL_{\rm dep}=\operatorname{HSIC}(Z_H,Z_N),
\end{equation}
o un modelo adversario que intente predecir el dispositivo o el entorno a partir de la representación latente humana. Este término no garantiza independencia causal, pero proporciona un mecanismo empírico para evitar que el codificador utilice correlaciones espurias evidentes.

\section{Transferencia escalonada y comparación con modelos recientes}
\label{sec:transfer_updated_program}

La transferencia se distinguirá en cuatro niveles: región reservada de una escena;
nueva sesión o dispositivo; nueva distribución de objetos; y nueva habitación o banda.
Los cinco bloques de dc01 sólo respaldan el primero para los predictores evaluados.
Se reservarán sesiones completas incluso durante preentrenamiento autosupervisado;
ventanas solapadas y capturas contiguas no se repartirán entre entrenamiento y prueba.
La adaptación utilizará un presupuesto explícito de referencias y mediciones, incluyendo
el tiempo de registro geométrico y calibración, además del número de etiquetas.

En el cambio de objetos de S5, se comparará el material neuronal global B2 con una
variante local propuesta, $\mathbf x_{\rm loc}=T_{W\leftarrow o}^{-1}\mathbf x_W$.
Mover el objeto modificará su transformación, conservando identidad y parámetros
cuando el modelo lo permita. Esta variante no forma parte de B2 ya ejecutado.
Learnable WDT, RFScape y RadTwin se ubican en esta comparación editable; NeRF$^2$
se mantendrá como referencia de campo espacial aprendido. Se informarán error antes
de adaptar, curvas de adaptación y coste de actualización bajo el mismo acceso a
geometría y mediciones.

RFDT e InverTwin se incorporan como referencias de modelado diferenciable e inversión,
particularmente para el tratamiento de cambios de visibilidad y objetivos de ajuste
\citep{Chen2026RFDT,Chen2025InverTwin}. No se contarán sus contribuciones relacionadas
como validaciones independientes de esta tesis. Su evaluación inversa se condicionará
a haber contrastado sensibilidad y referencia instrumental; una buena reconstrucción
estática no garantiza un gradiente físico correcto. RayLoc conserva su función de
comparación inversa, mientras que WaveVerse orienta coherencia dinámica. Los resultados
externos de estos trabajos no validan automáticamente nuestro banco ni sus hipótesis.

La comparación aprendida incluirá el enfoque CSI de WiFi-JEPA, además de I-JEPA como
antecedente conceptual, autocodificación y métodos temporales. Las cámaras y el radar
podrán aportar supervisión durante entrenamiento; el experimento de despliegue declarará
si la inferencia usa sólo RF o requiere dichas modalidades. Se compararán escenarios
con y sin referencia de fase, para evitar atribuir al aprendizaje una mejora debida
exclusivamente a instrumentación adicional.

Con ello queda establecido el marco conceptual completo: modelo directo, representación,
operadores de observación, dinámica, observabilidad y aprendizaje. La parte siguiente
convierte estos elementos en preguntas de investigación, objetivos, hipótesis y un programa
experimental verificable.

\part{Programa doctoral y estrategia de validación}
\chapter{Planteamiento del problema, objetivos e hipótesis}
\label{ch:problem}

\section{Problema general}

El problema central se formula como sigue:
\begin{quote}
\emph{¿Qué representación física y latente del entorno, del canal y de la dinámica humana conserva la información necesaria para reconstruir observaciones inalámbricas y, simultáneamente, permite recuperar de forma robusta variables de presencia, movimiento y fisiología cuando cambian la geometría, la instrumentación, el protocolo y el dominio de despliegue?}
\end{quote}

La formulación evita restringirse a una sola tecnología. Wi-Fi y LoRa se estudian como operadores de observación diferentes de un mismo proceso físico. La cuestión prioritaria es qué variaciones pertenecen al entorno y al movimiento, cuáles dependen de la adquisición y qué evidencia permite distinguirlas. El modelo deberá conservar tanto estructura relativa estable como información coherente de la perturbación física.

\section{Preguntas específicas}

\begin{enumerate}[label=\textbf{PI\arabic*:},leftmargin=3em]
\item ¿Qué nivel de calibración física es necesario para que un gemelo reproduzca no sólo el canal complejo, sino también su sensibilidad a pequeñas perturbaciones geométricas?
\item ¿Qué representación del entorno ---trazado de rayos explícito, campo implícito, representación centrada en objetos o híbrida--- generaliza mejor con pocas mediciones?
\item ¿Puede separarse de manera útil la representación en $Z_E$, $Z_\phi$, $Z_R$, $Z_H$, $Z_F$ y $Z_N$, y distinguir dentro de la fase los componentes predecibles de aquellos que requieren calibración o invariancia?
\item ¿Qué información de sensado se pierde al pasar de muestras I/Q complejas a representaciones más comprimidas, como amplitud, fase, características demoduladas o RSSI?
\item ¿Cómo cambia la observabilidad de una misma perturbación humana entre Wi-Fi y LoRa?
\item ¿Cuándo un fallo de estimación es atribuible al algoritmo y cuándo a falta de observabilidad física?
\item ¿Puede la simulación diferenciable reducir de forma medible la cantidad de datos reales necesarios para transferir el sistema a una nueva escena?
\item ¿Cómo deben seleccionarse posiciones, enlaces, ganancias y configuraciones de radio para maximizar observabilidad de una región objetivo bajo restricciones de comunicación y energía?
\end{enumerate}

Los objetivos y las hipótesis que siguen traducen las preguntas anteriores en resultados
observables. Su formulación conserva la distinción entre lo que se propone construir y lo
que deberá demostrarse mediante los experimentos.

\section{Objetivo general}

Desarrollar y validar un marco físico--latente de gemelo digital inalámbrico para el sensado humano mediante infraestructura de comunicación existente, capaz de modelar el canal de forma diferenciable, separar los factores del entorno, la dinámica humana, la fisiología y la instrumentación, cuantificar la observabilidad y la incertidumbre, y transferirse progresivamente desde la simulación y los datos públicos hacia escenarios propios en interiores y exteriores con Wi-Fi y LoRa.

\section{Objetivos específicos}

\begin{enumerate}[label=\textbf{O\arabic*:},leftmargin=3em]
\item Construir un banco de pruebas de canal estático, con datos reales y simulados de DICHASUS dc01, y comparar el trazado de rayos nominal con la calibración de materiales y la calibración sensible a la fase, caracterizando por separado retardo efectivo, dirección común y estructura relativa.
\item Evaluar representaciones explícitas e implícitas del entorno en términos de error complejo, eficiencia de datos, generalización espacial y editabilidad, distinguiendo error bruto y error respecto de invariancias de fase explícitas.
\item Construir y registrar una geometría métrica con incertidumbre explícita, distinguir las escalas global, de antenas y de desplazamiento local, y evaluar G0--G4 junto con una referencia coherente y controles de repetibilidad.
\item Construir una representación común de escena con trayectorias específicas de cada radio, capaz de sintetizar observaciones Wi-Fi OFDM y LoRa CSS/IQ bajo modelos controlados de instrumentación e interferencia.
\item Introducir la dinámica de manera acumulativa: reflector mecánico, primitivas humanas, movimiento macroscópico, respiración, combinación de movimiento y respiración, y varios usuarios.
\item Definir y validar criterios de detección para presencia y métricas de observabilidad basadas en jacobianos e información de Fisher para movimiento y respiración.
\item Aprender representaciones factorizadas y multimodales que reduzcan la dependencia respecto del entorno y la instrumentación, y que preserven las variables humanas de interés.
\item Resolver el problema inverso con estimación de incertidumbre y capacidad explícita de abstención cuando el estado sea pobremente observable.
\item Optimizar redes distribuidas Wi-Fi/LoRa para maximizar la información útil de sensado dentro de áreas de interés, bajo restricciones de potencia, tráfico, número de nodos y cobertura de comunicación.
\end{enumerate}


\section{Hipótesis central y contrastes falsables}
\label{sec:updated_hypotheses}

Se plantea que una representación factorizada, robusta frente a transformaciones
instrumentales identificadas y sensible a cambios físicos del canal, puede mejorar
la transferencia entre días, dispositivos y escenarios sin perder la información
diferencial necesaria para el sensado humano. La referencia coherente y la geometría
métrica constituyen variables experimentales explícitas. La hipótesis no afirma
reconstrucción perfecta de fase absoluta a partir de LiDAR, ni identificabilidad
universal de la factorización latente.

Las siguientes seis hipótesis operacionalizan esta formulación. Se contrastarán con
intervalos pareados y un tamaño mínimo de efecto fijado antes de la evaluación
confirmatoria; resultados compatibles con ausencia de mejora se informarán como tales.
\begin{description}
\item[H1: robustez relativa.] Productos espaciales/espectrales y distancias respecto de
fase común reducirán discrepancia entre sesiones y dispositivos frente al canal bruto,
bajo perturbaciones compatibles con su grupo de invariancia. Se contrastará en S4.7a,
Esc. E0 y Esc. E5, conservando error bruto y sensibilidad física como controles.
\item[H2: sensibilidad coherente.] Una rama con referencia independiente preservará mejor
la respuesta a micromovimientos que magnitud o Gram espacial en configuraciones donde
la perturbación se exprese predominantemente en fase. Esc. E1--E2 contrastarán Jacobianos,
desplazamiento y forma de onda; se reportarán también geometrías de sensibilidad nula.
\item[H3: eficiencia física.] Geometría registrada y descriptores RT reducirán las
mediciones de adaptación necesarias respecto de métodos puramente basados en datos.
S5 y Esc. E5 compararán curvas error--datos, costes y G0--G4. $N_{90}$ se calculará
respecto de una referencia de rendimiento y tarea fijadas previamente; no se mezclará
entre tareas ni se escogerá retrospectivamente el mejor techo de cada método.
\item[H4: referencia instrumental.] Un canal de referencia calibrado reducirá la
variabilidad instrumental más que una corrección exclusivamente aprendida bajo reinicio,
deriva y cambio de sesión. Esc. E0 medirá error residual y repetibilidad; Esc. E1
comprobará que la reducción conserva respuesta al movimiento. Se incluirá el coste
de instrumentación adicional y el fracaso cuando la perturbación no sea común.
\item[H5: complementariedad.] La combinación robusta--coherente mejorará transferencia
frente a CSI bruto, magnitud sola y cada rama aislada con igual supervisión y presupuesto.
Esc. E2--E5 medirán error de tarea, incertidumbre y abstención; la fusión no se
considerará superior si su beneficio desaparece al controlar referencia y datos.
\item[H6: observabilidad orientada a la tarea.] Jacobianos y Fisher predecirán mejor las
condiciones de recuperación de una o dos personas que RSSI/SNR por sí solos. Esc. E1--E4
permitirán contrastar esa capacidad predictiva antes de utilizarla para seleccionar enlaces
y optimizar una red. Se evaluarán falsos diagnósticos de observabilidad y cobertura de
incertidumbre, además del error medio.
\end{description}

La interfaz multirradio, la inversión y el diseño de red se mantienen como objetivos;
su desarrollo depende de los contrastes anteriores. Un resultado favorable de H1 no
prueba H2, y una mejora de potencia en G4 no verifica H3 para canal complejo o fisiología.

\chapter{Metodología general}
\label{ch:methodology}

\section{Principio metodológico: incorporación gradual de complejidad}

La metodología incorpora una sola fuente principal de complejidad en cada fase. Los elementos validados se conservan, se documentan sus supuestos y se añade el componente siguiente. El avance no depende de alcanzar de manera aislada un valor de exactitud, sino de demostrar que el nuevo componente explica un fenómeno que el modelo anterior no podía representar. Esta lógica enlaza el modelo físico con la adquisición experimental y, posteriormente, con el problema inverso.

\begin{table}[ht]
\centering
\caption{Líneas metodológicas de la investigación doctoral.}
\label{tab:lineas_metodologicas}
\begin{tabularx}{\textwidth}{p{1.3cm}L L}
\toprule
\textbf{Línea} & \textbf{Pregunta científica} & \textbf{Resultado principal}\\
\midrule
L1 & ¿Puede el gemelo reproducir el canal real y conservar su fase? & Gemelo estático nominal y calibrado.\\
L2 & ¿Cómo debe representarse el entorno? & Comparación de $Z_E$ y análisis de eficiencia de datos.\\
L3 & ¿Cómo preservar la dinámica y las variaciones pequeñas de fase? & Gemelo dinámico para movimiento y respiración.\\
L4 & ¿Cómo transforma cada tecnología de radio la información física? & Modelos de observación para Wi-Fi y LoRa.\\
L5 & ¿Qué variables pueden recuperarse y con qué incertidumbre? & Solución del problema inverso y análisis de observabilidad.\\
L6 & ¿Cómo configurar una red para maximizar la información útil? & Diseño optimizado de la red distribuida.\\
\bottomrule
\end{tabularx}
\end{table}

\section{Cadena de evidencia y ejecución concurrente}
\label{sec:updated_method_chain}

La secuencia científica actualizada es geometría métrica, calibración macroscópica,
factorización de fase, representación robusta/coherente, fidelidad diferencial y,
finalmente, sensado e inversión. Cada enlace requiere evidencia propia. El resultado
de DICHASUS motiva separar retardo transferible, referencia de fase y residual;
no autoriza a atribuir toda discrepancia a un defecto geométrico o instrumental.

Se desarrollarán dos frentes concurrentes. El computacional completará S4.7a con
$\mathsf V_0$--$\mathsf V_7$, controles de invariancia y particiones espaciales;
S4.7b se justificará por el residual y S5 comparará representaciones. El experimental
construirá referencia y metrología, caracterizará Esc. E0 y contrastará Esc. E1--E2
antes de adquirir datos humanos. La preparación de S6 no esperará a implementar todas
las arquitecturas de S5. Los operadores Wi-Fi y LoRa se instrumentarán desde el inicio,
aunque su comparación de sensado exija una perturbación común ya validada.

\subsection{Pruebas de esfuerzo de la representación}
\label{sec:representation_stress_tests}

Antes de entrenar un predictor, cada $\mathsf V_i$ se someterá a transformaciones
sintéticas sobre el mismo canal. Para $H_{a,k,t}$ se definen:
\begin{align}
\mathsf T_1(H)&=e^{j\phi_t}H,
&\phi_t&\sim\mathcal U(-\pi,\pi),\nonumber\\
\mathsf T_2(H)_{a,k,t}&=e^{j\phi_{g(a),t}}H_{a,k,t},
&&\text{fase por arreglo},\nonumber\\
\mathsf T_3(H)_{a,k,t}&=e^{-j2\pi\nu_k\delta\tau_t}H_{a,k,t},
&&\text{retardo de referencia},\nonumber\\
\mathsf T_4(H)&=g_tH,
&&\text{ganancia},\label{eq:s47_stress_transforms}\\
\mathsf T_5(H)_{a,k,t}&=e^{j\theta_a}H_{a,k,t},
&&\text{desfase por cadena},\nonumber\\
\mathsf T_6(H)_{a,k,t}&=e^{j2\pi\Delta f\,t}H_{a,k,t},
&&\text{CFO temporal},\nonumber\\
\mathsf T_7(H)&=\mathcal Q\!\left(\mathcal I_{\rm IQ}(H)+n\right),
&&\text{IQ, cuantización y ruido}.\nonumber
\end{align}
$g_t$, $\delta\tau_t$, $\Delta f$ y el nivel de ruido se obtendrán de rangos
observados o de una condición experimental declarada. RadioRange orientará la
inyección controlada de perturbaciones, sin importar automáticamente su modelo de
ranging a la tarea de sensado \citep{Lyu2026RadioRange}.

Las identidades algebraicas se verificarán con tolerancia numérica relativa fijada
antes de la ejecución; los efectos de ruido, regularización y máscaras se resumirán
mediante distribuciones e intervalos agrupados. La matriz final indicará, para cada
par $(\mathsf V_i,\mathsf T_j)$, error relativo de invariancia y pérdida de sensibilidad.
No se asignará de antemano una etiqueta binaria de ``robusto'': por ejemplo, una matriz
de Gram por frecuencia rechaza fase común y retardo común, pero cambia ante fases de
cadena y puede ocultar movimiento compartido. La selección seguirá un frente de Pareto
entre rechazo de perturbaciones y retención de respuesta física.

\begin{table}[htbp]
\centering\small
\caption{Relación entre escenarios físicos nuevos y nomenclatura previa. Todos los escenarios E0--E5 son propuestos.}
\label{tab:scenario_program_map}
\begin{tabularx}{\textwidth}{P{1.6cm}L P{3.4cm}}
\toprule
Escenario & Pregunta y transición & Correspondencia\\
\midrule
Esc. E0 & Estabilidad con escena y terminales fijos; referencia instrumental. & L1/L4; soporte de S6/S7.\\
Esc. E1 & Desplazamiento conocido de reflector; respuesta diferencial. & E3.1; S6; experimento D.\\
Esc. E2 & Fantoma periódico y perturbaciones combinadas. & L3; S6/S7; D y F.\\
Esc. E3 & Una persona: movimiento, postura y respiración diferenciados. & E3.2--E3.5; F y G.\\
Esc. E4 & Dos personas: separabilidad antes de reconocimiento. & E3.6; experimento H.\\
Esc. E5 & Día, dispositivo, objetos, habitación y banda reservados. & S5/S7; C y G; soporte de I.\\
\bottomrule
\end{tabularx}
\end{table}

Los identificadores Esc. E0--E5 nombran condiciones físicas y no reemplazan los
subexperimentos E1.1--E3.6 ni los experimentos temáticos A--I. Las condiciones G0--G4
comparan geometría; $\mathsf V_i$ compara observables; L1--L6 nombra líneas científicas
y LT1--LT6, tareas administrativas. La \cref{sec:physical_scenarios} fija controles,
referencias y criterios de paso para evitar que estas dimensiones se confundan.

\section{Primera etapa: modelado de canal en escenarios controlados}
\label{sec:first_method_stage}

La primera etapa establece la relación entre escenario, representación electromagnética
y respuesta compleja antes de introducir una tarea humana. Se organiza en cinco niveles:
validación analítica; reconstrucción del entorno dc01; descomposición y estadística de fase;
predicción y transferencia espacial; y caracterización de los componentes todavía no
predichos. El \cref{ch:marco_canal} fundamenta propagación, interferencia y referencia
espectral; el \cref{ch:learning_foundations} establece los estimadores y la evaluación.

\begin{table}[htbp]
\centering\small
\caption{Escenarios de la primera etapa y alcance de su evidencia.}
\begin{tabularx}{\textwidth}{P{2.8cm}L L}
\toprule
Escenario & Propiedad que se contrasta & Estado y relación experimental\\
\midrule
Espacio libre y una reflexión & Escala, fase, retardo y Fresnel & Verificación analítica del simulador; protocolo de E1.1.\\
Multitrayectoria controlada & Interferencia y sensibilidad a fase por camino & S4.1: reproducción PEOC/UPEC/PEAC y estabilidad de estimación.\\
Escena nominal dc01 & Ganancia, magnitud y dispersión temporal & S1/B0--B2: ejecutados con partición original.\\
Frecuencias y antenas reservadas & Transferencia de fase afín y estructura residual & S4.2--S4.3: ejecutados; distintos accesos a CPO.\\
Regiones reservadas de dc01 & Predicción de retardo y utilidad en el canal & S4.4--S4.5: ejecutados con cinco bloques y exclusión de 0.5 m.\\
Capturas con tiempo recuperado & Concentración y continuidad de la fase común & S4.6a--d.2: ejecutados sobre B0; sin ventaja clara de los vecinos ensayados.\\
Otra configuración de objetos o escena & Editabilidad y transferencia del entorno & S5: propuesto; no demostrado por los bloques de dc01.\\
\bottomrule
\end{tabularx}
\end{table}

Esta organización distingue variación de condiciones dentro de una escena de transferencia
entre escenas. La calibración de materiales B0--B2 y el predictor de retardo se evalúan
con protocolos diferentes; el resultado espacial del segundo no acredita automáticamente
el del primero. S4.6 caracteriza el desfase residual y documenta controles
espaciales y temporales. S4.7a evaluará fidelidad relativa mediante invariantes;
S4.7b contrastará incertidumbre por trayectoria. Los criterios concretos y sus condiciones de acceso a datos se presentan en
la \cref{sec:future_route}.

Las comparaciones externas tienen una función específica: Hoydis fundamenta la reproducción
material; Ruah, la incertidumbre de fase; NeRF$^2$, la predicción espacial de campos;
Learnable WDT y RadTwin, la representación del entorno editable; WaveVerse, la coherencia
dinámica; y RayLoc, la inversión diferenciable \citep{Hoydis2024DiffRT,Ruah2024Phase,
Zhao2023NeRF2,Jiang2025LearnableWDT,Zhang2026RadTwin,Zheng2026WaveVerse,Han2026RayLoc}.
Su incorporación depende de una consulta y un observable comparables, no de una lista
indiferenciada de arquitecturas.

\subsection{Decisión metodológica después de S4.6}

La decisión siguiente parte de un resultado positivo: el retardo predicho ya
mejora la concordancia angular y la coherencia, tanto en interpolación densa
como en regiones reservadas. S4.7a aborda la fase común que todavía se estima
con la medición, y S4.7b la discrepancia relativa restante. Así, cada bloque
responde a una limitación específica de una corrección que ha demostrado
utilidad parcial, conforme al balance de la \cref{sec:synthesis}.


El diagnóstico comprende 10\,000 capturas, 640\,000 fases por enlace y un
manifiesto temporal cuya correspondencia de coordenadas es exacta en ese
subconjunto. La fase común se estima con las propias antenas de cada captura,
por lo que su eliminación es descriptiva. Las reglas de vecindad no acreditan
predictibilidad local; tampoco excluyen estructura condicionada a variables
no examinadas. La distinción entre fase efectiva residual y CPO físico evita
atribuir al entorno o al receptor un error cuya causa no ha sido identificada.

S4.7a comprobará productos espaciales y entre frecuencias, junto con una
métrica del canal definida respecto de una fase común por arreglo. Los productos
espaciales eliminan también retardo común, por lo que la contribución del
retardo predicho se medirá con productos entre frecuencias o con el canal
apilado. El protocolo completo se fija en la \cref{sec:s47a_protocol}.
S4.7b estudiará el residual restante y S5 comparará representaciones bajo el
mismo acceso a mediciones y el mismo grupo de invariancia. La validación
diferencial S6 determinará si esa invariancia conserva la señal de movimiento.

\section{Línea 1: gemelo físico estático}

El estado de avance de esta línea se presenta en el \cref{ch:avances}. S1/B0--B2 y S4.1--S4.6d.2 cuentan con ejecuciones documentadas. El avance incluye análisis residual, predicción de retardo, evaluación por bloques espaciales y estudio de la fase común; S4.6 completa el diagnóstico circular y temporal; S4.7a--b permanecen propuestos. Las actividades que siguen fijan el protocolo completo y sus verificaciones pendientes; su enumeración no presupone que todas se hayan ejecutado.

\subsection{E1.1: validación analítica mínima}

Antes de interpretar resultados de escenas complejas se validarán:
\begin{enumerate}
\item espacio libre: Friis, fase y delay;
\item una reflexión: coeficientes de Fresnel y dos rayos;
\item difracción en configuración controlada;
\item movimiento de reflector de $0.5$ a $50$ mm;
\item consistencia entre derivada numérica y derivada diferenciable.
\end{enumerate}
El criterio de validación exige que la salida compleja del simulador reproduzca las soluciones analíticas dentro de tolerancias numéricas definidas antes del experimento.

\subsection{E1.2: configuración nominal de DICHASUS dc01}

Se utilizaron posiciones reales y la escena INUE para construir un banco reportado de $10000$ posiciones con
\begin{equation}
\left(\vect p_i,H_i^{\rm real},\calP_i,H_i^{\rm RT}\right).
\end{equation}
La referencia B0 utiliza materiales nominales y una escala global congelada. La comparación adicional sin ajuste de escala constituye una ablación propuesta; no debe confundirse con B0. Se conservan los caminos trazados, su activación, el orden de frecuencias y la partición original 5000/100/4900.

\subsection{E1.3: particiones espaciales}

La partición original se conserva como referencia de interpolación densa. S4.4d documenta su proximidad espacial, y S4.4e--f añade cinco bloques con exclusión de 0.5 m. Para cada protocolo se distinguen:
\begin{align}
\Omega_{\rm ent}&,\\
\Omega_{\rm val}&,\\
\Omega_{\rm prueba}&,
\end{align}
con regiones separadas. Se estudiarán tres regímenes: interpolación, extrapolación dentro de una escena y transferencia entre escenas cuando existan conjuntos de datos compatibles.

\subsection{E1.4: calibración acumulativa}

Comparaciones:
\begin{enumerate}
\item RT nominal con materiales de catálogo;
\item RT + escala global;
\item materiales aprendidos;
\item materiales neuronales dependientes de la posición;
\item materiales + scattering/antena cuando sea identificable;
\item diagnóstico de CPO y retardo efectivo, con referencias oráculo y transferencia;
\item calibración sensible a la fase y modelado del error local de fase.
\end{enumerate}
Cada versión utilizará la misma partición de datos. B0, B1 y B2 designan respectivamente materiales nominales, globales aprendidos y neuronales espaciales; las identificaciones no deben confundirse con los identificadores de experimentos A--I. Los resultados y las particiones documentadas se detallan en el \cref{ch:avances}.

\subsection{Métricas}

\paragraph{NMSE complejo}
\begin{equation}
\operatorname{NMSE}_H=
\frac{\sum_i\|\widehat H_i-H_i\|_2^2}
{\sum_i\|H_i\|_2^2}.
\end{equation}

\paragraph{Error de magnitud}
\begin{equation}
E_A=\frac{1}{N}\sum_i
\big\||\widehat H_i|-|H_i|\big\|_1.
\end{equation}

\paragraph{Pérdida circular de fase}
\begin{equation}
\calL_{\phi,\rm circ}=\frac{1}{NK}
\sum_{i,k}\left[1-
\cos(\angle\widehat H_{ik}-\angle H_{ik})\right].
\end{equation}
Esta pérdida es adimensional. Para reportar errores angulares en grados se utiliza la distancia circular absoluta $E_{\phi,\rm deg}$ de la \cref{eq:error_angular_avances}; ambas métricas deben distinguirse. Se añade la coherencia compleja normalizada de la \cref{eq:coherencia_avances}, con declaración de los ejes de agregación. Los valores brutos, alineados sobre la propia observación y transferidos se informarán por separado.

\paragraph{Dispersión de retardos.}
A partir de $h=\operatorname{IFFT}(H)$ se compararán el perfil de potencia en función del retardo y su dispersión cuadrática media.

\paragraph{Eficiencia de datos}
Se repetirá calibración con $N\in\{25,50,100,250,500,1000,\ldots\}$ para obtener $E(N)$.

\section{Línea 2: comparación de representaciones del entorno}

\subsection{Familias}

La comparación inicial realizada comprende RT nominal B0, materiales globales B1 y materiales neuronales espaciales B2. S4.4--S4.5 comparan predictores auxiliares del retardo sobre B0; no constituyen todavía el banco de campos neuronales S5. Se mantiene como ampliación propuesta la comparación con un campo neuronal implícito tipo NeRF$^2$, un gemelo centrado en objetos, una representación informada por la geometría tipo RadTwin y una formulación híbrida propia basada en descriptores de trayectorias y objetos \citep{Zhao2023NeRF2,Jiang2025LearnableWDT,Zhang2026RadTwin}. B2 no se identifica con esas arquitecturas ni acredita por sí solo transferencia entre escenas.

\subsection{Comparación por capacidades y acceso a información}

\begin{table}[htbp]
\centering\small
\caption{Función de las familias de gemelo en la comparación propuesta S5.}
\label{tab:twin_capability_comparison}
\begin{tabularx}{\textwidth}{P{2.8cm}L L}
\toprule
Familia & Capacidad que se contrastará & Condición de comparación\\
\midrule
RT B0/B1/B2 & Propagación explícita y calibración material global/espacial. & Geometría, mecanismos y datos de ajuste declarados.\\
NeRF$^2$ & Predicción espacial mediante campo aprendido de mediciones. & Densidad de muestras, consulta y regiones reservadas iguales.\\
Learnable WDT & Transporte y representación del entorno centrada en objetos. & Acceso a objetos y coste de adaptación explícitos.\\
RFScape & Representación neuronal de objetos integrada en trazado RF. & Objetos trasladados con identidad conservada y prueba independiente.\\
RadTwin & Uso de información geométrica en escenas variables. & Igual disponibilidad de geometría y observables de evaluación.\\
Híbrido propuesto & Trayectorias, coordenadas locales y dos ramas de fase. & Ablaciones de geometría, referencia, latente y supervisión.\\
\bottomrule
\end{tabularx}
\end{table}

Estas familias se apoyan en los trabajos de referencia \citep{Zhao2023NeRF2,Jiang2025LearnableWDT,
Chen2025RFScape,Zhang2026RadTwin}. La tabla define preguntas experimentales, no
un orden de superioridad publicado. Cuando un método requiera información adicional,
se comparará dentro del régimen correspondiente y se cuantificará el beneficio de
esa información. Sionna RT constituye el antecedente canónico del simulador diferenciable
\citep{Hoydis2023SionnaRT}; sus desarrollos posteriores no cambian retrospectivamente
la versión con la que se obtuvieron los resultados documentados.

\subsection{Tareas}

\begin{enumerate}
\item reconstrucción de canal en posiciones conocidas;
\item interpolación espacial;
\item extrapolación espacial;
\item adaptación con pocas mediciones y coste de calibración declarado;
\item cambio de objetos/materiales;
\item transferencia de escena;
\item inversión de parámetros físicos.
\end{enumerate}

\subsection{Métricas adicionales}

Además del error de canal se reportarán memoria, tiempo de inferencia, número de mediciones de calibración, error bajo edición y calidad del Jacobiano respecto de cambios físicos.

\section{Línea 3: gemelo dinámico}

\subsection{E3.1: reflector mecánico}

Se utilizará una superficie de geometría conocida y desplazamiento medido con
transmisor y receptor fijos. La rejilla inicial será
\begin{equation}
\delta x\in\{0,0.25,0.5,1,2,5,10\}\;\mathrm{mm},
\end{equation}
restringida por la incertidumbre efectiva de la etapa. Se compararán $H(t)$,
$\Delta H$, fase, Doppler y $\partial\mathcal V/\partial x$; pasos mayores podrán
introducirse como extensión macroscópica. Esc. E1 de la \cref{sec:scenario_e1}
fija repeticiones, referencias y controles de histéresis.

\subsection{E3.2: primitivas humanas}

Progresión:
\begin{equation}
\text{placa}\rightarrow
\text{elipsoide}\rightarrow
\text{cápsula}\rightarrow
\text{cuerpo articulado}\rightarrow
\text{SMPL}.
\end{equation}
La finalidad es cuantificar el grado de detalle geométrico necesario para preservar las derivadas relevantes para el sensado.

\subsection{E3.3: movimiento macroscópico}

Se simularán acciones de marcha, giro, sedestación, incorporación y movimiento de brazos. WiMANS y los conjuntos de datos con esqueleto corporal o video proporcionarán trayectorias de movimiento y referencias externas, pero no se utilizarán necesariamente para la calibración física.

\subsection{E3.4: respiración}

Se impondrá
\begin{equation}
x_{\rm chest}(t)=A_r\sin(2\pi f_rt+\phi)
\end{equation}
como estímulo inicial del fantoma, con amplitud pico $A_r\in[0.5,10]$ mm y
frecuencia $f_r\in[0.1,0.6]$ Hz. Estos rangos definen la intervención mecánica de
Esc. E2, no una imposición de amplitud torácica a participantes. Después se utilizarán
formas no sinusoidales, pausas mecánicas y, en Esc. E3, referencias respiratorias
reales autorizadas. Se contrastarán forma de onda, amplitud, fase y frecuencia.
BreatheSmart orientará condiciones controladas y Wi-Actigraph aportará contexto de
sueño y movimiento; su función no sustituye la referencia sincronizada del banco.

\subsection{E3.5: movimiento + respiración}

Se combinarán locomoción y postura con el micromovimiento torácico. Se compararán un método de filtrado de referencia, métodos de separación de fuentes y representaciones contrastivas inspiradas en MoVi-Fi.

\subsection{E3.6: varios usuarios}

La secuencia incluirá dos personas con movimientos distintos, señales respiratorias
con frecuencias diferentes o similares y movimiento macroscópico simultáneo, según
Esc. E4. Las interrupciones de una fuente se estudiarán primero en simulación o con
fantomas; no se interpretarán como validación de apnea humana. Antes de evaluar las
redes se verificará el condicionamiento de la FIM conjunta.

\section{Línea 4: modelos de observación}

\subsection{Wi-Fi}

Se generarán configuraciones en 2.4, 5 y 6 GHz, con anchos de banda de 20, 40 y 80 MHz, y se variarán las subportadoras y los enlaces. Se añadirán perturbaciones controladas: CPO, STO, SFO, ruido, pérdida de paquetes y variación entre dispositivos. UWB-Fi se considerará como una extensión para evaluar diversidad frecuencial dispersa.

\subsection{LoRa}

Se generarán muestras I/Q para distintas configuraciones de $f_c$, factor de dispersión, ancho de banda y calendario de paquetes. Se estudiarán los desajustes de frecuencia, la interferencia LoRa/FSK, las colisiones y el tráfico irregular. OpenLoRa permitirá validar la estadística de radio y de las perturbaciones instrumentales, mientras que SenLoRa y los estudios de sensado LoRa de largo alcance orientarán los experimentos con personas y tráfico ambiente.

\subsection{Experimento comparativo entre tecnologías de radio}

La misma geometría humana $G_H(t)$ se renderizará bajo Wi-Fi y LoRa, produciendo $\calP_{\rm WiFi}$ y $\calP_{\rm LoRa}$ por separado. Se compararán información de detección $D_{\rm presence}$ y matrices normalizadas $\mat F_{\rm motion}$ y $\mat F_{\rm respiration}$. Presencia se tratará como hipótesis discreta; movimiento y respiración como estados continuos.

\section{Línea 5: solución del problema inverso}

\subsection{Inversión por optimización}

Para estado $\vect q$,
\begin{equation}
\vect q^\star=\arg\min_{\vect q}
D\left(Y_{\rm real},F_\theta(\vect q)\right)+R(\vect q).
\end{equation}
Se estudiará inicialización jerárquica: posición/orientación gruesa y posteriormente micromovimiento.

\subsection{Inversión amortizada}

Un codificador aproxima
\begin{equation}
q_\phi(\vect q\mid Y)
\end{equation}
y produce media y covarianza. La calibración de incertidumbre se evaluará con cobertura empírica de intervalos y error de calibración.

\subsection{Abstención}

Si la FIM es degenerada o la incertidumbre posterior supera un umbral, el sistema debe poder abstenerse. La tasa de cobertura--riesgo será una métrica principal, especialmente en fisiología.

\section{Línea 6: diseño y despliegue de la red}

Sea $u$ la configuración de nodos, ganancia, frecuencia y parámetros de radio. Se plantea
\begin{equation}
u^\star=\arg\max_u
\int_{\Omega_{\rm AoI}}
w(\vect x)
\Phi\left(\mat F_{\bar q}(\vect x;u)\right)d\vect x
\label{eq:deployment_opt}
\end{equation}
sujeto a restricciones de potencia, número de nodos, energía, tráfico y PRR. Para una selección discreta de enlaces $S$:
\begin{equation}
\max_{S:|S|\le K}
\sum_{\vect x\in\Omega}
\log\det\left(
\mat I+\sum_{l\in S}\mat F_{\bar q,l}(\vect x)
\right).
\label{eq:doptimal}
\end{equation}
La solución se comparará con la ubicación de nodos basada en RSSI/SNR y con heurísticas geométricas.

\section{Optimización multiobjetivo y funciones de pérdida}

\subsection{Por qué el problema es multiobjetivo}

La tesis no persigue una sola propiedad. La reconstrucción del canal, la coherencia de fase, la invariancia frente al dominio, la preservación de información fisiológica y la discriminación de perturbaciones instrumentales pueden inducir gradientes en conflicto. La función escalar básica es
\begin{equation}
\calL=\sum_i\lambda_i\calL_i,
\end{equation}
pero los pesos $\lambda_i$ no deben elegirse arbitrariamente.

\subsection{Pérdidas propuestas}

\begin{align}
\calL_H & = \operatorname{NMSE}(\widehat H,H),\\
\calL_A & = \||\widehat H|-|H|\|_1,\\
\calL_\phi & = \E[1-\cos(\angle\widehat H-\angle H)],\\
\calL_{\rm CIR} & = \|\widehat h-h\|_1,\\
\calL_{\rm dyn} & = \|\Delta_t\widehat H-\Delta_tH\|_1,\\
\calL_J & = \|\widehat J-J\|_F^2,\\
\calL_{\rm physio} & = 1-\rho(\widehat x_{\rm chest},x_{\rm chest})
+\lambda_{dtw}\operatorname{DTW}(\widehat x_{\rm chest},x_{\rm chest}),\\
\calL_{\rm align}&=\|z_{RF}^{\rm comun}-z_T^{\rm comun}\|_2^2,\\
\calL_{\rm cyc}&=\|\widehat z-z\|_2^2.
\end{align}

\subsection{Selección de pesos}

Se compararán tres estrategias: pesos fijos normalizados, ponderación adaptativa mediante GradNorm y optimización multiobjetivo mediante proyección de gradientes \citep{Chen2018GradNorm,Yu2020PCGrad}. Para dos objetivos con $g_i^Tg_j<0$, PCGrad proyecta
\begin{equation}
g_i\leftarrow g_i-
\frac{g_i^Tg_j}{\|g_j\|_2^2}g_j.
\end{equation}
No se supone de antemano que PCGrad sea superior; se incluye porque el posible conflicto entre la fidelidad del modelo directo y la utilidad para el sensado constituye una hipótesis científica explícita.

\subsection{Frontera de Pareto}

Se evaluará además error temporal y amplitud sin deformación del eje de tiempo
en prueba; la correlación y DTW no serán las únicas métricas de forma de onda.

Se analizará si minimizar $\calL_H$ coincide con minimizar $\calL_{\rm sensado}$. Si ambas metas difieren, se presentarán las soluciones de Pareto y se discutirá el grado de fidelidad física necesario para cada tarea. Esta evaluación puede revelar que un modelo de canal ligeramente menos preciso, pero con jacobianos correctos, resulte más adecuado para resolver el problema inverso de sensado.

\chapter{Diseño experimental, conjuntos de datos y protocolo de validación}
\label{ch:experimental_design}

\section{Función metodológica de los conjuntos de datos}

La selección de conjuntos de datos depende de observables, referencias,
ventanas temporales y condiciones de adquisición. MM-Fi aporta supervisión
multimodal de actividad y postura; eHealth CSI aporta actividad humana,
sin sustituir una referencia de expansión torácica sincronizada
\citep{Yang2023MMFi,Galdino2023eHealth}.

\begin{table}[ht]
\centering
\caption{Función de los conjuntos de datos y de las fuentes experimentales en la tesis.}
\begin{tabularx}{\textwidth}{p{4.0cm}L L}
\toprule
\textbf{Rol} & \textbf{Fuentes principales} & \textbf{Uso}\\
\midrule
Calibración física & DICHASUS dc01; escenas adicionales & Comparación del gemelo estático con mediciones reales.\\
Radio e instrumentación & OpenLoRa; metodología UWB-Fi & Desajustes de frecuencia, colisiones y muestreo.\\
Dinámica macroscópica & WiMANS, MM-Fi y datos con postura & Trayectorias de movimiento y evaluación fuera de distribución.\\
Fisiología controlada & BreatheSmart & Respiración con condiciones conocidas.\\
Sueño/movimiento & Wi-Actigraph & Evaluación según etiquetas de sueño y movimiento disponibles; sin extrapolación clínica.\\
Varios usuarios & WiMANS, WiMUSE, MURAL-Fi & Separabilidad y asociación.\\
Sensado humano con LoRa & SenLoRa; estudios a través de paredes & Respiración, marcha y tráfico ambiente.\\
Adquisición propia & laboratorio y piloto & Correspondencia con el gemelo y validación final.\\
\bottomrule
\end{tabularx}
\end{table}

\subsection{Fuentes adicionales y límites de validación}

La prioridad inmediata comprende DICHASUS para fase, CAEZ como segunda fuente CSI
sujeta a auditoría de sus subconjuntos, Indoor Radio Mapping para geometría--RSSI
y OpenLoRa para el operador de adquisición. El resto se incorporará cuando responda
a una hipótesis concreta, evitando confundir disponibilidad de datos con cobertura
experimental del problema doctoral.

\begin{table}[htbp]
\centering\small
\caption{Fuentes complementarias: evidencia disponible y uso propuesto.}
\label{tab:additional_datasets}
\begin{tabularx}{\textwidth}{P{3.0cm}L L}
\toprule
Fuente & Uso previsto & Límite de interpretación\\
\midrule
CAEZ \citep{ETHCAEZ,Zumegen2024CAEZWiFi} & Contraste de representaciones con CSI real adicional. & Auditar radios, calibración, geometría y sesiones antes de definir transferencia.\\
DICHASUS d010 \citep{DICHASUSd010Web} & Sensibilidad de $\mathsf V_i$ ante el movimiento medido del transmisor. & Es otra campaña y geometría; no representa movimiento humano ni valida respiración.\\
Indoor Radio Mapping \citep{Milosheski2026IndoorMapping} & Registro de nubes de puntos y predicción de RSSI. & RSSI no valida fase ni micromovimiento.\\
WiSegRT \citep{Zhang2023WiSegRT} & CIR sintética en escenas segmentadas; ablaciones geométricas. & Simulación multiescena, no evidencia de transferencia real.\\
WiInSim \citep{QualcommWiInSim} & Diversidad de escenas y propagación sintética. & Las particiones deben reservar escenas; no sustituye medición.\\
RF-MatID \citep{Chen2026RFMatID} & Identificación de materiales y evaluación de etiquetas/priores. & Clasificar un material no identifica su permitividad efectiva en el recinto.\\
DeepSense 6G \citep{Alkhateeb2022DeepSense} & Diseño multimodal y transferencia en contextos propios del conjunto. & Su contexto de movilidad no es un banco de respiración interior.\\
\bottomrule
\end{tabularx}
\end{table}

BreatheSmart orienta el diseño de perturbaciones respiratorias controladas; sus
etiquetas o resultados de clasificación no equivalen a la referencia mecánica sincronizada
de Esc. E2. MM-Fi, WiMANS, WiMUSE y MURAL-Fi se utilizarán para representación humana,
asociación y generalización según sus anotaciones, después de validar el observable.
MoViFi y ViMo mantienen su naturaleza radar y no se reclasifican como mediciones CSI.
La literatura de respiración y varios usuarios seguirá evaluándose en el nivel de
la tarea; no servirá como sustituto de geometría y canal medidos conjuntamente.

En d010 se auditarán primero las marcas temporales, la posición, el muestreo y las
compensaciones disponibles. Sólo después se construirán pares cuasiestáticos y pares
con desplazamiento del transmisor, sin compartir capturas entre particiones. Para cada
$\mathsf V_i$ se compararán distancia frente a desplazamiento, tamaño del efecto,
intervalos agrupados y capacidad de distinguir cambio físico de repetición. Esta prueba
complementa dc01, pero no estima $\partial H/\partial\mathbf q_H$: cambia el terminal
móvil y potencialmente todas sus trayectorias, no un dispersor humano con Tx/Rx fijos.

\section{Principio de no-oráculo}

Los valores de referencia se utilizarán para la evaluación y la supervisión explícita cuando corresponda, pero no para seleccionar a posteriori el canal, la subportadora o la componente que minimice el error. Toda selección destinada al despliegue deberá basarse en métricas disponibles sin recurrir a esos valores, como la observabilidad estimada, la estabilidad, la incertidumbre o criterios establecidos con el conjunto de validación.

Se permiten referencias oráculo claramente identificadas para diagnosticar límites de alineación, como S4.2b, pero no se interpretan como procedimientos de despliegue. En S4.2d se reserva un grupo de antenas para evaluar transferencia del retardo; el CPO se ajusta con las antenas evaluadas, de modo que se transfiere únicamente el retardo. En S4.2e, retardo y CPO se ajustan sólo en las frecuencias de entrenamiento y se transfieren sin reajuste a las de prueba. Una partición por antenas o frecuencias dentro de la misma observación no sustituye la evaluación en posiciones o sesiones independientes.

\section{Referencia de fase, trazabilidad y evaluación relativa}

En S4.6, $\phi_0^\star$ y $\psi$ se construyen con el canal medido y retardo
oráculo. Su concentración y la reducción tras centrado son diagnósticos;
no se presentan como rendimiento predictivo independiente. Los estadísticos
marginales por antena se calculan descriptivamente en el conjunto completo.
Si se reutilizan como calibración de un método de despliegue deberán volver
a estimarse sólo con entrenamiento y quedar congelados durante la prueba.

S4.6d.1 recupera el tiempo físico a partir de coordenadas únicas en el
subconjunto de 10\,000 capturas. El error XYZ nulo acredita la correspondencia
de registros, no la exactitud física absoluta del posicionamiento o del reloj.
Las observaciones se ordenan por su marca temporal; los índices del TFRecord
no sustituyen el tiempo. Los dos arreglos comparten identificador y tiempo de
captura y no se cuentan como duplicados temporales independientes.

S4.7a evaluará primero cada arreglo por separado. Se fijarán máscara, frecuencias,
pares entre frecuencias, normalización y tratamiento del ruido antes de comparar
métodos. Los productos espaciales $G$ se utilizarán para fidelidad relativa;
los productos $K$ y la distancia respecto de fase común, para separar esa
fidelidad de la contribución del retardo. La identidad de $G$ con y sin retardo
común será un control de implementación. Se conservarán errores brutos y de
potencia para que una invariancia amplia no oculte una limitación del modelo.

Las pruebas negativas de continuidad se acompañarán, en su extensión
confirmatoria, de intervalos pareados, réplicas aleatorias y conteos de pares
por intervalo. Se distinguirá ausencia de ventaja observada de equivalencia
estadística y se reservarán campañas o referencias adicionales para contrastar
el origen instrumental. La nomenclatura S4.6d.2 identifica el último bloque
realizado; S4.7a y S4.7b identifican protocolos propuestos.


\section{Escenarios físicos E0--E5 y criterios de avance}
\label{sec:physical_scenarios}

La campaña propia mantendrá transmisor y receptor fijos durante los contrastes de
movimiento de objetos. Se identificarán geometría, antenas, polarización, potencia,
forma de onda y referencia de adquisición. El estado documentado termina en S4.6d.2;
los siguientes escenarios especifican un programa por ejecutar. Los tamaños y rangos
iniciales se revisarán mediante un piloto de precisión antes del estudio confirmatorio.

\subsection{Esc. E0: caracterización instrumental}
\label{sec:scenario_e0}

Se adquirirán intervalos iniciales de 10--30 minutos en habitación vacía, con terminales
y cables inmovilizados. Se repetirán condiciones de encendido, reconexión, cambio de
día y dispositivo, registrando temperatura y tiempo de estabilización. Se modificará
un factor principal a la vez; el orden se alternará o aleatorizará cuando sea posible,
para no confundir deriva temporal con condición experimental. Se conservará una
condición de referencia al inicio y al final de cada bloque.

Se compararán canal bruto, corrección aprendida, referencia física calibrada y productos
relativos. Las salidas serán deriva de fase y retardo, estabilidad de amplitud,
variabilidad entre cadenas, pérdidas de paquetes y distribución de $D_{\rm rep}$.
La estabilidad intrasesión se separará de reproducibilidad tras reinicio. Los cocientes
sólo se interpretarán bajo sus supuestos de perturbación común. El criterio de paso
será conocer el piso instrumental y la incertidumbre temporal para la escala de
movimiento de Esc. E1; si éstos son mayores que la señal esperada, se corregirá el
banco o se ampliará justificadamente la perturbación antes de aprender un estimador.

\subsection{Esc. E1: reflector con desplazamiento medido}
\label{sec:scenario_e1}

Se utilizará una etapa lineal con encoder o micrómetro calibrado y fijación rígida.
Se propone una rejilla inicial de desplazamientos de 0, 0.25, 0.5, 1, 2, 5 y 10 mm,
restringida a pasos resolubles por la metrología real. La posición comandada al motor
no será la referencia: se medirá el desplazamiento, incluyendo holgura, histéresis
y retorno al origen. Se ensayarán ambos sentidos y al menos 20 repeticiones iniciales
por condición; esta cifra caracteriza repetibilidad y no sustituye un cálculo de
precisión o potencia estadística a partir del piloto.

Se considerarán una trayectoria dominante, una reflexión controlada y multitrayectoria,
variando orientación del reflector y posición respecto de zonas de Fresnel. Las
condiciones nulas y movimientos de baja sensibilidad geométrica se conservarán como
controles. Se evaluarán $\Delta H$ con referencia común, $\Delta\mathsf V_i$,
$\widehat J_{\mathcal V,q}$, error de fase inducida y recuperación del desplazamiento.
LiDAR y radar complementan el contexto geométrico y dinámico, pero no reemplazan
la referencia mecánica submilimétrica. El criterio de paso será una respuesta física
repetible distinguible de E0, con incertidumbre y discrepancia del modelo cuantificadas,
no únicamente una correlación elevada de curvas filtradas.

\subsection{Esc. E2: fantoma respiratorio}
\label{sec:scenario_e2}

Se programará una perturbación mecánica periódica con amplitudes iniciales de
0.5--10 mm y frecuencias de 0.1--0.6 Hz, más formas no sinusoidales, pausas y
movimiento lento superpuesto. Se declarará si la amplitud es pico o pico a pico;
por defecto será pico respecto del estado medio. Se medirán posición y tiempo con
referencia independiente, sin identificar automáticamente estos rangos con desplazamiento
torácico de toda persona. El periodo de observación contendrá suficientes ciclos de
la frecuencia mínima y se registrará la duración exacta, así como los transitorios.

Se compararán magnitud, rama coherente, productos relativos y fusión, primero con
estimadores clásicos y después con modelos aprendidos. Las métricas incluirán error
de desplazamiento, frecuencia, amplitud, desfase temporal, coherencia y forma de onda,
así como detección de pausas mecánicas. Filtros, bandas y reglas de selección se fijarán
con validación; toda alineación temporal diagnóstica adicional se reportará por separado.
Recuperar sólo BPM puede ocultar una forma de onda incorrecta o el armónico dominante.
Se avanzará a Esc. E3 cuando la referencia temporal y física esté caracterizada y
la recuperación sea reproducible en condiciones reservadas. El resultado será una
validación mecánica, sin afirmaciones clínicas sobre apnea.

\subsection{Esc. E3: una persona}
\label{sec:scenario_e3}

Tras la aprobación ética y el consentimiento correspondientes, se separarán respiración
en reposo, cambios de postura/orientación, marcha y movimiento de brazos; posteriormente
se combinarán perturbaciones. La referencia de expansión respiratoria deberá sincronizarse
y caracterizarse. RGB-D o visión aportarán postura y actividad cuando el protocolo lo
permita, y el radar contrastará dinámica desde vistas calibradas. Se analizarán artefactos,
fallos de referencia y regiones no observables antes de entrenar un clasificador.

Se reservarán participantes y sesiones completos, y se reportará desempeño por condición,
incertidumbre, abstenciones y tasa de registros excluidos. La inferencia final indicará
expresamente sus modalidades necesarias. Una prueba de concepto de laboratorio no
se presentará como diagnóstico médico ni como validación del despliegue comunitario.

\subsection{Esc. E4: dos personas}

Después de Esc. E1--E3, se evaluarán dos personas estáticas, una en movimiento,
ambas con respiración observable y posiciones cruzadas u ocluidas. Primero se estudiará
separabilidad de las respuestas y condicionamiento del Jacobiano conjunto; después,
asociación de fuentes y errores individuales. Se comparará contra una fuente dominante
y contra condiciones de señales semejantes. Tres vistas radar, si se dispone de ellas,
requerirán registro y sincronización; su número por sí solo no garantiza separación.
El criterio de éxito será delimitar condiciones recuperables y fallos, antes de ampliar
usuarios o tareas fisiológicas.

\subsection{Esc. E5: transferencia}

Se introducirán cambios de día, dispositivo, antena, mobiliario y habitación de manera
separable antes de combinarlos. Se compararán inferencia sin adaptación, calibración
física y adaptación con presupuestos crecientes de datos. La evaluación multibanda
repetirá el experimento canónico del reflector y el fantoma, ajustando el rango de
movimiento a la longitud de onda y documentando antenas, ancho de banda y potencia.
Un cambio de banda modifica el canal físico: la transferencia debe conservar relaciones
físicas o eficiencia de adaptación, no exigir idénticos valores de $H$.

Una extensión con la UPNA y el Dr. Francisco Falcone se plantea como colaboración
por concretar, para repetir Esc. E1--E2 en las bandas e instrumentación que se acuerden,
incluidas eventualmente ondas milimétricas o sub-THz. No constituye una campaña ya
realizada ni una disponibilidad instrumental o financiación confirmadas. La comparación
externa empleará el mismo protocolo, referencias y criterios de exclusión del banco local.

\section{Trazabilidad del conjunto de datos propio}
\label{sec:own_dataset_schema}

Cada registro se asociará a la clave escena, distribución, sesión, captura, dispositivo
y enlace. La relación entre archivos se establecerá mediante identificadores persistentes,
no por el orden de lectura. Se conservarán muestras RF originales y productos derivados
separados, con versiones de procesamiento y sumas de verificación. Las modalidades
visuales y humanas seguirán las condiciones de retención, acceso y seudonimización
aprobadas; conservar señal RF cruda no exige conservar imágenes identificables sin límite.

\begin{table}[htbp]
\centering\small
\caption{Especificación lógica del conjunto multimodal propuesto.}
\label{tab:own_dataset_schema}
\begin{tabularx}{\textwidth}{P{2.5cm}L}
\toprule
Componente & Contenido mínimo y finalidad\\
\midrule
Escena & Nubes de puntos, mallas G0--G4, materiales, dimensiones de control,
transformaciones y unidades; reconstrucción geométrica auditable.\\
Radio & CSI Wi-Fi, I/Q LoRa/SDR, telemetría de red y datos radar disponibles;
forma de onda, ganancias, canales, paquetes y máscaras.\\
Referencia física & Posición mecánica medida, postura autorizada, expansión respiratoria
y eventos; incertidumbre y procedimiento de calibración.\\
Simulación & Trayectorias, CIR/CFR, configuración RT, semillas y versiones;
correspondencia de consulta con cada captura.\\
Calibración & Relojes, cables, antenas, referencia coherente, sesgos y temperatura;
parámetros ajustados y subconjunto utilizado.\\
Metadatos & Identidad de sesión y dispositivo, participante seudonimizado,
particiones, control de calidad, procedencia y sincronización.\\
\bottomrule
\end{tabularx}
\end{table}

El flujo por sesión incluirá comprobación de montaje y referencia, registro basal,
bloques con condiciones identificadas y control final de deriva. Se anotarán cambios
de AGC y pérdidas de captura, evitando rellenar silenciosamente intervalos ausentes.
Cada canal temporal conservará su reloj original y su transformación al reloj global.
Los modelos geométricos, parámetros RF y normalizaciones ajustados se congelarán antes
de aplicar el protocolo de prueba.

\section{Particiones para la evaluación}

Se reportarán como mínimo:
\begin{enumerate}
\item posición no observada durante el entrenamiento;
\item participante no observado durante el entrenamiento;
\item distribución de mobiliario no observada;
\item dispositivo no observado;
\item habitación no observada;
\item transferencia de interiores a exteriores;
\item transferencia de uno a varios usuarios, cuando proceda.
\end{enumerate}
No se distribuirá al azar entre entrenamiento y prueba información altamente correlacionada de un mismo segmento o posición.

\section{Repetibilidad y estadística}

Cada configuración se repetirá en varias sesiones y con distintas semillas de entrenamiento. Se informarán las distribuciones, los intervalos de confianza y los tamaños del efecto cuando corresponda. Para comparar métodos entre condiciones se priorizarán pruebas pareadas y correcciones por comparaciones múltiples, evitando conclusiones basadas sólo en la mejor ejecución.

\section{Ablaciones acumulativas}

Cada experimento o publicación deberá identificar el aporte de un componente concreto. Por ejemplo:
\begin{equation}
\begin{aligned}
\text{geometría y RT nominal}
&\rightarrow \text{calibración macroscópica}\\
&\rightarrow \text{referencia y factorización de fase}\\
&\rightarrow \text{ramas robusta y coherente}\\
&\rightarrow \text{fidelidad diferencial}
\rightarrow \text{representación humana}.
\end{aligned}
\end{equation}
La misma lógica se aplicará a LoRa: I/Q ideal $\rightarrow$ CFO y ruido $\rightarrow$ temporización de paquetes $\rightarrow$ interferencia $\rightarrow$ tráfico ambiente.

\section{Programa experimental y criterios de validación}

\subsection{Experimento A: gemelo nominal para dc01}

\textbf{Pregunta:} ¿la geometría y la implementación están correctamente alineadas con mediciones reales?

\textbf{Entradas:} escena INUE, posiciones DICHASUS, CSI corregido y trayectorias RT.

\textbf{Salidas:} $H_{\rm real}$, $H_{\rm RT}$, PDP, RMS-DS y errores por posición.

\textbf{Criterio de validación:} flujo de procesamiento completamente reproducible y discrepancia entre simulación y medición cuantificada. No se exige una coincidencia perfecta.

\textbf{Avance reportado:} banco de $10000$ posiciones y métricas nominales S1/B0 disponibles en el \cref{ch:avances}. Se documentan versiones, escala, mecanismos activos, partición original y comprobaciones de correspondencia de identificadores; la reproducción independiente debe utilizar los mismos artefactos.

\subsection{Experimento B: calibración física y de fase}

\textbf{Pregunta:} ¿qué parte de la discrepancia entre simulación y medición proviene de los materiales y qué parte de errores geométricos o de fase?

\textbf{Métodos de comparación:} parámetros ITU, ajuste exclusivo de escala, materiales globales y neuronales, diagnóstico de fase constante/retardo y calibración con incertidumbre por trayectoria.

\textbf{Criterio de validación:} mejora consistente en la partición espacial de prueba, no sólo en la de entrenamiento.

\textbf{Avance reportado:} B1/B2 mejoran las métricas macroscópicas; S4.2--S4.3 caracterizan fase afín y residual. S4.4--S4.5 evalúan predicción de retardo y efecto sobre el canal bajo particiones densas y por bloques. La generalización material de B1/B2 no se confunde con la del predictor de retardo sobre B0. S4.6 añade la caracterización de la fase común, la recuperación de marcas temporales y los controles de vecindad. PEAC se ha reproducido en un ejemplo controlado; su aplicación a dc01 permanece propuesta.

\subsection{Experimento C: eficiencia de datos}

\textbf{Pregunta:} ¿cuántas mediciones necesita cada representación?

\textbf{Salida:} curvas error vs. número de posiciones de calibración, $N_{90}$ y área bajo curva.

\subsection{Experimento D: fidelidad diferencial}

\textbf{Pregunta:} ¿el gemelo predice los cambios producidos por perturbaciones pequeñas?

\textbf{Diseño:} reflector mecánico con desplazamiento de referencia conocido.

\textbf{Métrica:} $E_J$, correlación de $\Delta H$ y error de fase inducida.

\subsection{Experimento E: representación realista de la capa física de LoRa}

\textbf{Pregunta:} ¿el modelo de observación reproduce la estructura estadística de las muestras I/Q reales?

\textbf{Conjuntos de datos:} OpenLoRa y capturas propias.

\textbf{Métricas:} PSD, CFO, espectro después de retirar el chirrido, SNR, detección de paquetes y distancias entre distribuciones.

\subsection{Experimento F: misma perturbación, dos tecnologías de radio}

\textbf{Pregunta:} ¿qué radio conserva mejor presencia, movimiento y respiración bajo la misma escena?

\textbf{Diseño:} perturbación física común y dos modelos de observación.

\textbf{Métrica:} FIM, recuperación de la forma de onda y robustez frente a perturbaciones instrumentales.

\subsection{Experimento G: transferencia de simulación a mediciones reales}

\textbf{Pregunta:} ¿el preentrenamiento con información física reduce la cantidad de datos reales necesaria?

\textbf{Comparaciones:} entrenamiento exclusivo con datos reales, entrenamiento exclusivo con datos sintéticos, datos sintéticos con adaptación y datos sintéticos con destilación multimodal.

\subsection{Experimento H: varios usuarios}

\textbf{Pregunta:} ¿cuándo son físicamente separables dos fuentes humanas?

\textbf{Análisis:} valores singulares del jacobiano y de la FIM conjunta, separación ciega de fuentes y modelos aprendidos.

\subsection{Experimento I: optimización del despliegue}

\textbf{Pregunta:} ¿la ubicación de nodos y la ganancia orientadas a la tarea superan la planificación basada en RSSI?

\textbf{Escenarios:} red distribuida en interiores y, posteriormente, un piloto en exteriores.

\section{Criterios confirmatorios y resultados negativos}

Los umbrales numéricos de error y la mejora mínima relevante se fijarán a partir de
precisión instrumental y objetivos de la tarea, antes de abrir la prueba confirmatoria.
No se escogerán para hacer pasar un método después de observar sus resultados. El
piloto servirá para estimar dispersión y número de sesiones requeridas; las semillas
de aprendizaje no sustituyen réplicas físicas independientes.

Las comparaciones serán pareadas por estado y sesión, con remuestreo agrupado y
análisis por región, orientación y potencia. Se informará el número de sesiones,
posiciones, participantes, pares y exclusiones que sustentan cada estadístico. Para
S4.7a se conservarán las 4900 posiciones de prueba de la partición densa y los cinco
bloques de 2000 posiciones del protocolo espacial, con su entrenamiento correspondiente;
no se confundirá este último total de 10\,000 evaluaciones con la partición original.
B1/B2 requerirán reentrenamiento bajo bloques antes de atribuirles esa generalización.

Si la representación robusta mejora estabilidad pero reduce la respuesta mecánica,
se conservará como rama de contexto y se recurrirá a la rama coherente para la tarea.
Si la referencia no es estable, se revisará E0; si el RT acierta potencia pero falla
la respuesta diferencial, se revisarán geometría local, trayectorias, dispersión y
operador antes de aumentar capacidad neuronal. Si las dos fuentes no son separables,
se informará la limitación y se evaluará otra disposición de enlaces. Estas decisiones
permiten que un resultado negativo refine la metodología sin transformarse en una
conclusión universal sobre la imposibilidad del sensado.

\section{Síntesis gráfica prevista}

La narrativa experimental se condensará en cinco figuras cuando existan los resultados
correspondientes: (i) taxonomía de fase física, geométrica, instrumental y de
procesamiento; (ii) matriz cuantitativa $\mathsf V_0$--$\mathsf V_7$ frente a
$\mathsf T_1$--$\mathsf T_7$; (iii) frente de Pareto entre rechazo de perturbaciones
y sensibilidad geométrica, usando d010 sólo si supera la auditoría; (iv) fidelidad
B0--B2 bajo protocolos denso y bloqueado, distinguiendo modelos reentrenados; y
(v) transferencia desde DICHASUS hacia perturbaciones controladas, conjunto externo
y sensado humano futuro. Estas figuras no se rellenarán con resultados esperados ni
mezclarán estadísticas incompatibles: cada panel declarará datos, partición, unidad
de réplica y acceso a referencias.

\chapter{Avances experimentales del modelado de canal: calibración,
estructura de fase y transferencia espacial}
\label{ch:avances_rf_twin_completo}
\label{ch:avances}

% ----------------------------------------------------------------------
\FloatBarrier
\section{Objetivo científico del bloque experimental}
% ----------------------------------------------------------------------

La secuencia experimental desarrollada hasta ahora estudia una pregunta que
precede al objetivo final de sensado fisiológico:

\begin{quote}
\emph{¿Qué parte de la respuesta compleja de un canal inalámbrico puede
reconstruirse, calibrarse, descomponerse y predecirse a partir de un gemelo
electromagnético antes de introducir explícitamente personas, movimiento o
fisiología?}
\end{quote}

La motivación es metodológica. En un sistema de sensado por radiofrecuencia, el observable que
finalmente contiene información humana es una perturbación pequeña del canal.
La concordancia estática proporciona una primera comprobación del entorno y de
la cadena de síntesis. Su insuficiencia exige delimitar qué predicciones siguen
siendo contrastables mediante observables relativos o perturbaciones controladas.
Sin embargo, la condición opuesta tampoco es suficiente: un modelo que
reconstruye bien el canal estático puede no preservar la sensibilidad del canal
frente a una perturbación humana. Por ello, la tesis no busca únicamente
minimizar una pérdida de reconstrucción, sino descomponer la fidelidad en
propiedades físicamente distintas.

La progresión experimental seguida puede resumirse como

\begin{equation}
\begin{aligned}
\Gamma&\longrightarrow\mathcal P_{\rm RT}\longrightarrow H_{\rm RT}
\longrightarrow\Delta\phi,\\
\Delta\phi&\longrightarrow\{\Delta\tau_{\rm eff},\phi_0,r_\phi\}
\longrightarrow\widehat{\Delta\tau}_{\rm eff}\longrightarrow\widetilde H_{\rm RT}.
\end{aligned}
\label{eq:avance_general}
\end{equation}

donde $\Gamma$ representa la escena, $\mathcal P_{\rm RT}$ el conjunto de
trayectorias trazadas, $H_{\rm RT}$ la respuesta compleja sintetizada,
$\Delta\tau_{\rm eff}$ una pendiente espectral efectiva, $\phi_0$ un desfase
constante residual y $r_\phi$ el componente no explicado por la pendiente común y los interceptos
por enlace. El número de parámetros depende de los ejes de ajuste; no se trata
de dos parámetros escalares para todo el tensor.

Este programa se relaciona directamente con cuatro líneas de literatura.
Primero, Hoydis et al.~\cite{Hoydis2024DiffRT} muestran que un trazador de rayos
diferenciable permite calibrar propiedades electromagnéticas, dispersión difusa y
patrones de antena usando mediciones. Segundo, Ruah et
al.~\cite{Ruah2024Phase} muestran que pequeñas discrepancias geométricas pueden
hacer que la fase simulada de cada trayectoria sea suficientemente incorrecta
como para sesgar la calibración, motivando métodos con incertidumbre de fase. Tercero,
NeRF$^2$~\cite{Zhao2023NeRF2}, Learnable Wireless Digital
Twins~\cite{Jiang2025LearnableWDT} y RadTwin~\cite{Zhang2026RadTwin} proponen
representaciones neuronales alternativas del campo RF o de las interacciones
electromagnéticas. Cuarto, WaveVerse~\cite{Zheng2026WaveVerse} muestra que la
coherencia temporal de fase resulta crítica cuando el simulador se utiliza en
escenarios dinámicos y de sensado.

La contribución experimental actual caracteriza la estructura de discrepancia
y su transferencia parcial. La conservación conjunta del canal y de sus
perturbaciones humanas constituye la hipótesis que esos avances permiten
formular con mayor precisión, y permanece sujeta a validación diferencial.

% ----------------------------------------------------------------------

\FloatBarrier
\section{Alcance del registro y convenciones de evaluación}
\label{sec:metricas_avances}

Este capítulo integra los resultados documentados hasta S4.6d.2, con precisión
numérica conservada en las tablas de cada ejecución. El análisis de S4.3--S4.5
se realiza sobre la referencia B0, salvo indicación expresa; no debe extenderse
su validación regional a materiales B1/B2 sin reentrenarlos bajo ese protocolo.
La ejecución documentada no equivale a una reproducción independiente de
los tensores. Los scripts y artefactos que permiten verificarla se identifican
en la \cref{sec:artifacts}.

Para cada enlace $v=(i,u,a)$ y conjunto de frecuencias válidas $\mathcal K_v$,
se definen
\begin{align}
E_{\phi,v}&=\frac{180}{\pi|\mathcal K_v|}\sum_{k\in\mathcal K_v}
\left|\wrap_{(-\pi,\pi]}\bigl(\angle H^{\rm real}_{v,k}-\angle\widehat H_{v,k}\bigr)\right|,
\label{eq:error_angular_avances}\\
\rho_v&=\frac{\left|\sum_{k\in\mathcal K_v}H^{\rm real}_{v,k}\widehat H_{v,k}^*\right|}
{\sqrt{\sum_k|H^{\rm real}_{v,k}|^2}\sqrt{\sum_k|\widehat H_{v,k}|^2}}.
\label{eq:coherencia_avances}
\end{align}
La mediana del MAE angular por enlace difiere de la mediana de todos los
residuos elementales. Las tablas que indican fase mediana o coherencia mediana
resumen los indicadores de enlace; los cuantiles deben interpretarse sobre
esa misma unidad. Las máscaras de amplitud nula o insuficiente, su umbral y
las exclusiones deben acompañar los resultados, pues la fase no está definida
en amplitud nula.

El NMSE de magnitud por enlace se expresa como
\begin{equation}
\mathrm{NMSE}_{|H|,v}=10\log_{10}
\frac{\sum_k(|\widehat H_{v,k}|-|H^{\rm real}_{v,k}|)^2}
{\sum_k|H^{\rm real}_{v,k}|^2},
\end{equation}
y el complejo sustituye el numerador por
$\sum_k|\widehat H_{v,k}-H^{\rm real}_{v,k}|^2$. Se distinguen esas reducciones
por enlace de una razón de sumas globales. Las comparaciones de potencia y
RMS-DS conservan el MAE y correlaciones del registro; la conversión a dB y
el orden de agregación forman parte de la especificación del evaluador.

La coherencia $\rho_v$ es invariante a una rotación constante del enlace.
Por ello el CPO puede reducir el error angular sin modificar esa coherencia.
Para $R_{v,k}=H^{\rm real}_{v,k}\widehat H_{v,k}^*$, un retardo candidato
produce
\begin{equation}
C_v(\tau)=\sum_k R_{v,k}e^{j2\pi\nu_k\tau},\qquad
\widehat\phi_{0,v}(\tau)=\angle C_v(\tau),\qquad
\widetilde H_{v,k}=\widehat H_{v,k}e^{j\widehat\phi_{0,v}-j2\pi\nu_k\widehat\tau}.
\label{eq:cpo_delay_avances}
\end{equation}
El módulo del producto residual pondera frecuencias por las amplitudes de ambos
canales. Maximizar $\sum_a|C_{i,u,a}(\tau)|$ minimiza el error cuadrático
complejo respecto de las rotaciones permitidas; no equivale necesariamente a
maximizar la media no ponderada de coherencias normalizadas ni a minimizar
el error angular absoluto. Los resultados deben interpretarse bajo el objetivo
y los grados de libertad efectivamente utilizados.

En S4.2b todos los parámetros de fase utilizan la observación evaluada y
constituyen una referencia oráculo. S4.2d transfiere retardo entre antenas con
CPO libre en evaluación. S4.2e transfiere retardo e intercepto entre frecuencias
sin reajuste. S4.5 predice retardo desde descriptores y ajusta CPO con el canal
evaluado. La denominación de cada condición conserva esa diferencia de acceso.

\FloatBarrier
\section{Banco común S1: DICHASUS dc01 y gemelo estático nominal}
\label{sec:bank_dc01}
% ----------------------------------------------------------------------

\FloatBarrier
\subsection{Datos, dimensionalidad y escena}

El banco utiliza DICHASUS dc01, perteneciente a la familia dcxx. Esta familia
contiene mediciones de canal con posicionamiento y un modelo tridimensional
del entorno, por lo que permite contrastar una respuesta compleja real con un
modelo geométrico de propagación. En el procesamiento utilizado en esta tesis
se trabajan $N=10\,000$ posiciones. Para cada posición se conserva una
respuesta compleja medida

\begin{equation}
H_i^{\rm real}\in\mathbb C^{2\times 32\times 1024},
\end{equation}

correspondiente a dos arreglos de 32 antenas y 1024 subportadoras
\citep{Euchner2023DICHASUSdcxx}. La portadora
de la escena es \SI{3.438}{GHz} y el ancho de banda es \SI{50}{MHz}.

La unidad básica de trabajo puede representarse como

\begin{equation}
\mathcal D_{S1}
=
\left\{
\left(
\mathbf p_i,\,
H_i^{\rm real},\,
\mathcal P_i^{\rm RT}
\right)
\right\}_{i=1}^{10000}.
\label{eq:s1_data_object}
\end{equation}

El trazado nominal se ejecutó con profundidad máxima cinco y produjo, como
máximos estructurales del TFRecord, 198 trayectorias especulares, 6 difractadas
y 445 de dispersión difusa. El tensor almacenado se validó para las 10\,000
posiciones. Aunque las trayectorias de dispersión difusa se conservan en el archivo,
los canales evaluados en B0--B2 y S4 se reconstruyen con
\path{cir(scattering=False)}; por ello, los caminos activos en el cálculo de
$H_{\rm RT}$ son las contribuciones especulares y difractadas.

Con $\nu_k=f_k-f_{\rm ref}$ y ganancias complejas que incorporan la fase
de propagación referida a $f_{\rm ref}$, el canal sintetizado tiene la forma

\begin{equation}
H_{i,u,a}(f_k)
=
\sum_{\ell\in \mathcal P_{i,u,a}}
\alpha_{i,u,a,\ell}(f_k)
\exp\left[-j2\pi\nu_k\tau_{i,u,a,\ell}\right],
\label{eq:synth_channel}
\end{equation}

donde $u\in\{1,2\}$ identifica el arreglo y $a\in\{1,\ldots,32\}$ su
antena. Las etiquetas Tx1/Tx2 de los registros se conservan como identificadores
de los dos bloques de enlaces; no implican dos transmisores móviles adicionales.
Si el simulador invierte los extremos por reciprocidad, debe mantenerse explícita
la correspondencia entre índices de adquisición y de síntesis.

\FloatBarrier
\subsection{Entorno de software y trazabilidad}

La ejecución principal se realizó con Python 3.10, TensorFlow 2.15.1 y
Sionna 0.19.2. El trazado de rayos se apoya en la escena
\path{inue_simple}. Debido a diferencias de API respecto del repositorio
original, se introdujo un adaptador de compatibilidad para el formato de
\path{trace_paths}, pero el módulo de síntesis posterior quedó congelado. La respuesta
simulada utilizada en B1/B2 se calcula mediante:

\begin{enumerate}
\item colocación de receptores;
\item reconstrucción de los caminos trazados;
\item \path{scene.compute_fields};
\item \path{paths.cir(scattering=False)};
\item conversión a OFDM mediante \path{cir_to_ofdm_channel};
\item división por $\sqrt{1024}$;
\item escalamiento global equivalente del canal medido.
\end{enumerate}

La escala nominal congelada fue

\begin{equation}
\alpha_{B0}=5.762854637936243\times10^{-10},
\end{equation}

cuya conversión $10\log_{10}\alpha_{B0}$ es $-92.3936$ dB y

\begin{equation}
\sqrt{\alpha_{B0}}
=
2.4005946425700954\times10^{-5}.
\end{equation}

La escala es un factor de potencia de normalización y su raíz actúa sobre
amplitudes; no constituye por sí sola una potencia recibida en dBm. En esta
serie B0 ya incorpora la escala congelada: no representa una comparación sin
normalización. B1 modifica también la escala; su mejora no debe atribuirse
exclusivamente a materiales sin una ablación de ganancia. Deben conservarse
las convenciones de normalización de ambos canales y su conjunto de ajuste.

% ----------------------------------------------------------------------
\FloatBarrier
\section{B0--B2: reconstrucción estática del entorno}
\label{sec:b012_detailed}
% ----------------------------------------------------------------------

\FloatBarrier
\subsection{B0: RT nominal con materiales ITU}

B0 utiliza geometría y materiales nominales. Su objetivo no es ser el mejor
modelo, sino fijar un punto de referencia para cuantificar la discrepancia
entre una escena plausible y la medición real.

Sobre 4900 posiciones de prueba se obtuvieron:

\begin{align}
\rho_P &= 0.699401,\\
{\rm MAE}_P &= 6.290937~{\rm dB},\\
\rho_{\tau_{\rm RMS}} &= 0.173353,\\
{\rm MAE}_{\tau_{\rm RMS}} &= 43.211357~{\rm ns},\\
{\rm mediana~NMSE}_{|H|} &= -3.173375~{\rm dB},\\
{\rm mediana~coh} &= 0.156483,\\
{\rm mediana~MAE_{\phi}} &= 89.925659^\circ.
\end{align}

El resultado es físicamente informativo: la escena nominal recupera una parte
considerable de la variación espacial de potencia, pero reproduce pobremente
la estructura temporal fina y prácticamente no está alineada en fase.

\FloatBarrier
\subsection{B1: calibración global de materiales}

B1 sigue la idea de RT diferenciable de Hoydis et
al.~\cite{Hoydis2024DiffRT}. Se optimizan propiedades globales por objeto o
material manteniendo fija la geometría. El entrenamiento usó lotes de ocho posiciones,
tasa de aprendizaje $10^{-3}$ y la partición original del banco: 5000
posiciones de entrenamiento, 100 de validación y 4900 de prueba. Las 10\,000
posiciones constituyen el banco completo, no el tamaño del conjunto de entrenamiento.

Los parámetros finales reportados fueron aproximadamente:

\begin{table}[!htbp]
\centering
\caption{Parámetros electromagnéticos aprendidos en B1.}
\label{tab:b1_materials}
\begin{adjustbox}{max width=\textwidth}
\begin{tabular}{lrr}
\toprule
Objeto/material & $\epsilon_r$ & $\sigma$ [S/m]\\
\midrule
Piso & 1.0010 & 0.0010\\
Techo & 70.4308 & 8.089956\\
Muros & 7.3262 & 1.952099\\
\bottomrule
\end{tabular}
\end{adjustbox}
\end{table}

La escala final fue $8.191072820018519\times10^{-10}$. Los valores ajustados
se interpretan como parámetros efectivos del modelo. En particular, la permitividad
elevada del techo no identifica por sí sola una propiedad intrínseca del material: puede
compensar geometría, patrones de antena u otras interacciones no modeladas.

En prueba:

\begin{align}
\rho_P &= 0.736297, &
{\rm MAE}_P &= 3.815674~{\rm dB},\\
\rho_{\tau_{\rm RMS}} &= 0.278345, &
{\rm MAE}_{\tau_{\rm RMS}} &= 27.320232~{\rm ns},\\
{\rm mediana~NMSE}_{|H|} &= -3.918742~{\rm dB},\\
{\rm coh} &= 0.155111,\\
{\rm MAE_{\phi}} &= 89.953903^\circ.
\end{align}

Por tanto, aprender parámetros materiales mejora claramente potencia,
dispersión de retardos y magnitud, pero no la fase bruta.

\FloatBarrier
\subsection{B2: materiales neuronales dependientes de posición}

B2 introduce una función neuronal espacial de 57\,860 parámetros que devuelve
propiedades electromagnéticas efectivas en función de la posición de
interacción:

\begin{equation}
\boldsymbol\theta_M(\mathbf x)
=
g_\psi\bigl(\gamma(\mathbf x)\bigr).
\end{equation}

Esta formulación permite absorber heterogeneidad espacial que un único valor
por material no puede representar. Es importante no confundir B2 con
NeRF$^2$: B2 conserva el soporte físico de trazado de rayos y aprende parámetros de
interacción/material; NeRF$^2$ aprende directamente una representación
continua del campo RF~\cite{Zhao2023NeRF2}.

B2 alcanza:

\begin{align}
\rho_P &= 0.870918, &
{\rm MAE}_P &= 2.223648~{\rm dB},\\
\rho_{\tau_{\rm RMS}} &= 0.646920, &
{\rm MAE}_{\tau_{\rm RMS}} &= 16.498459~{\rm ns},\\
{\rm mediana~NMSE}_{|H|} &= -4.455138~{\rm dB},\\
{\rm NMSE_{H}} &= 3.115020~{\rm dB},\\
{\rm coh} &= 0.145279,\\
{\rm MAE_{\phi}} &= 89.961143^\circ.
\end{align}

\begin{table}[!htbp]
\centering
\small
\caption{Comparación consolidada B0--B2.}
\label{tab:b012_consolidated}
\begin{adjustbox}{max width=\textwidth}
\begin{tabular}{lccc}
\toprule
Métrica & B0 & B1 & B2\\
\midrule
Correlación potencia $\uparrow$ & 0.699 & 0.736 & \textbf{0.871}\\
MAE potencia [dB] $\downarrow$ & 6.29 & 3.82 & \textbf{2.22}\\
Correlación RMS-DS $\uparrow$ & 0.173 & 0.278 & \textbf{0.647}\\
MAE RMS-DS [ns] $\downarrow$ & 43.21 & 27.32 & \textbf{16.50}\\
NMSE magnitud med. [dB] $\downarrow$ & $-3.17$ & $-3.92$ & \textbf{$-4.46$}\\
Coherencia compleja bruta $\uparrow$ & 0.156 & 0.155 & 0.145\\
Error angular bruto [$^\circ$] $\downarrow$ & 89.93 & 89.95 & 89.96\\
Error angular tras CPO [$^\circ$] & 81.42 & 81.39 & 81.58\\
\bottomrule
\end{tabular}
\end{adjustbox}
\end{table}

La conclusión de B0--B2 es una de las primeras observaciones fuertes de la
tesis:

La mejora de potencia, magnitud y dispersión de retardos no implica, en esta
comparación, una mejora de la concordancia angular bruta.

B2 obtiene los mejores agregados de potencia, dispersión de retardos y magnitud,
pero eso no resuelve la correspondencia compleja.

% ----------------------------------------------------------------------
\FloatBarrier
\section{S4.1: reproducción de calibración con incertidumbre de fase}
\label{sec:s41_full}
% ----------------------------------------------------------------------

Ruah et al.~\cite{Ruah2024Phase} identifican una debilidad del enfoque que
optimiza materiales directamente contra el canal complejo: un error geométrico
pequeño cambia la longitud de trayectoria y, por tanto, la fase
$2\pi d/\lambda$. El material puede terminar absorbiendo un error cuya causa es
geométrica o instrumental.

Se reprodujeron tres esquemas del trabajo:

\begin{description}
\item[PEOC:] Fase Error-Oblivious Calibration, que confía en la fase RT.
\item[UPEC:] Uniform Fase Error Calibration, que elimina prácticamente la
información de fase de trayectoria.
\item[PEAC:] Fase Error-Aware Calibration, con errores locales modelados
mediante una distribución circular y optimización variacional.
\end{description}

La distribución base puede escribirse

\begin{equation}
p(\epsilon|\mu,\kappa)
=
\frac{\exp[\kappa\cos(\epsilon-\mu)]}
{2\pi I_0(\kappa)}.
\end{equation}

Se evaluaron varias semillas y, entre otros, los regímenes $\kappa=0$ y
$\kappa=2.253$. El resultado experimental fue que UPEC resultó especialmente
estable para la permitividad, PEOC mostró sesgo sistemático cuando la fase
simulada era incorrecta y PEAC pudo recuperar soluciones muy precisas, pero con
sensibilidad a inicialización y cuencas de fallo. La conductividad fue
considerablemente menos identificable.

La lección no fue ``PEAC gana'', sino:

La expresividad del modelo de fase debe evaluarse junto con su estabilidad
de optimización e identificabilidad.

Esto motivó no aplicar PEAC inmediatamente a dc01, sino caracterizar primero
qué estructura de error de fase existe realmente.

% ----------------------------------------------------------------------
\FloatBarrier
\section{S4.2: descomposición de baja dimensión de la fase}
\label{sec:s42_full}
% ----------------------------------------------------------------------

Se define

\begin{equation}
R_{i,u,a,k}
=
H^{\rm real}_{i,u,a,k}
H^{{\rm RT}*}_{i,u,a,k}.
\end{equation}

La hipótesis inicial es

\begin{equation}
\angle R_{i,u,a,k}
\simeq
\phi_{0,i,u,a}
-
2\pi\nu_k\Delta\tau_{{\rm eff},i,u}
+
r_{i,u,a,k}
\pmod{2\pi}.
\label{eq:low_dim_phase}
\end{equation}

$\Delta\tau_{\rm eff}$ debe denominarse \emph{retardo espectral relativo efectivo}; no es correcto llamarlo STO físico. Las mediciones utilizadas ya
incluyen compensación de STO/CPO y el parámetro ajustado es relativo al modelo
RT concreto.

\FloatBarrier
\subsection{S4.2a: diferencias adyacentes}

A partir de

\begin{equation}
R_{k+1}R_k^*
\simeq
B_k e^{-j2\pi\Delta f\Delta\tau_{\rm eff}},
\end{equation}

se obtiene un estimador local de pendiente. Los resultados agregados fueron:

\begin{table}[!htbp]
\centering
\caption{S4.2a: error angular tras estimación local del retardo.}
\label{tab:s42a_results}
\begin{adjustbox}{max width=\textwidth}
\begin{tabular}{lrrr}
\toprule
Modelo & CPO [$^\circ$] & común adyacente [$^\circ$] & por enlace [$^\circ$]\\
\midrule
B0 & 81.42 & 66.05 & 66.23\\
B1 & 81.39 & 66.48 & 65.31\\
B2 & 81.58 & 54.99 & 60.71\\
\bottomrule
\end{tabular}
\end{adjustbox}
\end{table}

En este estimador B2 parecía particularmente favorable. Sin embargo, el
resultado dependía del procedimiento local elegido.

\FloatBarrier
\subsection{S4.2b: retardo oráculo de alineación óptima}

Se define

\begin{equation}
\widehat\tau_{i,u}
=
\arg\max_\tau
\sum_a
\left|
\sum_k
R_{i,u,a,k}
e^{j2\pi\nu_k\tau}
\right|.
\label{eq:s42b_objective}
\end{equation}

Se utilizó una rejilla espectral de 8192 puntos, un paso de búsqueda de 2.5 ns
y un intervalo de $\pm500$ ns. El relleno con ceros interpola el objetivo y
no aumenta la resolución física: $1/B=20$ ns para 50 MHz. El periodo de
ambigüedad es $1/\Delta f=20.48$ microsegundos con la rejilla nativa. La
selección intercalada modifica la ambigüedad y requiere conservar las frecuencias
originales y el intervalo común de búsqueda.

\begin{table}[!htbp]
\centering
\caption{S4.2b: alineación global oráculo. Los errores angulares se expresan en grados; las coherencias son adimensionales.}
\label{tab:s42b_results}
\begin{adjustbox}{max width=\textwidth}
\begin{tabular}{lrrrr}
\toprule
Modelo & fase CPO & fase oráculo & $\rho$ bruto & $\rho$ oráculo\\
\midrule
B0 & 81.42 & \textbf{31.63} & 0.156 & \textbf{0.745}\\
B1 & 81.39 & 37.15 & 0.155 & 0.696\\
B2 & 81.58 & 38.17 & 0.145 & 0.683\\
\bottomrule
\end{tabular}
\end{adjustbox}
\end{table}

Este experimento cambia la interpretación de B0--B2. B0, inferior en
reconstrucción macroscópica, resulta el más fácil de alinear en fase mediante
una transformación oráculo. Por ello:

La capacidad de alineación a posteriori es una propiedad condicionada a la
transformación y los datos disponibles; no sustituye la fidelidad predictiva del modelo.

\FloatBarrier
\subsection{S4.2c: acuerdo del retardo efectivo entre modelos}

Los retardos oráculo presentan:

\begin{table}[!htbp]
\centering
\caption{S4.2c: acuerdo entre modelos.}
\label{tab:s42c_results}
\begin{adjustbox}{max width=\textwidth}
\begin{tabular}{lrr}
\toprule
Par & Pearson & diferencia mediana\\
\midrule
B0--B1 & 0.886 & 0--2.5 ns\\
B0--B2 & 0.834 & 0--2.5 ns\\
B1--B2 & 0.819 & 0--2.5 ns\\
\bottomrule
\end{tabular}
\end{adjustbox}
\end{table}

Las tres estimaciones quedan dentro de 10 ns en 69.46\% de los casos. Esto
apoya la existencia de un componente compartido de baja dimensión, aunque con
una cola dependiente del modelo.

\FloatBarrier
\subsection{S4.2d: transferencia entre antenas}

Se estima el retardo usando 16 antenas y se aplica a las otras 16. El CPO de
las antenas evaluadas sigue siendo libre, de modo que esta prueba transfiere
retardo, no fase completa.

\begin{table}[!htbp]
\centering
\caption{S4.2d: transferencia entre antenas. Los errores angulares se expresan en grados; las coherencias son adimensionales.}
\label{tab:s42d_results}
\begin{adjustbox}{max width=\textwidth}
\begin{tabular}{lrrrrr}
\toprule
Modelo & CPO & transferido & oráculo & $\rho_{\rm tr}$ & $\rho_{\rm or}$\\
\midrule
B0 & 81.42 & 35.52 & 30.49 & 0.715 & 0.754\\
B1 & 81.39 & 42.19 & 35.34 & 0.652 & 0.708\\
B2 & 81.58 & 41.56 & 37.11 & 0.654 & 0.690\\
\bottomrule
\end{tabular}
\end{adjustbox}
\end{table}

La correlación del retardo ajuste/oráculo es 0.833, 0.762 y 0.789 para B0, B1 y
B2, respectivamente, con diferencia mediana de 2.5 ns.

\FloatBarrier
\subsection{S4.2e: transferencia estricta en frecuencia}

Esta prueba es más estricta porque tanto $\tau$ como $\phi_0$ se estiman
exclusivamente en las frecuencias de ajuste y se transfieren sin recalibración
a las frecuencias reservadas.

\begin{table}[!htbp]
\centering
\scriptsize
\caption{S4.2e: transferencia estricta en frecuencia. Los retardos se expresan en ns y los errores angulares en grados.}
\label{tab:s42e_results}
\begin{adjustbox}{max width=\textwidth}
\begin{tabular}{llrrrrrr}
\toprule
Partición & Modelo & fase bruta & fase transf. & fase oráculo &
$\rho_{\rm tr}$ & $\rho_{\rm or}$ & $|\Delta\tau|_{\rm med}$\\
\midrule
pares$\to$impares & B0 & 89.93 & 31.65 & 31.65 & 0.745 & 0.745 & 0\\
 & B1 & 89.96 & 37.17 & 37.17 & 0.695 & 0.695 & 0\\
 & B2 & 89.96 & 38.18 & 38.18 & 0.683 & 0.683 & 0\\
impares$\to$pares & B0 & 89.92 & 31.62 & 31.62 & 0.745 & 0.745 & 0\\
 & B1 & 89.95 & 37.14 & 37.13 & 0.696 & 0.696 & 0\\
 & B2 & 89.96 & 38.16 & 38.16 & 0.683 & 0.683 & 0\\
bajas$\to$altas & B0 & 89.92 & 44.80 & 25.49 & 0.785 & 0.791 & 2.5\\
 & B1 & 89.88 & 50.00 & 30.35 & 0.742 & 0.749 & 2.5\\
 & B2 & 89.88 & 49.64 & 32.28 & 0.723 & 0.733 & 2.5\\
altas$\to$bajas & B0 & 89.90 & 45.33 & 25.95 & 0.791 & 0.799 & 2.5\\
 & B1 & 90.00 & 50.00 & 30.60 & 0.752 & 0.759 & 2.5\\
 & B2 & 90.01 & 49.78 & 32.79 & 0.732 & 0.742 & 2.5\\
\bottomrule
\end{tabular}
\end{adjustbox}
\end{table}

La interpolación intercalada alcanza prácticamente el oráculo. La extrapolación
de media banda es más difícil pese a que la diferencia de retardo mediana es
sólo 2.5 ns. A 25 MHz:

\begin{equation}
2\pi(25\times10^6)(2.5\times10^{-9})
\approx 22.5^\circ.
\end{equation}

Por tanto, errores de pocos nanosegundos tienen una consecuencia angular
macroscópica.

% ----------------------------------------------------------------------
\FloatBarrier
\section{S4.3: estructura del residual sensible a la fase}
\label{sec:s43_full}
% ----------------------------------------------------------------------

\FloatBarrier
\subsection{S4.3a: diagnóstico espectral}

Después de eliminar el mejor retardo común y CPO se analiza el residual
$r_\phi$. Sobre B0, 4900 posiciones y 313\,600 enlaces de antena:

\begin{table}[!htbp]
\centering
\caption{S4.3a: estadísticos del residual de fase.}
\label{tab:s43a_stats}
\begin{adjustbox}{max width=\textwidth}
\begin{tabular}{lr}
\toprule
Indicador & Valor\\
\midrule
Mediana del MAE angular por enlace & 31.634$^\circ$\\
Q25 & 18.643$^\circ$\\
Q75 & 51.994$^\circ$\\
Q90 & 71.166$^\circ$\\
Media & 37.087$^\circ$\\
Resultante circular mediana & 0.7860\\
Varianza circular mediana & 0.2140\\
MAE de incrementos de fase & 2.407$^\circ$\\
Concentración de incrementos & 0.997427\\
Coherencia residual & 0.745215\\
\bottomrule
\end{tabular}
\end{adjustbox}
\end{table}

La resultante de incrementos espectrales por separación fue:

\begin{table}[!htbp]
\centering
\caption{S4.3a: resultante circular de incrementos por separación espectral.}
\label{tab:s43a_lag}
\begin{adjustbox}{max width=\textwidth}
\begin{tabular}{cr}
\toprule
Separación [subportadoras] & $C_\phi(L)$ mediano\\
\midrule
1 & 0.997427\\
2 & 0.995886\\
4 & 0.991220\\
8 & 0.976974\\
16 & 0.938767\\
32 & 0.870279\\
64 & 0.838135\\
128 & 0.784616\\
256 & 0.738365\\
512 & 0.761427\\
\bottomrule
\end{tabular}
\end{adjustbox}
\end{table}

Los incrementos presentan elevada concentración a pequeña separación y una
estructura espectral que no queda descrita por el error angular marginal.
Estos resultados motivan contrastar modelos de dependencia; una resultante
no nula a separaciones grandes no rechaza por sí sola independencia, según
la distinción de la \cref{sec:circular_foundations}.

\FloatBarrier
\subsection{S4.3b: complejidad geométrica}

Las medianas de los descriptores de trayectorias fueron aproximadamente:
$n_{\rm spec}=168$, $n_{\rm diff}=6$, $n_{\rm active}=174$,
$n_{\rm scat}=304$ y extensión de retardo activo 687.36 ns. Los conteos brutos
presentan asociaciones débiles; el extensión de retardo, por ejemplo, tiene
Spearman $\simeq-0.015$ con el residual de fase.

Esto descarta una hipótesis demasiado simple: el residual no queda explicado
por ``tener más trayectorias''.

\FloatBarrier
\subsection{S4.3c: multitrayecto efectivo ponderado por energía}

Se construyen descriptores como

\begin{align}
D_1 &= \max_\ell p_\ell,\\
D_3 &= \sum_{\ell\in{\rm top3}} p_\ell,\\
D_5 &= \sum_{\ell\in{\rm top5}} p_\ell,\\
N_{\rm eff} &= \frac{1}{\sum_\ell p_\ell^2},\\
\mathcal H_p &=
-\frac{\sum_\ell p_\ell\log p_\ell}{\log L}.
\end{align}

\begin{table}[!htbp]
\centering
\caption{S4.3c: complejidad multitrayecto efectiva.}
\label{tab:s43c_stats}
\begin{adjustbox}{max width=\textwidth}
\begin{tabular}{lr}
\toprule
Descriptor & Mediana\\
\midrule
$n_{\rm energy>0}$ & 174\\
Fracción dominante $D_1$ & 0.2660\\
Fracción de los tres caminos dominantes & 0.5692\\
Fracción de los cinco caminos dominantes & 0.7099\\
Número efectivo de trayectorias $N_{\rm eff}$ & 7.1667\\
Entropía normalizada & 0.5183\\
dispersión RMS de retardos ponderado & 15.2584 ns\\
$n_{10{\rm dB}}$ & 9\\
$n_{20{\rm dB}}$ & 17\\
\bottomrule
\end{tabular}
\end{adjustbox}
\end{table}

A nivel de posición, entre otros resultados:

\begin{align}
\rho_S(\mathcal H_p,E_\phi) &= +0.491,\\
\rho_S(\mathcal H_p,C_\phi(64)) &= -0.679,\\
\rho_S(D_1,C_\phi(64)) &= +0.670,\\
\rho_S(\tau_{\rm RMS},C_\phi(16)) &= -0.520,\\
\rho_S(\tau_{\rm RMS},C_\phi(64)) &= -0.437.
\end{align}

Los descriptores de distribución energética muestran asociaciones de mayor
magnitud con la estructura residual que los conteos geométricos considerados.
Esta comparación es descriptiva y no atribuye causalidad a un descriptor aislado.

\FloatBarrier
\subsection{S4.3d: efectos entre posiciones y dentro de posición}

Se separa la variación entre posiciones de la variación entre enlaces dentro
de una misma posición. Aparecen inversiones de signo:

\begin{table}[!htbp]
\centering
\scriptsize
\caption{S4.3d: ejemplos de efectos multiescala.}
\label{tab:s43d}
\begin{adjustbox}{max width=\textwidth}
\begin{tabular}{lrrr}
\toprule
Relación & Entre posiciones & Dentro de posición & Enlaces agrupados\\
\midrule
$N_{\rm eff}\leftrightarrow E_\phi$ & -0.461 & +0.247 & +0.265\\
$N_{\rm eff}\leftrightarrow C_\phi(64)$ & +0.492 & -0.361 & -0.474\\
$\mathcal H_p\leftrightarrow E_\phi$ & +0.521 & +0.249 & +0.372\\
$\mathcal H_p\leftrightarrow C_\phi(64)$ & -0.700 & -0.341 & -0.425\\
$\tau_{\rm RMS}\leftrightarrow C_\phi(16)$ & -0.589 & -0.130 & -0.312\\
$\tau_{\rm RMS}\leftrightarrow C_\phi(64)$ & -0.575 & -0.182 & -0.278\\
\bottomrule
\end{tabular}
\end{adjustbox}
\end{table}

El caso de $N_{\rm eff}$ tiene una interpretación tipo Simpson: comparar dos
entornos y comparar dos enlaces del mismo entorno no representa la misma
pregunta física.

\FloatBarrier
\subsection{S4.3e: bootstrap agrupado por posición}

Se ejecutaron 2000 remuestreos a nivel de receptor, evitando tratar antenas
correlacionadas como observaciones independientes. Entre los intervalos de
confianza obtenidos:

\begin{table}[!htbp]
\centering
\scriptsize
\caption{S4.3e: asociaciones con bootstrap por posición.}
\label{tab:s43e}
\begin{adjustbox}{max width=\textwidth}
\begin{tabular}{lccc}
\toprule
Relación & Entre posiciones [95\% CI] & Dentro de posición [95\% CI] & Fracción reportada\\
\midrule
$\mathcal H_p\leftrightarrow E_\phi$ &
+0.521 [0.499,0.543] & +0.304 [0.297,0.314] & 82.7\%\\
$\mathcal H_p\leftrightarrow C_\phi(64)$ &
-0.700 [-0.717,-0.681] & -0.557 [-0.565,-0.547] & 83.0\%\\
$D_1\leftrightarrow C_\phi(64)$ &
+0.754 [0.739,0.768] & +0.540 [0.530,0.550] & 84.2\%\\
$N_{\rm eff}\leftrightarrow E_\phi$ &
-0.461 [-0.485,-0.439] & +0.295 [0.285,0.303] & 81.6\%\\
$N_{\rm eff}\leftrightarrow C_\phi(64)$ &
+0.492 [0.475,0.508] & -0.552 [-0.560,-0.543] & 83.2\%\\
\bottomrule
\end{tabular}
\end{adjustbox}
\end{table}

Las reducciones de S4.3d y S4.3e se conservan separadas: sus coeficientes
intrapositivos no son numéricamente intercambiables. El análisis dentro de posición
puede calcularse sobre residuos centrados o resumir coeficientes por posición,
como se distingue en la \cref{sec:multilevel_foundations}. La configuración de
cada script debe fijar cuál se utiliza. La última columna reproduce la fracción
denominada ``mismo signo'' en el registro, cuyo referente de comparación debe
especificarse antes de interpretarla; no representa concordancia entre las dos
columnas anteriores, que presentan signos opuestos para $N_{\rm eff}$.

S4.3 documenta asociaciones entre el residual de fase y la distribución de
energía y retardo de las trayectorias. Los intervalos del remuestreo por
posición cuantifican estabilidad bajo ese esquema, sin eliminar la posible
dependencia espacial entre posiciones ni identificar una causa única.

% ----------------------------------------------------------------------
\FloatBarrier
\section{S4.4: predictibilidad del retardo efectivo}
\label{sec:s44_full}
% ----------------------------------------------------------------------

\FloatBarrier
\subsection{S4.4a: conjunto de datos de predicción}

La variable objetivo por posición y Tx es el retardo de alineación óptima de S4.2b:
\begin{equation}
y_{i,u}=\Delta\tau_{{\rm eff},i,u}^{\star}.
\end{equation}

Los descriptores provienen únicamente del RT. Se calculan 48 variables a partir de
estadísticos (media, mediana, desviación e IQR) de concentración de potencia,
número efectivo de trayectorias, entropía, dispersión RMS de retardos, conteos dentro de umbrales de dB y
potencia total sobre las 32 antenas.

\begin{table}[!htbp]
\centering
\caption{S4.4a: conjunto de datos para predicción del retardo efectivo.}
\label{tab:s44a_split}
\begin{adjustbox}{max width=\textwidth}
\begin{tabular}{lrrr}
\toprule
Partición & posiciones & filas posición--Tx & descriptores RT\\
\midrule
Entrenamiento & 5000 & 10000 & 48\\
Validación & 100 & 200 & 48\\
Prueba & 4900 & 9800 & 48\\
\bottomrule
\end{tabular}
\end{adjustbox}
\end{table}

La comprobación de correspondencia interna con S4.3a produjo error máximo de posición cero y
diferencia de variable objetivo mediana/máxima cero.

\FloatBarrier
\subsection{S4.4b: métodos de referencia}

\begin{table}[!htbp]
\centering
\scriptsize
\caption{S4.4b: métodos de referencia de predicción en la partición original. Los errores de retardo se expresan en ns.}
\label{tab:s44b_results}
\begin{adjustbox}{max width=\textwidth}
\begin{tabular}{lrrrrrrrr}
\toprule
Modelo & MAE & MedAE & RMSE & $R^2$ & Pearson & Spearman &
$\le5$ns & $\le10$ns\\
\midrule
Constante & 47.38 & 20.00 & 73.43 & -0.118 & -- & -- & 0.186 & 0.288\\
Ridge pos & 49.98 & 42.67 & 62.40 & 0.193 & 0.440 & 0.397 & 0.050 & 0.109\\
Ridge RT & 19.57 & 11.53 & 34.30 & 0.756 & 0.870 & 0.832 & 0.247 & 0.444\\
Ridge pos+RT & 19.30 & 11.52 & 34.00 & 0.760 & 0.872 & 0.831 & 0.245 & 0.445\\
RF pos+RT & \textbf{10.89} & \textbf{4.74} & \textbf{25.37} &
\textbf{0.867} & \textbf{0.931} & \textbf{0.928} & \textbf{0.515} &
\textbf{0.724}\\
\bottomrule
\end{tabular}
\end{adjustbox}
\end{table}

\FloatBarrier
\subsection{S4.4c: ablación no lineal}

\begin{table}[!htbp]
\centering
\caption{S4.4c: Random Forest con distintos grupos de descriptores. Los errores de retardo se expresan en ns.}
\label{tab:s44c_results}
\begin{adjustbox}{max width=\textwidth}
\begin{tabular}{lrrrrrr}
\toprule
Modelo & MAE & MedAE & Q90 & $R^2$ & Pearson & Spearman\\
\midrule
RF posición & 11.15 & 5.09 & 24.66 & 0.867 & 0.931 & 0.929\\
RF RT & 11.24 & 5.20 & 24.79 & 0.864 & 0.930 & 0.925\\
RF pos+RT & \textbf{10.89} & \textbf{4.74} & \textbf{24.12} &
0.867 & 0.931 & 0.928\\
\bottomrule
\end{tabular}
\end{adjustbox}
\end{table}

En RT, las variables de mayor importancia predictiva son $n_{20\,\mathrm{dB}}$, fracción de potencia dominante, número efectivo de trayectorias y fracciones de los tres y cinco caminos dominantes. Este resultado es
coherente con S4.3. Sin embargo, posición sola obtiene prácticamente la precisión semejante, lo que motiva una auditoría dla partición.

\FloatBarrier
\subsection{S4.4d: auditoría de proximidad espacial}

La distancia del prueba a su vecino entrenamiento más cercano es:

\begin{align}
{\rm mediana} &= 0.048289~{\rm m},\\
Q_{25}&=0.016304~{\rm m},\\
Q_{75}&=0.107514~{\rm m},\\
Q_{90}&=0.184428~{\rm m},\\
Q_{95}&=0.252456~{\rm m},\\
{\rm max}&=0.494789~{\rm m}.
\end{align}

Un método de referencia que simplemente copia el variable objetivo del vecino entrenamiento más cercano
alcanza MedAE 2.5 ns; Spearman 0.899 en Tx1 y 0.931 en Tx2. Además, el error
del RF aumenta con distancia. Para RF pos+RT, la MedAE pasa de 2.70 ns en el
primer cuartil a 6.90 ns en el cuarto.

Por ello, la partición original mide principalmente interpolación espacial densa y
no debe presentarse como una prueba fuerte de generalización espacial.

\FloatBarrier
\subsection{S4.4e: bloques espaciales con margen de exclusión}

El eje $y$ tiene el mayor extensión, aproximadamente 39.99 m. Se divide en cinco
bloques contiguos de 2000 posiciones y se elimina del entrenamiento cualquier
posición a menos de 0.5 m del bloque prueba.

\begin{table}[!htbp]
\centering
\caption{S4.4e: geometría de las cinco particiones espaciales.}
\label{tab:s44e_blocks}
\begin{adjustbox}{max width=\textwidth}
\begin{tabular}{crrrr}
\toprule
Bloque & Posiciones de prueba & Entrenamiento conservado & removidos & NN mediana [m]\\
\midrule
0 & 2000 & 7917 & 83 & 7.331\\
1 & 2000 & 7749 & 251 & 1.677\\
2 & 2000 & 7603 & 397 & 2.131\\
3 & 2000 & 7605 & 395 & 3.314\\
4 & 2000 & 7916 & 84 & 2.520\\
\bottomrule
\end{tabular}
\end{adjustbox}
\end{table}

\FloatBarrier
\subsection{S4.4f: generalización espacial bloqueada}

\begin{table}[!htbp]
\centering
\scriptsize
\caption{S4.4f: MAE/MedAE por bloque para el retardo efectivo. Los errores de retardo se expresan en ns.}
\label{tab:s44f_results}
\begin{adjustbox}{max width=\textwidth}
\begin{tabular}{clrrrr}
\toprule
Bloque & Modelo & MAE & MedAE & $R^2$ & Spearman\\
\midrule
0 & RF posición & 48.66 & 28.73 & 0.146 & 0.692\\
0 & RF RT & 21.76 & 16.02 & 0.782 & 0.874\\
0 & RF pos+RT & \textbf{20.66} & \textbf{13.07} & 0.785 & 0.890\\
1 & RF posición & \textbf{13.77} & \textbf{9.13} & 0.224 & 0.710\\
1 & RF RT & 15.07 & 10.65 & 0.185 & 0.585\\
1 & RF pos+RT & 14.43 & 10.04 & 0.235 & 0.630\\
2 & RF posición & 17.44 & 10.31 & -0.046 & 0.480\\
2 & RF RT & 15.83 & 10.63 & 0.346 & 0.619\\
2 & RF pos+RT & \textbf{15.03} & \textbf{9.09} & 0.363 & 0.638\\
3 & RF posición & 38.74 & 22.33 & 0.463 & 0.649\\
3 & RF RT & 27.35 & 17.52 & 0.716 & 0.832\\
3 & RF pos+RT & \textbf{25.81} & \textbf{16.26} & 0.738 & 0.834\\
4 & RF posición & 19.19 & 13.32 & 0.885 & 0.820\\
4 & RF RT & \textbf{16.45} & 13.93 & 0.930 & 0.779\\
4 & RF pos+RT & 17.50 & 15.48 & 0.927 & 0.784\\
\bottomrule
\end{tabular}
\end{adjustbox}
\end{table}

La siguiente tabla reporta la media de los indicadores de las cinco particiones;
no corresponde a recalcular una mediana global de todas las filas:

\begin{table}[!htbp]
\centering
\caption{S4.4f: agregado de cinco particiones espaciales. Los errores de retardo se expresan en ns.}
\label{tab:s44f_aggregate}
\begin{adjustbox}{max width=\textwidth}
\begin{tabular}{lrrrr}
\toprule
Modelo & MAE medio & MedAE media & $R^2$ medio & Spearman medio\\
\midrule
RF posición & 27.56 & 16.76 & 0.334 & 0.670\\
RF RT & 19.29 & 13.75 & 0.592 & 0.738\\
RF pos+RT & \textbf{18.69} & \textbf{12.79} & \textbf{0.610} &
\textbf{0.755}\\
\bottomrule
\end{tabular}
\end{adjustbox}
\end{table}

RT reduce aproximadamente 30\% el MAE medio frente a sólo posición. La
ganancia es particularmente notable en bloque 0, donde la separación mediana
al entrenamiento es 7.33 m. Este resultado aporta evidencia, en los bloques de dc01 evaluados, de que los
descriptores físicos de $\mathcal P_{\rm RT}$ contienen información
transferible sobre $\Delta\tau_{\rm eff}$ más allá de la proximidad espacial.

% ----------------------------------------------------------------------
\FloatBarrier
\section{S4.5: efecto del retardo predicho sobre el canal}
\label{sec:s45_full}
% ----------------------------------------------------------------------

\FloatBarrier
\subsection{S4.5a: partición original}

El predictor de $\tau$ se aplica sobre el CFR y se recalculan MAE angular y
coherencia. El CPO por enlace todavía se estima desde la medición prueba, por lo
que el experimento aísla la transferibilidad del retardo, no una corrección
completamente sin acceso a la medición evaluada.

\begin{table}[!htbp]
\centering
\caption{S4.5a: resultados sobre el canal corregido en 4900 posiciones. Los errores angulares se expresan en grados; las coherencias son adimensionales.}
\label{tab:s45a_results}
\begin{adjustbox}{max width=\textwidth}
\begin{tabular}{lrrrr}
\toprule
Método & fase med. & fase Q75 & $\rho$ med. & $\rho$ Q25\\
\midrule
RAW & 89.93 & 95.83 & 0.156 & 0.047\\
CPO & 81.42 & 87.96 & 0.156 & 0.047\\
Pred. $\tau$ + CPO & \textbf{45.59} & 70.07 & \textbf{0.634} & 0.361\\
Oráculo $\tau$ + CPO & 31.63 & 51.99 & 0.745 & 0.559\\
\bottomrule
\end{tabular}
\end{adjustbox}
\end{table}

El error de $\tau$ es MedAE 4.7399 ns, Q75 11.1313 ns, Q90 24.1221 ns y
MAE 10.8904 ns. El predictor mejora el error angular respecto de CPO en 81.50\% de
los enlaces y la coherencia en 83.47\%.

Definiendo

\begin{equation}
\eta_\phi
=
\frac{
E_{\phi,\rm CPO}-E_{\phi,\rm pred}
}{
E_{\phi,\rm CPO}-E_{\phi,\mathrm{or}}
},
\end{equation}

se obtiene aproximadamente $\eta_\phi=0.72$. Es una razón de mejoras de
agregados, no el promedio de fracciones por enlace. La razón sólo se interpreta
cuando el denominador es positivo y suficientemente distinto de cero; no está
forzada al intervalo $[0,1]$ para cualquier predictor o métrica. Para coherencia, la fracción
análoga es aproximadamente 0.81. El predictor recupera, por tanto, una parte
sustancial de la corrección oráculo.

\FloatBarrier
\subsection{S4.5b: generalización por bloques espaciales sobre el canal corregido}

S4.5b utiliza las predicciones de S4.4f, sin reentrenar durante la evaluación
del canal. La trazabilidad fue comprobada explícitamente:
los errores máximos de posición entre el manifiesto y los TFRecords fueron
$3.55\times10^{-15}$ m en entrenamiento, $3.55\times10^{-15}$ m en validación y
$2.22\times10^{-16}$ m en prueba.

\begin{table}[!htbp]
\centering
\scriptsize
\caption{S4.5b: resultados sobre el canal corregido por partición espacial. Los retardos se expresan en ns y los errores angulares en grados.}
\label{tab:s45b_fold_results}
\begin{adjustbox}{max width=\textwidth}
\begin{tabular}{clrrrrrr}
\toprule
Bloque & Predictor & $\tau$ MedAE & fase med. & $\rho$ med. &
$\eta_\phi$ & $\eta_\rho$ & $P(E_\phi<CPO)$\\
\midrule
0 & posición & 28.73 & 83.28 & 0.121 & 0.017 & 0.037 & 0.544\\
0 & RT & 16.02 & 72.07 & 0.350 & 0.224 & 0.381 & 0.692\\
0 & pos+RT & \textbf{13.07} & \textbf{66.49} & \textbf{0.427} & 0.328 & 0.496 & 0.726\\
1 & posición & \textbf{9.13} & \textbf{55.70} & \textbf{0.535} & 0.440 & 0.549 & 0.665\\
1 & RT & 10.65 & 58.74 & 0.507 & 0.361 & 0.481 & 0.641\\
1 & pos+RT & 10.04 & 57.57 & 0.519 & 0.391 & 0.510 & 0.649\\
2 & posición & 10.31 & 60.91 & 0.482 & 0.371 & 0.504 & 0.655\\
2 & RT & 10.63 & 60.09 & 0.486 & 0.394 & 0.514 & 0.647\\
2 & pos+RT & \textbf{9.09} & \textbf{56.83} & \textbf{0.525} & 0.481 & 0.604 & 0.686\\
3 & posición & 22.33 & 76.31 & 0.261 & 0.155 & 0.261 & 0.703\\
3 & RT & 17.52 & 73.62 & 0.289 & 0.204 & 0.303 & 0.751\\
3 & pos+RT & \textbf{16.26} & \textbf{72.18} & \textbf{0.317} & 0.229 & 0.344 & 0.770\\
4 & posición & 13.32 & 67.27 & 0.427 & 0.321 & 0.514 & 0.799\\
4 & RT & 13.93 & \textbf{65.39} & \textbf{0.444} & 0.352 & 0.537 & 0.840\\
4 & pos+RT & 15.48 & 68.24 & 0.407 & 0.305 & 0.486 & 0.821\\
\bottomrule
\end{tabular}
\end{adjustbox}
\end{table}

\begin{table}[!htbp]
\centering
\caption{S4.5b: agregado sobre cinco particiones espaciales. Los retardos se expresan en ns y los errores angulares en grados.}
\label{tab:s45b_aggregate}
\begin{adjustbox}{max width=\textwidth}
\begin{tabular}{lrrrrrr}
\toprule
Predictor & $\tau$ MedAE & $\tau$ MAE & fase med. & $\rho$ med. &
$\eta_\phi$ & $\eta_\rho$\\
\midrule
Posición & 14.83 & 27.56 & 69.72 & 0.377 & 0.235 & 0.373\\
RT & 13.32 & 19.29 & 65.97 & 0.427 & 0.310 & 0.458\\
Pos+RT & \textbf{12.30} & \textbf{18.69} & \textbf{64.08} &
\textbf{0.448} & \textbf{0.348} & \textbf{0.493}\\
\bottomrule
\end{tabular}
\end{adjustbox}
\end{table}

En S4.5b los indicadores se calculan sobre las predicciones agrupadas de los
bloques. Por ello la mediana de retardo 12.30 ns no debe confundirse con
la media de medianas 12.79 ns de S4.4f.

El método de referencia CPO agregado es 81.49 grados y coherencia 0.157; el oráculo alcanza
31.46 grados y 0.746. Por tanto, en regiones reservadas el
predictor combinado recupera 34.8\% del beneficio angular oráculo y 49.3\% del
beneficio de coherencia.

La comparación entre S4.5a y S4.5b cuantifica la pérdida de transferibilidad:

\begin{table}[!htbp]
\centering
\caption{Comparación entre interpolación densa y extrapolación espacial.}
\label{tab:dense_vs_blocked}
\begin{adjustbox}{max width=\textwidth}
\begin{tabular}{lrr}
\toprule
Indicador & Partición densa & Bloques espaciales\\
\midrule
$\tau$ MedAE [ns] & 4.74 & 12.30\\
Fase predicha [deg] & 45.59 & 64.08\\
Coherencia predicha & 0.634 & 0.448\\
$\eta_\phi$ & 0.720 & 0.348\\
$\eta_\rho$ & 0.812 & 0.493\\
\bottomrule
\end{tabular}
\end{adjustbox}
\end{table}

La corrección de retardo es, por tanto, \emph{transferible pero localmente
condicionada}. La dependencia regional motiva estudiar el soporte espacial de calibración
de $Z_\phi^{\rm pred}$. Los cinco bloques no identifican por sí solos una
longitud universal de generalización.

% ----------------------------------------------------------------------
\FloatBarrier
\section{S4.6: naturaleza del desfase constante residual}
\label{sec:s46_full}
% ----------------------------------------------------------------------

Los experimentos S4.2--S4.5 mostraron que una parte importante de la discrepancia
compleja puede representarse mediante una pendiente espectral efectiva
$\Delta\tau_{\rm eff}$. Sin embargo, aun después de compensar esa pendiente se
requiere un término de fase constante por enlace para alcanzar la referencia
oráculo utilizada en la descomposición. S4.6 estudia la naturaleza de ese segundo
grado de libertad antes de intentar incorporarlo a una representación latente o
predecirlo mediante aprendizaje automático.

Las etiquetas \texttt{Tx0} y \texttt{Tx1} se conservan en las tablas porque así
figuran en el registro S4.6; aquí indexan los dos arreglos de 32 antenas del
tensor, no dos transmisores móviles adicionales. El bloque caracteriza la
referencia B0 y no extiende automáticamente sus resultados a B1/B2.


Para cada posición $i$, arreglo $u$ y antena $a$ se define
\begin{equation}
\phi_{0,i,u,a}^{\star}
=
\angle\left[
\sum_k
H^{\rm real}_{i,u,a,k}
H^{{\rm RT}*}_{i,u,a,k}
\exp\left(j2\pi\nu_k\tau_{i,u}^{\star}\right)
\right],
\label{eq:s46_phi0_def}
\end{equation}
donde $\tau_{i,u}^{\star}$ es el retardo óptimo según el criterio de S4.2b. El
superíndice $\star$ recuerda que $\phi_0$ es un descriptor posterior de la
medición y no una cantidad disponible a priori para una consulta nueva.
Se utiliza $\nu_k=f_k-f_c$, con la misma referencia que en S4.2.
Cambiar el origen frecuencial cambia también el intercepto; una comparación de
$\phi_0^\star$ o $\psi$ exige conservar esa convención y las máscaras de validez.

Es importante evitar una identificación causal prematura. DICHASUS dcxx
proporciona estimaciones de desfases de fase y tiempo para compensar el canal de
referencia y advierte que dichas estimaciones son \emph{best effort}, que son
más fiables dentro de un mismo arreglo que entre arreglos distantes y que el
canal de referencia puede fluctuar temporalmente \cite{DICHASUSdcxxWeb}. Por
otro lado, Ruah et al.\ muestran que errores geométricos pequeños pueden inducir
errores de fase importantes en los caminos simulados \cite{Ruah2024Phase}. En
consecuencia, $\phi_0^\star$ se interpreta en esta etapa como \emph{desfase constante residual efectivo entre la medición y el trazado}; no como el CPO
físico original del receptor.

\FloatBarrier
\subsection{S4.6a: construcción del conjunto completo de fases residuales}

La prueba preliminar de integridad reportó 1536 observaciones y una resultante
circular global $R=0.05123$. El registro suministrado contiene una discrepancia en la descripción
del tamaño: 24 posiciones en cada una de tres particiones, con dos arreglos y
32 antenas, producirían $24\times3\times2\times32=4608$ observaciones;
1536 corresponden a 24 posiciones en total. Se conserva el conteo reportado,
sin asignar retrospectivamente una distribución por partición no corroborada.
Esta prueba comprueba el procesamiento a pequeña escala y no sustenta las
conclusiones estadísticas del conjunto completo.

Posteriormente se ejecutó el cálculo completo sobre las 10\,000 posiciones del banco. El conjunto final contiene
\begin{equation}
10000\times2\times32=\boxed{640000}
\end{equation}
observaciones de $\phi_0^\star$. Las cuentas por partición son 320\,000 para entrenamiento,
6400 para validación y 313\,600 para prueba. La resultante circular global cae a
\begin{equation}
R_{\rm global}=0.001088,
\end{equation}
con rango prácticamente completo $[-\pi,\pi]$.

\begin{table}[!htbp]
\centering
\caption{S4.6a: construcción del conjunto de datos de desfase constante residual.}
\label{tab:s46a_dataset}
\begin{adjustbox}{max width=\textwidth}
\begin{tabular}{lrr}
\toprule
Partición & Filas & Interpretación\\
\midrule
Entrenamiento & 320\,000 & 5000 posiciones $\times2\times32$\\
Validación & 6400 & 100 posiciones $\times2\times32$\\
Prueba & 313\,600 & 4900 posiciones $\times2\times32$\\
\midrule
Total & 640\,000 & 10\,000 posiciones $\times2\times32$\\
\bottomrule
\end{tabular}
\end{adjustbox}
\end{table}

La concentración global casi nula es incompatible con una fase universal
fuertemente concentrada en estas observaciones, pero no demuestra uniformidad ni
independencia circular: diferentes antenas o posiciones
podrían poseer desfases concentrados alrededor de medias distintas que se cancelan
al agregarse. Por ello S4.6b descompone explícitamente las escalas de variación.

\FloatBarrier
\subsection{S4.6b: descomposición circular por arreglo, antena y posición}

Se probó inicialmente la jerarquía descriptiva
\begin{equation}
\phi_{0,i,u,a}
\simeq
\beta_{u,a}+\psi_{i,u}+\epsilon_{i,u,a}
\pmod{2\pi},
\label{eq:s46_hierarchy}
\end{equation}
donde $\beta_{u,a}$ sería un sesgo fijo por antena, $\psi_{i,u}$ una rotación
común a las 32 antenas de un mismo arreglo en una posición y
$\epsilon_{i,u,a}$ un residual específico del enlace.
La suma admite la libertad de referencia
$(\beta_{u,a},\psi_{i,u})\mapsto(\beta_{u,a}+c_u,\psi_{i,u}-c_u)$.
El procedimiento descrito a continuación fija una convención mediante medias
marginales por antena; no equivale a un ajuste conjunto identificable de
componentes físicas. Cuando la resultante marginal es casi nula, su dirección
media es inestable y debe tratarse como parte del análisis de sensibilidad.


En la prueba preliminar, las resultantes medianas por antena parecían moderados
($R\approx0.198$ para Tx0 y $0.150$ para Tx1), pero esta impresión desapareció
con el conjunto completo. Sobre 10\,000 posiciones, la concentración por arreglo es
mínima: $R=0.00594$ para Tx0 y $0.00455$ para Tx1. Más importante, la resultante
mediana por antena es sólo $0.01024$ para Tx0 y $0.00630$ para Tx1; incluso el
máximo entre las 64 combinaciones arreglo--antena permanece por debajo de $0.023$.

\begin{table}[!htbp]
\centering
\small
\caption{S4.6b: estabilidad circular por antena en el conjunto completo.}
\label{tab:s46b_ant}
\begin{adjustbox}{max width=\textwidth}
\begin{tabular}{lrrrrr}
\toprule
Arreglo & mediana $R$ & Q25 & Q75 & Q90 & máximo\\
\midrule
Tx0 & 0.01024 & 0.00631 & 0.01405 & 0.01789 & 0.02029\\
Tx1 & 0.00630 & 0.00482 & 0.00887 & 0.01422 & 0.02243\\
\bottomrule
\end{tabular}
\end{adjustbox}
\end{table}

La dirección marginal por antena se expresa como
\begin{equation}
\beta_{u,a}=\angle\left[\frac{1}{N}\sum_{i=1}^{N}
e^{j\phi_{0,i,u,a}^{\star}}\right],\qquad N=10\,000.
\end{equation}
Después de retirar descriptivamente esa media circular de cada antena,
el error absoluto mediano permanece en
\begin{equation}
\operatorname{Med}|\wrap(\phi_0-\beta)|=88.99^{\circ},
\end{equation}
con una resultante global $0.00905$. Por tanto, una tabla estática de desfases por
antena no constituye una explicación útil del fenómeno.

La situación cambia al analizar simultáneamente las 32 antenas de una misma
posición y arreglo. Una vez retirado $\beta$, se define
\begin{equation}
\psi_{i,u}
=
\angle\left[
\frac{1}{32}\sum_{a=1}^{32}
\exp\left(j[\phi_{0,i,u,a}-\beta_{u,a}]\right)
\right].
\label{eq:s46_psi}
\end{equation}
Se denota por $R_{{\rm within},i,u}$ el módulo del promedio complejo
que aparece dentro del argumento de esta expresión.
La resultante circular dentro de posición presenta mediana $0.4216$, Q25
$0.3013$, Q75 $0.5499$ y Q90 $0.6897$. Es decir, aunque la orientación absoluta
presenta concentración global muy baja al recorrer el conjunto de datos, las antenas de una misma
captura muestran una dirección circular común apreciable.

Al sustraer esa dirección común, el residual absoluto mediano disminuye de
$88.99^{\circ}$ a $45.97^{\circ}$ y la resultante del residual centrado aumenta a
$0.4291$. Tx1 presenta una estructura común algo más marcada que Tx0: el error
final es $42.49^{\circ}$ y $R=0.507$ frente a $50.42^{\circ}$ y $R=0.351$ para
Tx0.

\begin{table}[!htbp]
\centering
\caption{S4.6b: descomposición descriptiva completa de $\phi_0^\star$.}
\label{tab:s46b_decomposition}
\begin{adjustbox}{max width=\textwidth}
\begin{tabular}{lr}
\toprule
Indicador & Resultado\\
\midrule
Resultante global bruta & 0.00109\\
Error mediano tras $\beta_{u,a}$ & $88.99^{\circ}$\\
Resultante mediana dentro de posición & 0.4216\\
Q25 / Q75 de la resultante dentro de posición & 0.3013 / 0.5499\\
Q90 de la resultante dentro de posición & 0.6897\\
Error mediano tras $\beta+\psi$ & $45.97^{\circ}$\\
Resultante residual tras $\beta+\psi$ & 0.4291\\
Tx0 error final / resultante & $50.42^{\circ}$ / 0.351\\
Tx1 error final / resultante & $42.49^{\circ}$ / 0.507\\
\bottomrule
\end{tabular}
\end{adjustbox}
\end{table}

Las tres particiones reproducen prácticamente el mismo patrón: resultantes
dentro de posición medianas de 0.421--0.432 y errores finales medianos de
44.5--46.3 grados. La estabilidad entre particiones refuerza que no se trata de una
peculiaridad del subconjunto utilizado.

Este resultado modifica la hipótesis original: $\beta_{u,a}$ no es un término
dominante, mientras que $\psi_{i,u}$ es una dirección común estimada a posteriori dentro de
cada arreglo. Sin embargo, como $\psi$ se estima usando las propias 32 antenas de
la posición evaluada, S4.6b constituye una evaluación descriptiva con acceso oráculo y no demuestra
transferibilidad.
Además, esa dirección minimiza la pérdida de coseno sobre las propias antenas
utilizadas para estimarla; la mejora del ajuste no es una validación fuera de muestra. Para cuantificar cuánto excede esa reducción el ajuste a
fluctuaciones finitas se requiere un control que permute posiciones de cada
antena, conserve sus marginales y repita todo el centrado. La baja concentración
marginal tampoco descarta sesgos relativos estables ocultos por rotaciones
comunes variables: esa hipótesis se evalúa con diferencias entre antenas en
S4.7a. Las 640\,000 fases no constituyen 640\,000 capturas independientes.


\FloatBarrier
\subsection{S4.6c: continuidad espacial de la dirección común}

La siguiente pregunta fue si $\psi_{i,u}$ puede interpretarse como un campo
circular suave del entorno. Dado que $\psi\in\mathbb S^1$, la distancia relevante
es
\begin{equation}
d_{\phi}(\psi_i,\psi_j)
=
\left|
\angle\left(e^{j(\psi_i-\psi_j)}\right)
\right|.
\label{eq:circ_distance_psi}
\end{equation}

Se tomó para cada punto de prueba su vecino de entrenamiento espacialmente más cercano dentro
del mismo arreglo. La mediana de distancia entre entrenamiento y prueba reproduce la auditoría de
S4.4d: 0.04829 m. A pesar de esa proximidad, el error circular mediano de copiar
$\psi$ del vecino es $89.75^{\circ}$ para Tx0 y $88.85^{\circ}$ para Tx1. La
media de $\cos(\Delta\psi)$ es 0.0071 y 0.0098, respectivamente, prácticamente
el valor nulo esperado sin alineación útil.

No se observa una asociación monotónica apreciable entre distancia y error: Spearman es $-0.0093$ para Tx0
y $0.0020$ para Tx1. Aun en el cuartil más cercano, donde la distancia mediana
es apenas 7.8 mm, el error continúa alrededor de 90 grados. Tampoco la
concentración dentro de posición muestra una asociación monotónica apreciable con el error del vecino espacial:
Spearman entre $R_{\rm within}$ y error es aproximadamente 0.011--0.013.

\begin{table}[!htbp]
\centering
\small
\caption{S4.6c: evaluación espacial de la fase común mediante el vecino próximo.}
\label{tab:s46c_spatial}
\begin{adjustbox}{max width=\textwidth}
\begin{tabular}{lrr}
\toprule
Métrica & Tx0 & Tx1\\
\midrule
Distancia NN mediana [m] & 0.04829 & 0.04829\\
Error circular NN mediano [$^\circ$] & 89.75 & 88.85\\
Error Q25 [$^\circ$] & 44.37 & 43.95\\
Error Q75 [$^\circ$] & 134.37 & 134.05\\
Media de $\cos(\Delta\psi)$ & 0.0071 & 0.0098\\
Spearman distancia--error & -0.0093 & 0.0020\\
$P(|\Delta\psi|\leq30^\circ)$ & 0.167 & 0.173\\
$P(|\Delta\psi|\leq45^\circ)$ & 0.253 & 0.256\\
\bottomrule
\end{tabular}
\end{adjustbox}
\end{table}

El control entre arreglos en la misma posición produce MedAE $82.45^{\circ}$, media
$85.39^{\circ}$, $P(|\Delta\psi|\leq30^{\circ})=0.192$ y
$P(|\Delta\psi|\leq45^{\circ})=0.286$. Por tanto, tampoco aparece una fase
ambiental común fuerte compartida por ambos arreglos.

El contraste con $\Delta\tau_{\rm eff}$ es notable: en S4.4d, la vecindad
espacial era altamente informativa para el retardo efectivo, mientras que para
$\psi$ el vecino próximo no aporta una ventaja clara incluso en el cuartil con
distancias milimétricas. Esto permite
separar empíricamente dos tipos de discrepancia de fase que no deberían
representarse de la misma manera.

\FloatBarrier
\subsection{S4.6d.1: recuperación del tiempo original y auditoría de trazabilidad}

La ausencia de continuidad espacial hizo plausible una hipótesis alternativa:
$\psi$ podría estar asociado al estado temporal de adquisición o al canal de
referencia. El TFRecord enriquecido de S1 no conservaba el campo \texttt{time},
por lo que fue necesario recuperarlo desde el TFRecord DICHASUS original.

La inspección del archivo fuente confirmó los campos \texttt{cfo},
\texttt{csi}, \path{gt-interp-age-tachy}, \texttt{pos-tachy}, \texttt{snr} y
\texttt{time}. La implementación de \texttt{load\_dataset()} mostró además que:
(i) la calibración modifica únicamente CSI; (ii) la posición se transforma al
sistema de coordenadas de la escena; (iii) el filtro utilizado para S1 es
$-15<y<25$; y (iv) la generación de los 10\,000 registros de caminos trazados ocurre antes
del \texttt{shuffle} utilizado posteriormente para entrenamiento, validación y prueba.

Se reprodujo el parser y el filtro originales para extraer las primeras 10\,000
posiciones filtradas junto con su marca temporal. La comparación contra el TFRecord
S1 produjo error máximo de coordenadas exactamente 0 m. Las 10\,000 posiciones
físicas son únicas y cada una aparece exactamente dos veces en S4.6, una por arreglo.
Esto permitió resolver el tiempo por correspondencia exacta de coordenadas, sin depender de
reconstruir retrospectivamente la permutación de TensorFlow.

Durante una primera prueba se intentó reproducir la permutación aleatoria mediante un conjunto
artificial de índices con semilla 42. Aunque el conjunto cubría 0--9999 exactamente
una vez, la primera fila no coincidió con la asignación histórica y apareció una
discrepancia de 13.26 m. Ese intento fue descartado. No indica un problema en las
particiones anteriores; únicamente demuestra que reconstruir a posteriori el orden
del procesamiento mediante un conjunto artificial de índices no era una estrategia de trazabilidad
suficientemente fuerte. La correspondencia exacta por posición eliminó esa ambigüedad.

El manifiesto temporal definitivo contiene 20\,000 filas posición--arreglo, sin
marcas temporales faltantes y con error máximo XYZ de 0 m. Las 10\,000 posiciones
abarcan
\begin{equation}
8287.224\;{\rm s}\leq t\leq9972.360\;{\rm s},
\end{equation}
con una duración reportada de 1685.137 s ($\approx28.09$ min). Los extremos redondeados dan una diferencia 0.001 s menor, compatible con su precisión de presentación.
Ordenadas cronológicamente las 10\,000 capturas únicas, la
separación mediana entre observaciones es 0.04785 s. El orden del TFRecord no es
cronológico: se detectaron 2192 diferencias no positivas y un salto mínimo de
$-1635.744$ s. Por ello todos los análisis temporales posteriores usan
\texttt{time} explícito y nunca el índice del archivo como sustituto del tiempo físico.

\begin{table}[!htbp]
\centering
\caption{S4.6d.1: auditoría de trazabilidad temporal.}
\label{tab:s46d1_time}
\begin{adjustbox}{max width=\textwidth}
\begin{tabular}{lr}
\toprule
Indicador & Resultado\\
\midrule
Posiciones fuente únicas & 10\,000\\
Filas del manifiesto posición--arreglo & 20\,000\\
Marcas temporales faltantes & 0\\
Error máximo XYZ fuente--S4.6 & 0 m\\
Tiempo mínimo & 8287.224 s\\
Tiempo máximo & 9972.360 s\\
Duración cubierta & 1685.137 s\\
$\Delta t$ cronológico mediano & 0.04785 s\\
Diferencias no positivas en orden de archivo & 2192\\
Mayor salto negativo en orden de archivo & -1635.744 s\\
\bottomrule
\end{tabular}
\end{adjustbox}
\end{table}

Las tres particiones cubren prácticamente toda la ventana temporal: prueba
8292.4--9971.2 s, entrenamiento 8287.2--9972.4 s y validación 8336.0--9967.3 s. Esto
debilita una explicación trivial basada en que cada partición pertenezca a una época
diferente de la campaña.

\FloatBarrier
\subsection{S4.6d.2: continuidad temporal de la dirección común}

Con la marca temporal real disponible se compararon tres predictores diagnósticos:
vecino de entrenamiento más cercano en el tiempo, vecino de entrenamiento más cercano en el espacio y
una muestra aleatoria de entrenamiento. El objetivo no fue proponer un predictor de
despliegue, sino determinar si $\psi$ conserva estructura temporal local.

Para Tx0, el vecino temporal está a una separación mediana de 0.624 s y produce
error circular mediano $90.34^{\circ}$. Para Tx1 el resultado es
$90.56^{\circ}$. Las medianas son próximas a las del vecino espacial y de la referencia
aleatoria; sin intervalos pareados ni un margen de equivalencia prefijado, esta
proximidad no constituye una prueba formal de equivalencia. La media de $\cos(\Delta\psi)$ es $-0.0054$ y $-0.0034$,
y la resultante de la diferencia es sólo 0.0141 y 0.0043.

\begin{table}[!htbp]
\centering
\small
\caption{S4.6d.2: comparación temporal, espacial y aleatoria.}
\label{tab:s46d2_nn}
\begin{adjustbox}{max width=\textwidth}
\begin{tabular}{llrrr}
\toprule
Arreglo & Referencia & MedAE [$^\circ$] & MAE [$^\circ$] & Media cos.\\
\midrule
Tx0 & Vecino temporal & 90.34 & 90.49 & -0.0054\\
Tx0 & Vecino espacial & 89.75 & 89.39 & 0.0071\\
Tx0 & Aleatorio & 90.84 & 90.35 & -0.0048\\
\midrule
Tx1 & Vecino temporal & 90.56 & 90.34 & -0.0034\\
Tx1 & Vecino espacial & 88.85 & 89.33 & 0.0098\\
Tx1 & Aleatorio & 89.18 & 89.00 & 0.0126\\
\bottomrule
\end{tabular}
\end{adjustbox}
\end{table}

La falta de alineación útil también aparece en el intervalo temporal más
corto examinado. Esta escala no excluye dependencia a separaciones menores
que las representadas en los pares analizados. Para pares separados alrededor de 47.9 ms, el error
circular mediano es $88.29^{\circ}$ para Tx0 y $88.26^{\circ}$ para Tx1, con media del
coseno 0.008 en ambos casos. Aumentar la separación temporal desde decenas de milisegundos hasta
centenas de segundos no produce una tendencia sistemática: los errores se
mantienen aproximadamente entre 87 y 91 grados y las medias del coseno oscilan
alrededor de cero.

\begin{table}[!htbp]
\centering
\scriptsize
\caption{S4.6d.2: función temporal de estructura de la fase común. Los errores angulares se expresan en grados; la media del coseno es adimensional.}
\label{tab:s46d2_structure}
\begin{adjustbox}{max width=\textwidth}
\begin{tabular}{lrrrr}
\toprule
Intervalo [s] & Tx0 MedAE & Tx0 media cos. & Tx1 MedAE & Tx1 media cos.\\
\midrule
0--0.075 & 88.29 & 0.008 & 88.26 & 0.008\\
0.075--0.150 & 89.73 & 0.001 & 88.85 & 0.006\\
0.150--0.300 & 91.10 & -0.011 & 88.80 & 0.011\\
0.300--0.600 & 90.47 & -0.003 & 91.24 & -0.008\\
0.600--1.200 & 89.33 & 0.012 & 90.43 & -0.009\\
1.200--2.500 & 87.22 & 0.021 & 89.03 & 0.008\\
2.500--5 & 89.54 & 0.002 & 88.51 & 0.010\\
5--10 & 89.73 & 0.000 & 90.24 & 0.002\\
10--30 & 91.11 & -0.005 & 88.87 & 0.011\\
30--60 & 89.99 & 0.005 & 89.70 & -0.002\\
60--180 & 90.42 & -0.002 & 89.93 & -0.005\\
180--600 & 87.12 & 0.023 & 89.29 & 0.006\\
\bottomrule
\end{tabular}
\end{adjustbox}
\end{table}

Tampoco se detecta relación monotónica entre separación temporal y error:
Spearman es 0.021 para Tx0 y -0.018 para Tx1. Para descartar que una aparente
continuidad temporal estuviera en realidad inducida por movimiento espacial, se
repitió el análisis restringiendo el vecino temporal a pares separados al menos
0.25 m o 0.5 m. Los errores permanecen alrededor o por encima de 90 grados. Por
ejemplo, con separación $\geq0.5$ m se obtienen medianas de $93.52^{\circ}$ para
Tx0 y $99.63^{\circ}$ para Tx1.

Finalmente, el control simultáneo entre Tx0 y Tx1 produce error circular mediano
$82.45^{\circ}$, media $85.39^{\circ}$, media del coseno 0.0632 y resultante de la
diferencia 0.1044. Los indicadores muestran una alineación media débil entre arreglos. Sin un
contraste distribucional no se atribuye a este resultado una conclusión formal
sobre uniformidad o independencia.

\FloatBarrier
\subsection{Conclusión de S4.6: fase común sin predictibilidad local demostrada}

Los resultados S4.6a--d.2 documentan una concentración absoluta baja por antena,
una dirección común descriptiva dentro de cada arreglo y la ausencia de una
ventaja clara de copiar esa dirección desde vecinos espaciales o temporales.
Estos controles no identifican su causa física, no prueban independencia
estadística y no excluyen predictores basados en metadatos adicionales o en
una referencia de adquisición diferente.

Bajo la referencia espectral adoptada, la discrepancia admite la organización
\begin{equation}
\begin{aligned}
\angle\left(H^{\rm real}_{i,u,a,k}H^{{\rm RT}*}_{i,u,a,k}\right)
&\simeq -2\pi\nu_k\Delta\tau_{{\rm eff},i,u}
+\psi_{i,u}+r^{\rm rel}_{i,u,a,k}\pmod{2\pi},\\
r^{\rm rel}_{i,u,a,k}&=\beta_{u,a}+\epsilon_{i,u,a}+r_{\phi,i,u,a,k}
\pmod{2\pi}.
\end{aligned}
\label{eq:s46_final_factorization}
\end{equation}
El residual relativo incluye las fases específicas de antena y el residual
espectral de S4.3; no puede atribuirse íntegramente a trayectorias individuales
sin un análisis adicional. Omitir $\beta$ como predictor dominante no significa
que todos los sesgos de antena sean físicamente nulos.

El retardo efectivo mostró transferencia espacial parcial. Para $\psi$, en
cambio, las referencias locales ensayadas no justifican todavía entrenar un
modelo de alta capacidad con posición o tiempo como únicas entradas. La
clasificación de $\psi$ como parámetro perturbador es una decisión operacional
para el siguiente experimento, condicionada al modelo, la referencia de fase
y las variables observadas. No constituye una demostración de impredecibilidad
para cualquier representación o sistema de adquisición.

S4.6 también modifica la interpretación de $Z_\phi$. En lugar de tratar todos
los grados de libertad de fase como latentes predictibles, se propone
\begin{equation}
Z_\phi
=
\left[
Z_\phi^{\rm pred},
Z_\phi^{\rm nuisance},
Z_\phi^{\rm residual}
\right],
\end{equation}
con $\Delta\tau_{\rm eff}\in Z_\phi^{\rm pred}$,
$\psi_{i,u}\in Z_\phi^{\rm nuisance}$ y la estructura por trayectoria o enlace de S4.3
dentro de $Z_\phi^{\rm residual}$.

Este resultado es metodológicamente relevante para comparar representaciones
como NeRF$^2$ \cite{Zhao2023NeRF2}, Learnable WDT
\cite{Jiang2025LearnableWDT} o RadTwin \cite{Zhang2026RadTwin}. Un modelo del
entorno no debería penalizarse exclusivamente mediante una fase absoluta cuya
predictibilidad no se ha demostrado mediante los controles de posición y tiempo
ejecutados. El benchmark deberá distinguir fidelidad compleja bruta de fidelidad
compleja respecto de una clase de equivalencia declarada.


% ----------------------------------------------------------------------

\section{Síntesis científica de S1--S4.6}
\label{sec:synthesis}
% ----------------------------------------------------------------------

\textbf{El resultado principal de S4 es una mejora comprobada de la concordancia
de fase mediante una corrección parcialmente predecible.} El retardo efectivo
predicho reduce el error angular del canal y aumenta su coherencia en posiciones
reservadas. La rotación constante continúa estimándose con la medición evaluada;
por tanto, el alcance demostrado es una corrección transferible condicionada,
y no la predicción completa de la fase absoluta.

\subsection{Tres preguntas que delimitan el aporte}

\begin{enumerate}
\item \textbf{¿Existe concordancia de fase en bruto?} B0--B2 presentan errores
angulares próximos a 90 grados, pese a las mejoras de potencia, magnitud y
RMS-DS de B1/B2. La calibración macroscópica no resuelve esa discrepancia.
\item \textbf{¿Puede alinearse el canal con la medición?} Sí. Una pendiente
espectral común al arreglo, junto con rotaciones constantes por enlace,
reduce sustancialmente el error. Es una referencia diagnóstica obtenida
utilizando el canal real, que revela una estructura de corrección de baja dimensión.
\item \textbf{¿Puede predecirse esa corrección sin consultar el canal objetivo?}
Parcialmente. S4.4 predice el retardo desde posición y descriptores RT; S4.5
comprueba que ese retardo mejora el canal. La fase constante aún se ajusta
con la medición de prueba y su dirección común no mostró transferencia útil
mediante los vecinos de posición o tiempo de S4.6.
\end{enumerate}

La respuesta a la segunda pregunta no sustituye a la tercera. A su vez,
la falta de predicción completa de fase no invalida la mejora conseguida
mediante el retardo predicho. Son niveles distintos de evidencia.

\subsection{Mejora de fase alcanzada y condiciones de acceso a datos}

El \cref{tab:phase_contribution_summary} reúne las comparaciones que permiten
cuantificar el aporte predictivo. En cada protocolo se mantiene el mismo
tratamiento del desfase constante por enlace, denominado CPO en los evaluadores;
el cambio entre referencia y predictor es la incorporación del retardo predicho.
Los indicadores son medianas de errores angulares por enlace y de coherencias,
según la \cref{sec:metricas_avances}.

\begin{table}[!htbp]
\centering\small
\caption{Mejora de fase en S4.5 y origen de la corrección. El desfase constante
por enlace se ajusta con la medición evaluada en todas las filas.}
\label{tab:phase_contribution_summary}
\begin{tabularx}{\textwidth}{P{2.4cm}Lrr}
\toprule
Protocolo & Corrección & Error [$^\circ$] & Coherencia\\
\midrule
Denso & Sólo CPO & 81.42 & 0.156\\
Denso & Retardo predicho + CPO & \textbf{45.59} & \textbf{0.634}\\
Denso & Retardo oráculo + CPO & 31.63 & 0.745\\
\midrule
Por bloques & Sólo CPO & 81.49 & 0.157\\
Por bloques & Retardo predicho + CPO & \textbf{64.08} & \textbf{0.448}\\
Por bloques & Retardo oráculo + CPO & 31.46 & 0.746\\
\bottomrule
\end{tabularx}
\end{table}

En interpolación densa, el predictor reduce la mediana angular en 35.83 grados
respecto de la referencia con CPO y recupera aproximadamente el 72\% del
beneficio angular oráculo y el 81\% del beneficio de coherencia. Bajo los cinco
bloques espaciales, la reducción es de 17.41 grados y las fracciones recuperadas
son 34.8\% y 49.3\%, respectivamente. Estas fracciones comparan diferencias de
indicadores agregados; no son probabilidades de éxito por enlace ni porcentajes
de fase física recuperada. La magnitud de la mejora es sustancial en esos
indicadores, sin atribuirle una significación estadística no contrastada.

La condición predictiva se refiere a $\widehat{\Delta\tau}_{\rm eff}$:
su cálculo en prueba no utiliza $H_{\rm real,test}$. El canal corregido completo
sí utiliza esa medición para estimar el CPO. La coherencia normalizada es
invariante a una rotación constante por enlace; por ello su aumento confirma
el efecto de la corrección espectral y no puede atribuirse al ajuste de CPO.
La mediana de 45.59 grados de S4.5 y la de 45.97 grados de S4.6 describen
objetos distintos: la primera resume errores del CFR corregido, mientras que
la segunda resume fases constantes después del centrado circular entre antenas.
No forman una secuencia de mejoras acumulativas sobre una misma métrica.

\subsection{De la alineación a una estructura parcialmente transferible}

S4.1 motiva caracterizar el error de fase antes de utilizarlo para ajustar
materiales: los ensayos de PEOC/UPEC/PEAC muestran sensibilidad al tratamiento
de esa incertidumbre y a la optimización. Su función es diagnóstica, coherente
con la calibración estudiada por Ruah et al.\ \citep{Ruah2024Phase}; no establecen
un ganador universal para dc01.

S4.2 aporta evidencia complementaria de una pendiente espectral reutilizable.
Los retardos estimados correlacionan entre B0/B1/B2 (0.886, 0.834 y 0.819).
El retardo ajustado con 16 antenas mejora las otras 16 hasta 35.52 grados en
B0, conservando CPO libre en las antenas evaluadas. En la prueba espectral
más estricta, retardo e intercepto se transfieren juntos: pares a impares
alcanzan 31.65 grados, iguales al oráculo a la precisión reportada. Esto
acredita interpolación espectral intercalada; la transferencia entre medias
bandas presenta errores mayores. Un error de retardo de 2.5 ns induce 22.5
grados a 25 MHz de separación, explicando la sensibilidad a esa extrapolación.

S4.4 añade la evidencia predictiva. La proximidad mediana de 4.83 cm en la
partición original identifica un régimen de interpolación densa. Con regiones
reservadas, el MAE medio pasa de 27.56 ns usando posición a 19.29 ns usando
RT y a 18.69 ns combinando ambos. En el bloque 0, con separación mediana de
7.33 m, posición y RT alcanzan 48.66 y 21.76 ns. Estos resultados muestran
utilidad de los descriptores del trazado fuera de la vecindad inmediata de
calibración, dentro de dc01; no equivalen a transferencia entre edificios.
S4.5 completa la evidencia al vincular esa predicción auxiliar con la mejora
del canal descrita en el \cref{tab:phase_contribution_summary}.

\subsection{Naturaleza diferenciada de los componentes de fase}

La discrepancia se organiza operacionalmente como
\begin{equation}
\begin{aligned}
\Delta\phi_{i,u,a,k}\simeq{}&
\underbrace{-2\pi\nu_k\Delta\tau_{{\rm eff},i,u}}
_{\text{parcialmente predecible}}
+\underbrace{\psi_{i,u}}_{\text{estimada con la captura}}\\
&+\underbrace{r^{\rm rel}_{i,u,a,k}}
_{\text{estructura relativa aún no modelada por completo}}
\pmod{2\pi}.
\end{aligned}
\label{eq:phase_contribution_factorization}
\end{equation}
Esta ecuación resume funciones inferenciales, no identifica tres causas físicas
independientes. Conforme a la \cref{eq:s46_final_factorization},
$r^{\rm rel}=\beta+\epsilon+r_\phi$ incluye las fases relativas por antena y el
residual espectral. Así se evita confundir el residual de S4.3, obtenido tras
ajustar interceptos por enlace, con el que queda al retirar sólo una rotación
común al arreglo.

S4.3 caracteriza estructura espectral considerable: la resultante de incrementos
es aproximadamente 0.997 a una subportadora, 0.939 a 16, 0.838 a 64 y 0.738 a
256. Las asociaciones con distribución energética son más informativas que el
conteo bruto de caminos. Sin embargo, concentración y asociación no prueban
por sí solas ausencia de ruido, rechazo de independencia o causalidad física.
Ese residual sigue siendo un objeto de modelado, motivando S4.7b.

S4.6 estima una dirección común dentro de cada captura, con concentración
mediana 0.4216. Su copia desde vecinos espaciales o temporales produce errores
cercanos a 90 grados, sin ventaja clara frente al control aleatorio. La decisión
actual es tratar $\psi$ como parámetro perturbador respecto del predictor del
entorno. Esto no identifica el CPO físico del hardware ni demuestra imposibilidad
universal de predicción. S4.7a evaluará fidelidad relativa sin exigir reproducir
esa referencia angular, controlando qué información útil se pierde al eliminarla.

\FloatBarrier
\subsection{Balance de hipótesis y resultados}

El \cref{tab:phase_hypotheses_summary} resume el papel de cada experimento.
Se distingue evidencia favorable, alcance limitado y contraste aún pendiente;
la palabra rechazo se reserva para hipótesis con un procedimiento estadístico
explícito, no para la sola inspección de indicadores.

\begingroup\small
\begin{longtable}{P{4.1cm}P{1.6cm}P{7.9cm}}
\caption{Hipótesis de fase, experimentos y alcance de la evidencia.}
\label{tab:phase_hypotheses_summary}\\
\toprule
Hipótesis o pregunta & Bloque & Resultado y alcance\\
\midrule
\endfirsthead
\toprule
Hipótesis o pregunta & Bloque & Resultado y alcance\\
\midrule
\endhead
La discrepancia carece de estructura aprovechable & S4.2--3 &
La alineación y la estructura espectral muestran organización aprovechable;
no se excluyen componentes aleatorias.\\
Un desfase constante basta & S4.2 &
Insuficiente: el error de B0 permanece alrededor de 81.4 grados; añadir retardo
reduce sustancialmente la discrepancia.\\
Una pendiente espectral explica una fracción importante & S4.2b &
Apoyada por alineación oráculo; no es todavía predicción en una consulta nueva.\\
El retardo tiene estructura compartida entre modelos & S4.2c &
Correlaciones de 0.819--0.886; no implican identidad ni una causa física única.\\
El retardo se comparte entre antenas & S4.2d &
Transferencia parcial entre grupos de antenas, con CPO libre en evaluación.\\
La corrección se transfiere entre frecuencias & S4.2e &
Retardo e intercepto se transfieren sin reajuste; interpolación intercalada
más favorable que extrapolación entre medias bandas.\\
El residual es independiente e idénticamente distribuido & S4.3a &
La estructura observada motiva un contraste que preserve marginales y ajuste;
el rechazo formal de independencia no está establecido.\\
El número bruto de caminos explica el error & S4.3b &
Escasa capacidad explicativa en los indicadores examinados.\\
La distribución de energía está asociada al residual & S4.3c--e &
Asociaciones documentadas, dependientes del nivel de agregación;
no acreditan causalidad.\\
El retardo puede predecirse con variables disponibles & S4.4 &
Sí, desde posición y/o RT; el protocolo espacial determina el alcance.\\
El rendimiento denso representa generalización regional & S4.4d &
La proximidad de entrenamiento y prueba delimita ese rendimiento como
interpolación densa.\\
Los descriptores RT aportan fuera de la vecindad inmediata & S4.4f &
Reducción del MAE medio respecto de posición, en los cinco bloques de dc01.\\
El retardo predicho mejora el canal complejo & S4.5 &
Menor error angular y mayor coherencia; fase constante todavía ajustada
con la medición evaluada.\\
La mejora se conserva íntegramente al separar regiones & S4.5b &
La mejora persiste, pero disminuye: 64.08 grados y coherencia 0.448.\\
Una fase absoluta fija por antena explica la discrepancia & S4.6b &
Baja concentración marginal; no excluye sesgos relativos estables ocultos
por rotaciones comunes.\\
Existe una dirección común estimable dentro del arreglo & S4.6b &
Concentración descriptiva 0.4216; falta cuantificar el efecto del ajuste
sobre las mismas antenas mediante controles específicos.\\
Copiar la dirección común desde vecinos espaciales es útil & S4.6c &
No se observa una ventaja clara con las distancias y la regla evaluadas.\\
Copiarla desde vecinos temporales es útil & S4.6d.2 &
No se observa una ventaja clara en los intervalos examinados; no es una prueba
de impredecibilidad para cualquier modelo.\\
\bottomrule
\end{longtable}
\endgroup

La contribución consolidada es, por tanto, identificar una componente de
retardo de baja dimensión, parcialmente transferible y útil para mejorar el
CFR complejo; distinguirla de una dirección común condicionada a la captura;
y caracterizar la estructura relativa restante. Este avance establece qué
parte de la fase puede corregirse con información predictiva y qué parte
requiere otro tratamiento. La fidelidad diferencial ante movimiento humano
continúa pendiente: los avances de reconstrucción estática no la sustituyen.

% ----------------------------------------------------------------------
\FloatBarrier
\section{Cómo se relaciona este trabajo con las representaciones del estado del arte}
\label{sec:sota_positioning}
% ----------------------------------------------------------------------

\FloatBarrier
\subsection{NeRF2}

NeRF$^2$~\cite{Zhao2023NeRF2} introduce un campo neuronal de radiofrecuencia para
representar de forma continua el campo de radio y predecir señales en
posiciones no observadas. Su papel en esta tesis debe ser el de un método de referencia
\emph{campo neuronal}. No debe considerarse equivalente a B2: B2 modifica las
propiedades electromagnéticas dentro de un trazador de rayos; NeRF$^2$ aprende una
representación continua del campo.

El protocolo espacial desarrollado en S4.4d--f proporciona ahora una forma
mucho más rigurosa de evaluar NeRF$^2$ que una partición aleatoria. La comparación
debe incluir, al menos, interpolación espacial, extrapolación por bloques espaciales,
curvas de cantidad de datos y posteriormente fidelidad diferencial.

\FloatBarrier
\subsection{Learnable Wireless Digital Twins}

Learnable WDT~\cite{Jiang2025LearnableWDT} propone una arquitectura modular
que segmenta el entorno en objetos, aprende representaciones de sus
propiedades electromagnéticas y aprende efectos de interacción. Su motivación
incluye reutilizar conocimiento cuando cambia la disposición del entorno.

Este trabajo es especialmente relevante para cuestionar B2. La función de B2
está asociada a coordenadas globales; si un objeto se mueve, las propiedades
aprendidas no necesariamente se trasladan con él. Por ello la comparación
correcta debe introducir una intervención física o simulada:

\begin{equation}
E_0
\xrightarrow{\text{mover objeto}}
E_1.
\end{equation}

Se propone comparar B2 global, B2 con coordenadas locales al objeto y
Learnable WDT bajo sin adaptación y adaptación con pocas mediciones.

\FloatBarrier
\subsection{RadTwin}

RadTwin~\cite{Zhang2026RadTwin} condiciona explícitamente la representación en
la geometría de la escena mediante nubes de puntos y utiliza información de trazado de rayos para limitar las contribuciones relevantes antes de un decodificador neuronal.
Su objetivo de generalización a configuraciones dinámicos lo convierte en un método de referencia
natural cuando la tesis pase de posiciones nuevas dentro de una escena a
cambios del propio entorno.

\FloatBarrier
\subsection{WaveVerse}

WaveVerse~\cite{Zheng2026WaveVerse} combina generación de mundos 4D con un
simulador de RF coherente en fase y presenta experimentos de dinámica,
conformación de haz, formación de imágenes y seguimiento respiratorio. Por ello la novedad de esta
tesis no puede formularse simplemente como ``simular respiración mediante RT
coherente''. La pregunta diferencial debe ser:

\begin{quote}
¿Qué fidelidad estática, de fase y diferencial debe conservar el gemelo para
que la señal de sensado simulada concuerde con mediciones independientes y
siga siendo utilizable bajo perturbaciones no deseadas, cambios de escena y hardware?
\end{quote}

\FloatBarrier
\subsection{RayLoc y problemas inversos}

RayLoc~\cite{Han2026RayLoc} formula localización como la inversión de un
simulador RT diferenciable. Para esta tesis es una referencia importante al
pasar de

\begin{equation}
\mathbf q\rightarrow H
\end{equation}

a

\begin{equation}
H\rightarrow\widehat{\mathbf q},
\end{equation}

pero no debe utilizarse todavía como método de referencia del modelo directo de B0--B2. Su relevancia
aumentará cuando la variable desconocida sea posición/movimiento humano.

% ----------------------------------------------------------------------
\FloatBarrier
\section{Conexión con sensado Wi-Fi y LoRa}
\label{sec:wifi_lora}
% ----------------------------------------------------------------------

Las revisiones de sensado Wi-Fi \citep{Zhu2024Survey,Armenta2024Survey}
identifican como observable relevante la
variación de CFR inducida por cambios de propagación; el CSI disponible en
dispositivos comerciales es una muestra discreta de esa CFR. Este hecho
justifica trabajar primero con $H$ complejo antes de pasar al modelo de
instrumentación Wi-Fi.

LoRa requiere otra observación. El trabajo de Zhang et
al.~\cite{Zhang2020LoRa} demuestra experimentalmente que el movimiento de un
objetivo puede inducir variaciones detectables sobre señales LoRa a larga
distancia y a través de paredes, incluyendo respiración. SenLoRa
\cite{SenLoRa2025} formula una extracción de señal de sensado a partir de
la diferencia temporal del canal y subraya la necesidad de compensar CFO, SFO
y fluctuación de fase para sensado de movimientos pequeños. OpenLoRa
\cite{Mishra2023OpenLoRa} proporciona una plataforma, muestras I/Q y escenarios
experimentales uniformes para validar recepción/demodulación LoRa.

La arquitectura multirradio conserva una escena común $\Gamma(t)$ y genera
caminos electromagnéticos específicos de cada modalidad:
\begin{equation}
\Gamma(t)\longrightarrow\mathcal P_m(t)\longrightarrow H_m(t)
\xrightarrow{\mathcal O_m}Y_m(t),\qquad m\in\{\mathrm{WiFi},\mathrm{LoRa}\}.
\end{equation}
El observable Wi-Fi es una estimación discreta del canal y el observable LoRa
inicial es una secuencia I/Q. Frecuencia, antenas y propiedades dispersivas
impiden identificar $\mathcal P_{\rm WiFi}$ con $\mathcal P_{\rm LoRa}$.

% ----------------------------------------------------------------------
\FloatBarrier
\section{Factorización latente que emerge de los resultados}
\label{sec:latent_model}

Se mantiene la arquitectura global $Z=[Z_E,Z_\phi,Z_R,Z_H,Z_F,Z_N]$, pero
S4.6 exige distinguir la función inferencial de cada componente de fase:
\begin{equation}
Z_\phi=[Z_\phi^{\rm pred},Z_\phi^{\rm nuisance},Z_\phi^{\rm residual}].
\label{eq:phase_latent_roles}
\end{equation}
$Z_\phi^{\rm pred}$ contiene cantidades cuya predicción puede contrastarse con
información disponible en una consulta, como $\Delta\tau_{\rm eff}$.
$Z_\phi^{\rm nuisance}$ agrupa grados de libertad que, bajo los controles
actuales, requieren calibración, marginalización o invariancia, como
$\psi_{i,u}$. $Z_\phi^{\rm residual}$ contiene la discrepancia restante por
enlace, frecuencia y eventualmente trayectoria. El término inglés
\emph{nuisance} designa aquí un parámetro perturbador respecto de una tarea,
no una atribución causal al hardware.

Las categorías anteriores por escala ---componente de baja dimensión, entorno,
enlace y trayectoria--- se conservan como descriptores auxiliares. No equivalen
a la partición por predictibilidad: un término de baja dimensión puede carecer
de un predictor útil, y un residual de enlace puede conservar estructura.
$Z_E$ representa entorno; $Z_R$, configuración y referencia de radio;
$Z_H$, dinámica humana; $Z_F$, fisiología; y $Z_N$, ruido e interferencia.
La relación entre $Z_\phi$ y $Z_R$ exige una convención de referencia explícita
para evitar representar dos veces la misma rotación.

La clasificación no presupone independencia causal ni un codificador ya
aprendido. Los experimentos validan descriptores y reglas de tratamiento;
la representación latente se evaluará posteriormente por transferencia,
intervenciones y conservación de información útil para la tarea. La decisión
actual sobre $\psi$ podrá revisarse si nuevos metadatos o adquisiciones con
referencia estable muestran predictibilidad fuera de muestra.

\FloatBarrier
\section{Ruta experimental inmediata y futura}
\label{sec:future_route}
% ----------------------------------------------------------------------

\FloatBarrier
\subsection{S4.7a: fidelidad relativa e invariancia a la fase común}
\label{sec:s47a_protocol}

El siguiente bloque evalúa B0 sin convertir $\psi$ en una nueva etiqueta para
un predictor de alta capacidad. Las \cref{eq:spatial_gram_invariant,eq:cross_frequency_invariant,eq:phase_quotient_distance}
definen tres observables complementarios: productos espaciales $G_u(k)$,
productos entre frecuencias $K_u(k,l)$ y distancia del canal apilado respecto
de una única fase común por arreglo. Se evaluarán primero en la partición
original y luego en los bloques de S4.4f, con el predictor de retardo congelado
y entrenado bajo el protocolo correspondiente.

Para $G$, las versiones con y sin corrección de retardo común deben coincidir.
Esta igualdad será un control numérico del procesamiento y no una mejora
esperada. Para $K$ y para la distancia respecto de fase común, en cambio, el
retardo modifica la estructura espectral y la comparación resulta informativa.
El experimento se organiza como sigue:
\begin{enumerate}
\item cuantificar potencia, magnitud y error complejo bruto sobre un soporte común;
\item comparar productos espaciales medidos y simulados, separando diagonales
(potencia) y términos fuera de diagonal (estructura relativa);
\item evaluar $K$ y la distancia respecto de fase común sin corregir retardo,
con retardo predicho y con retardo oráculo usado sólo como referencia diagnóstica;
\item contrastar productos relativos entre antenas antes y después de una
calibración relativa fija, si ésta se estima exclusivamente con entrenamiento;
\item comprobar estabilidad por arreglo, región, potencia y concentración
interna, e informar las exclusiones y la distribución de errores.
\end{enumerate}

Una métrica propuesta para la estructura espacial es
\begin{equation}
E_{G,i,u}=\frac{\sum_{k\in\mathcal K}\|
\widehat G_{i,u}(k)-G^{\rm real}_{i,u}(k)\|_F^2}
{\sum_{k\in\mathcal K}\|G^{\rm real}_{i,u}(k)\|_F^2}.
\label{eq:s47_gram_error}
\end{equation}
Se reportará también una versión con normalización de energía, distinguiendo
el error de estructura del error de escala. Para $K$ se sustituye la suma por
un conjunto de pares $(k,l)$ y separaciones espectrales prefijados. Las
normalizaciones se definen sólo cuando el denominador es válido; los umbrales
se fijan con entrenamiento y validación. Las potencias bajas requieren atención
porque los productos de canales estimados contienen ruido y sesgo en la diagonal.

La distancia de la \cref{eq:phase_quotient_distance} aplica una sola rotación
al vector completo de cada arreglo. Es distinta del ajuste libre por enlace
de S4.5: su valor mide concordancia en una clase de equivalencia, sin presentar
la rotación óptima como predicción instrumental. Los errores se agregan por
posición y se comparan de forma pareada entre métodos; las 32 antenas no se
cuentan como réplicas independientes.

Los controles mínimos son: rotación común sintética, que debe dejar invariantes
$G$, $K$ y la distancia cociente; retardo común sintético, que debe conservar
$G$ y transformar $K$ según la \cref{eq:cross_frequency_invariant}; y rotaciones
independientes por antena, que deben alterar los términos relativos. Se
examinará además la calibración de diferencias de fase entre antenas: una
resultante marginal pequeña en S4.6b no excluye un sesgo relativo estable
oculto por $\psi$. El criterio de cierre será una separación reproducible entre
error de escala, error relativo y error de referencia, con incertidumbre pareada
y sin utilizar la medición de prueba para entrenar correctores.

Este bloque no demuestra todavía fidelidad de sensado. La
\cref{sec:invariance_sensing_limit} establece qué sensibilidad puede perderse
al eliminar fase común; S6 comparará el observable invariante y el canal con
referencia basal fija ante una perturbación mecánica conocida.

\FloatBarrier
\paragraph{Ampliación operacional de la comparación.}
El protocolo se implementará con $\mathsf V_0$--$\mathsf V_7$ de la
\cref{sec:representation_v_family}. El canal canónico sólo se considerará desplegable
cuando su referencia sea causalmente disponible; en ausencia de ella se conservará
como diagnóstico. Las 4900 posiciones de prueba densa y los cinco bloques de 2000
posiciones tendrán agregaciones separadas. La incorporación de CAEZ será una validación
adicional propuesta tras auditar su adquisición, no una extensión ya ejecutada de dc01.
E0--E2 compararán estabilidad y sensibilidad para decidir qué representación integra
cada rama. Así, el cierre estático de S4.7a no se confundirá con validación de sensado.

\subsection{S4.7b: incertidumbre circular y por trayectoria}

Después de retirar los componentes de baja dimensión, se comparará primero la
incertidumbre circular de la fase común o residual agregada y, posteriormente, la
incertidumbre por trayectoria de PEAC. No se confundirá una distribución posterior
estimada con el canal evaluado con una predicción disponible en despliegue. La
comparación distinguirá dos escenarios:

\begin{enumerate}
\item inferencia posterior de errores de trayectoria usando la medición
evaluada;
\item predicción transferible de una distribución de error desde descriptores
físicos del camino.
\end{enumerate}

Una extensión posible es

\begin{equation}
(\mu_\ell,\kappa_\ell)
=
g_\theta(
d_\ell,
\theta_\ell,
\text{tipo de interacción},
\text{material},
\text{contexto}
),
\end{equation}

evaluada sobre posiciones que no participaron en el ajuste.

Para una variable circular $\psi$, los baselines serán uniforme, media circular
global, von Mises global, von Mises condicionada y mezcla condicionada:
\begin{equation}
p(\psi\mid x)=\sum_{m=1}^{M}\pi_m(x)
\frac{\exp\{\kappa_m(x)\cos[\psi-\mu_m(x)]\}}
{2\pi I_0(\kappa_m(x))},
\qquad \sum_m\pi_m(x)=1.
\label{eq:s47b_von_mises_mixture}
\end{equation}
Las entradas candidatas incluirán potencia RT, dispersión de retardos, número efectivo
de trayectorias, fracción dominante, longitudes, interacciones, diferencias B0--B2,
posición y arreglo. Toda selección se ajustará dentro del entrenamiento de cada bloque.

Se reportarán log-verosimilitud circular negativa, error angular, calibración de
regiones predictivas y riesgo frente a cobertura al abstenerse. La cobertura se medirá
en posiciones o sesiones independientes, no sobre las muestras usadas para ajustar
$\kappa$. Se contrastará el modelo condicionado contra los baselines mediante diferencias
pareadas e intervalos agrupados. Si no mejora de forma reproducible, la conclusión será
que esos descriptores no identifican la fase residual bajo el protocolo, sin extrapolar
una imposibilidad universal ni forzar una estimación puntual sobreconfiante.

\FloatBarrier
\subsection{S5: comparación sistemática de representaciones}

El siguiente bloque mayor debe comparar:

\begin{equation}
\text{B0/B1/B2}
\;\;{\rm vs}\;\;
\text{NeRF}^2
\;\;{\rm vs}\;\;
\text{Learnable WDT}
\;\;{\rm vs}\;\;
\text{RadTwin}.
\end{equation}

Debe conservar un protocolo común:
\begin{enumerate}
\item mismo conjunto físico y mismas posiciones;
\item partición espacial densa;
\item reserva por bloques espaciales;
\item curvas de eficiencia de datos;
\item potencia, dispersión de retardos, magnitud, NMSE complejo y métricas de fase;
\item coste de entrenamiento/inferencia;
\item adaptación ante cambio de objetos;
\item posteriormente, fidelidad de $\Delta H$.
\end{enumerate}

NeRF$^2$ debe entrar primero, porque su contribución se puede evaluar en la
misma escena estática. Learnable WDT y RadTwin son más informativos cuando se
introducen cambios de disposición u objetos.

\FloatBarrier
\subsection{S6: fidelidad diferencial y movimiento controlado}

La hipótesis central de sensado se formulará mediante

\begin{equation}
\Delta H
=
H(\mathbf x_H+\Delta\mathbf x_H)
-
H(\mathbf x_H),
\end{equation}

y localmente

\begin{equation}
\Delta H
\simeq
J_H
\Delta\mathbf x_H,
\qquad
J_H=
\frac{\partial H}{\partial\mathbf x_H}.
\label{eq:jacobian_sensing}
\end{equation}

La secuencia experimental recomendada es:
\begin{enumerate}
\item reflector estático;
\item desplazamiento conocido;
\item movimiento sinusoidal conocido;
\item movimiento humano grueso con Tx/Rx fijos;
\item respiración de una persona;
\item respiración con movimientos corporales no asociados a la respiración;
\item múltiples personas.
\end{enumerate}

La calibración debe fijarse en el estado basal. Recalibrar libremente el gemelo
en cada instante podría absorber la perturbación física que se desea detectar.

\FloatBarrier
\subsection{S7: operador LoRa y validación multirradio}

La capa LoRa debe validarse inicialmente contra IQ real, aprovechando
OpenLoRa y conjuntos de datos con condiciones LOS/NLOS, interferencia y colisiones.
El módulo de síntesis debe aplicar el canal físico a chirps CSS y añadir explícitamente
CFO, STO/SFO, fluctuación de fase, SNR e interferencia. SenLoRa ofrece una referencia
directa sobre la necesidad de eliminar offsets antes de utilizar continuidad de
fase para sensado de movimientos pequeños.

% ----------------------------------------------------------------------
\FloatBarrier
\section{Criterios de cierre}
% ----------------------------------------------------------------------

\begingroup
\small
\begin{longtable}{P{2.0cm}P{4.5cm}P{7.1cm}}
\caption{Criterios propuestos para cerrar cada bloque experimental.}
\label{tab:closure_criteria}\\
\toprule
Etapa & Pregunta & Criterio de cierre\\
\midrule
\endfirsthead
\toprule
Etapa & Pregunta & Criterio de cierre\\
\midrule
\endhead
B0--B2 &
¿Mejora la reconstrucción del entorno? &
Métricas macro y de magnitud estables con artefactos reproducibles.\\
S4.1 &
¿La calibración es sensible a errores de fase? &
Reproducción estable de PEOC/UPEC/PEAC y caracterización de fallos.\\
S4.2 &
¿Existe una componente de baja dimensión de fase? &
Transferencia demostrada entre antenas y frecuencia, separando oráculo de
predicción.\\
S4.3 &
¿Qué estructura estadística conserva el residual? &
Asociaciones multitrayecto robustas, con controles de dependencia y remuestreo agrupado.\\
S4.4--S4.5 &
¿Puede predecirse $\Delta\tau_{\rm eff}$? &
Predicción del retardo sin usar su medición objetivo y beneficio en el canal
condicionado al CPO ajustado, bajo particiones densa y por bloques.\\
S4.6 &
¿Qué estructura conserva la fase constante? &
Diagnóstico circular y temporal documentado; alcance limitado a los controles ejecutados.\\
S4.7a &
¿Qué fidelidad conserva un observable invariante? &
Separar escala y estructura relativa; validar invariancia común y sensibilidad al retardo entre frecuencias.\\
S4.7b &
¿Hace falta incertidumbre por trayectoria? &
Mejora estable en datos reservados frente a PEOC/UPEC, sin usar la medición de prueba para ajustar la predicción.\\
S5 &
¿Qué representación generaliza mejor? &
Comparación común con NeRF$^2$, WDT y RadTwin.\\
S6 &
¿Se preserva $\Delta H$? &
Movimiento conocido reproducido sin recalibrar la señal objetivo.\\
S7 &
¿Se transfiere a radios de sensado reales? &
Wi-Fi/LoRa validados con operadores y perturbaciones instrumentales específicos.\\
\bottomrule
\end{longtable}
\endgroup

% ----------------------------------------------------------------------
\FloatBarrier
\section{Programas y artefactos experimentales}
\label{sec:artifacts}
\label{sec:trazabilidad_avances}
% ----------------------------------------------------------------------

El cuadro siguiente conserva los identificadores del registro experimental. Para S4.6 se identifican los productos descritos en el documento de resultados; sus nombres de archivo y versiones de programa no se especifican en el material recibido y deberán vincularse al manifiesto de ejecución.
Las rutas corresponden al entorno de ejecución de los experimentos; su inclusión
documenta la procedencia de los resultados y permite localizar los artefactos
necesarios para una reproducción independiente.

\begingroup
\small
\begin{longtable}{P{2.6cm}P{11.0cm}}
\caption{Trazabilidad resumida de programas y artefactos.}
\label{tab:artifacts}\\
\toprule
Bloque & Programas y artefactos principales\\
\midrule
\endfirsthead
\toprule
Bloque & Programas y artefactos principales\\
\midrule
\endhead
S1 &
\path{s1_01_check_dc01.py},
\path{s1_02_trace_real_point.py},
\path{gen_dataset.py};
TFRecord final de 10\,000 posiciones en
\path{/mnt/c/rf-sensing-data/S1_dc01/traced_paths/}.\\
B0 &
\path{/itu_baseline/test_eval_full/}; escala congelada y métricas por
enlace.\\
B1 &
\path{s2_01_fit_learned_materials.py},
\path{s2_02_evaluate_learned_materials_test.py};
estado entrenado \path{/learned_materials/final/learned_materials}.\\
B2 &
\path{neural_materials.py};
estado entrenado \path{/neural_materials/final/neural_materials}.\\
S4.2b &
\path{s4_02b_coherence_optimal_delay_dc01.py}.\\
S4.2e &
\path{s4_02e_frequency_heldout_delay_dc01.py}.\\
S4.3a &
\path{s4_03a_residual_phase_diagnostics_dc01.py}.\\
S4.3c &
\path{s4_03c_inspect_path_cir_dc01.py}.\\
S4.3d &
\path{s4_03d_multilevel_phase_analysis_dc01.py}.\\
S4.3e &
\path{s4_03e_cluster_bootstrap_dc01.py}.\\
S4.4a &
\path{s4_04a_build_tau_prediction_dataset_dc01.py};
\path{/s4_4a_tau_prediction_dataset/B0/}.\\
S4.4b &
\path{s4_04b_tau_prediction_baselines_dc01.py}.\\
S4.4c &
\path{s4_04c_tau_nonlinear_ablation_dc01.py}.\\
S4.4d &
\path{s4_04d_spatial_generalization_audit_dc01.py}.\\
S4.4e &
\path{s4_04e_prepare_spatial_blocks_dc01.py}.\\
S4.4f &
\path{s4_04f_spatial_blocked_prediction_dc01.py};
60\,000 filas de predicción, sin identificadores ausentes ni claves duplicadas.\\
S4.5a &
\path{s4_05a_end_to_end_predicted_tau_dc01.py}.\\
S4.5b &
\path{s4_05b_spatial_blocked_end_to_end_dc01.py};
\path{/s4_5b_spatial_blocked_end_to_end/B0/}.\\
S4.6a--b &
Registro de 640\,000 fases por posición, arreglo y antena; agregados circulares
por partición y descomposición de la dirección común.\\
S4.6c &
Correspondencias entre vecinos de entrenamiento y prueba, errores angulares y
control entre arreglos.\\
S4.6d.1 &
Manifiesto de 20\,000 filas posición--arreglo, con tiempo recuperado y
correspondencia XYZ exacta; permutación retrospectiva descartada.\\
S4.6d.2 &
Comparaciones temporal, espacial y aleatoria; estructura por intervalos de tiempo
y controles con separación espacial mínima.\\
\bottomrule
\end{longtable}
\endgroup

% ----------------------------------------------------------------------
\FloatBarrier
\section{Hipótesis doctoral consolidada}
% ----------------------------------------------------------------------

La evidencia acumulada conduce a una hipótesis más precisa que
``construir un gemelo inalámbrico neuronal'':

\begin{quote}
Una representación útil para sensado por radiofrecuencia debe separar la estructura
electromagnética del entorno de los grados de libertad de fase y hardware,
generalizar fuera de los puntos de calibración y preservar la respuesta
diferencial del canal frente a perturbaciones humanas. La fidelidad de reconstrucción estática y la fidelidad de las perturbaciones
se evalúan por separado; ninguna sustituye la validación del observable
utilizado para sensado.
\end{quote}

Formalmente, dos modelos pueden satisfacer

\begin{equation}
{\rm NMSE}(H_{M_1},H_{\rm real})
<
{\rm NMSE}(H_{M_2},H_{\rm real}),
\end{equation}

pero invertir el orden en

\begin{equation}
\left\|
\frac{\partial H_{M_1}}{\partial x_H}
-
\frac{\partial H_{\rm real}}{\partial x_H}
\right\|
>
\left\|
\frac{\partial H_{M_2}}{\partial x_H}
-
\frac{\partial H_{\rm real}}{\partial x_H}
\right\|.
\end{equation}

Demostrar o refutar esa separación entre \emph{fidelidad de reconstrucción del canal} y \emph{fidelidad de sensado} constituye uno de los ejes más prometedores
de la tesis.

% ----------------------------------------------------------------------
\FloatBarrier
\section{Estado actual del programa}

El registro documenta B0--B2 y S4.1--S4.6d.2. S4.6 completa la caracterización
descriptiva y los controles locales previstos para el desfase constante:
la fase común al arreglo se estima a posteriori, pero su transferencia mediante
vecinos espaciales o temporales no muestra una ventaja clara en dc01. El cierre
es operacional respecto de esas preguntas, no una identificación definitiva
del origen de la fase ni una prueba de impredecibilidad universal.

El siguiente bloque es S4.7a, seguido del estudio residual por trayectoria
S4.7b y la comparación de representaciones S5:
\begin{equation}
\begin{aligned}
\mathrm{S4.7a}&\longrightarrow\mathrm{S4.7b}
\longrightarrow\mathrm{S5}\;\text{(representaciones)},\\
\mathrm{S6}&:\;\Delta H\text{ y observables invariantes con movimiento conocido},\\
\mathrm{S7}&:\;\text{operadores Wi-Fi y LoRa}.
\end{aligned}
\end{equation}
La preparación mecánica S6 puede avanzar una vez fijadas las convenciones de
fase, la calibración basal y el observable; no depende de completar todas las
arquitecturas de S5. El programa conserva la incorporación gradual de
complejidad, guiada por una limitación experimental concreta.

La actualización del programa se concreta en la \cref{tab:scenario_program_map}:
S4.7 avanza junto con diseño de referencia y Esc. E0; S6 se materializa en reflector
y fantoma E1--E2, y S7 compara operadores sobre esas perturbaciones. G0--G4 evaluará
la contribución geométrica en datos propios. E3--E5 añaden personas y transferencia
cuando exista evidencia suficiente. Ninguno de estos identificadores amplía el conjunto
de resultados ejecutados que se reporta en este capítulo.

\chapter{Contribuciones en desarrollo y alcance científico}
\label{ch:contributions}

\section{Aportes experimentales documentados sobre la fase}

El aporte actual es demostrar que una corrección de retardo parcialmente
predecible mejora la concordancia de fase del canal. La predicción completa de
fase absoluta permanece abierta; esta limitación no elimina el resultado
positivo de la corrección espectral. La \cref{sec:synthesis} reúne las cifras,
los protocolos y las hipótesis que sostienen esta distinción.

\begin{enumerate}
\item \textbf{Descomposición y transferencia de una corrección de baja dimensión.}
S4.2 distingue retardo efectivo e interceptos por enlace, encuentra acuerdo
parcial entre modelos y transfiere la corrección entre antenas y frecuencias.
La alineación oráculo define un límite diagnóstico; la transferencia estricta
entre frecuencias demuestra reutilización de los parámetros sin reajuste en
las frecuencias reservadas.
\item \textbf{Predicción de retardo con utilidad sobre el canal complejo.}
S4.4 identifica información predictiva en los descriptores RT y separa
interpolación densa de evaluación regional. S4.5 reduce el error angular de
81.42 a 45.59 grados en el protocolo denso y de 81.49 a 64.08 grados por
bloques, respecto de las referencias con CPO. La coherencia pasa de 0.156 a
0.634 y de 0.157 a 0.448, respectivamente. El retardo se predice sin el canal
objetivo, aunque el desfase constante continúa ajustándose con ese canal.
\item \textbf{Caracterización del error que permanece.}
S4.3 documenta estructura espectral y asociaciones con la energía multitrayecto;
S4.6 distingue una dirección común estimable en la captura cuya copia desde
vecinos espaciales o temporales no aporta una ventaja clara. La trazabilidad
temporal recuperada permite evaluar esa hipótesis con tiempo físico. Estos
hallazgos motivan invariantes de fase y modelado residual, sin identificar
una causa instrumental única ni demostrar independencia estadística.
\end{enumerate}

B0--B2 y S4.1 se sitúan como reproducción y adaptación de la calibración
diferenciable y del tratamiento de incertidumbre de fase
\citep{Hoydis2024DiffRT,Ruah2024Phase}. Las extensiones S4.2--S4.6 aportan
la caracterización y la evaluación anteriores sobre dc01. Su alcance de
novedad doctoral deberá contrastarse con métodos de referencia bajo protocolos
comunes; la ejecución de un diagnóstico no constituye por sí sola una
prioridad científica universal.

\section{Contribuciones doctorales por completar}

Sobre esta base, se mantienen las siguientes contribuciones doctorales como objetivos por completar:
\begin{enumerate}
\item una definición de gemelo inalámbrico orientado al sensado, basada en fidelidad de campo y fidelidad diferencial;
\item un banco de pruebas común para representaciones físicas y neuronales del entorno bajo regímenes de información explícitos;
\item una interfaz de trayectorias que desacople propagación de observación Wi-Fi/LoRa;
\item una representación factorizada del entorno, la estructura de fase, el ser humano, la fisiología, la radio y las perturbaciones instrumentales;
\item un marco de observabilidad basado en Fisher para determinar zonas y variables recuperables;
\item un procedimiento de transferencia de simulación a mediciones reales, evaluado mediante la eficiencia de datos y la generalización fuera de distribución;
\item una formulación del diseño de red orientada al sensado mediante criterios de diseño óptimo de experimentos;
\item un conjunto de datos multimodal propio que vincule la señal compleja de radiofrecuencia, el estado humano y referencias sincronizadas.
\end{enumerate}

La contribución principal no se define por lograr la mayor exactitud en un conjunto de datos particular, sino por establecer una metodología transferible para responder tres preguntas: \textit{qué información física existe}, \textit{qué representación la conserva} y \textit{qué configuración de infraestructura permite recuperarla de forma confiable}.

\section{Criterio de contribución de la actualización metodológica}

La propuesta actual concentra la contribución por demostrar en la relación entre
invariancia instrumental y sensibilidad física. El artefacto central será un banco
que compare ramas robusta y coherente bajo intervenciones conocidas, con geometría,
referencia y observación declaradas. Se exigirá mostrar qué información conserva cada
rama y en qué condiciones su combinación mejora transferencia o incertidumbre.
La mera combinación de RT, LiDAR y una red neuronal no bastará como novedad.

Tres contrastes harán revisable esa aportación: G0--G4 con ablación material para medir
el efecto geométrico; E0--E2 para separar repetibilidad de respuesta diferencial; y
E3--E5 para contrastar si ese conocimiento se transfiere a variables humanas y otros
dominios. El conjunto propio incluirá también condiciones nulas, fallos y límites
de referencia. La evidencia de dc01 sustenta la motivación y los protocolos de fase;
la validación de estas nuevas contribuciones permanece pendiente.

\chapter{Propuesta científica y correspondencia con el proyecto de Ciencia de Frontera}
\label{ch:secihti_mapping}

\section{Problema, hipótesis y alcance de la propuesta actualizada}

La propuesta aborda la dificultad de recuperar movimiento y fisiología mediante
infraestructura inalámbrica cuando la señal útil se superpone con multitrayectoria,
deriva instrumental y cambios del entorno. Los avances de dc01 muestran que mejorar
potencia y retardos no garantiza predecir fase bruta, y que una corrección espectral
parcialmente predecible sí aporta concordancia bajo las condiciones documentadas.
El siguiente problema científico consiste en determinar qué transformaciones pueden
compensarse o hacerse invariantes sin perder la respuesta física de interés.

Se propone un gemelo digital con geometría métrica, modelo de observación y dos ramas
complementarias: una robusta para estructura relativa y otra coherente para dinámica.
La hipótesis central y H1--H6 se formalizan en la \cref{sec:updated_hypotheses}.
El contraste requiere intervenir físicamente el entorno y caracterizar la referencia
de adquisición; no se resuelve únicamente añadiendo capacidad neuronal al ajuste
estático. La aportación esperada es un procedimiento reproducible que relacione
calibración, representación, sensibilidad y transferencia con incertidumbre explícita.

El programa conserva Wi-Fi y LoRa como radios de interés. SDR, radar, visión y
sensores de referencia se emplearán para validar y supervisar, declarando qué modalidades
son necesarias durante inferencia. La aplicación humana se desarrollará después del
reflector y el fantoma. El trabajo de laboratorio y la demostración comunitaria tendrán
alcances propios; no se atribuye validación clínica ni despliegue formal a la Etapa 1.

\section{Correspondencia científica y administrativa}

El programa aprobado se organiza en diseño, implementación experimental e integración
científica. La actualización precisa la secuencia técnica y los criterios de evidencia;
no constituye por sí misma una modificación del presupuesto o de compromisos autorizados.

\begin{table}[htbp]
\centering\small
\caption{Correspondencia entre etapas del proyecto y evidencia del programa actualizado.}
\label{tab:proposal_updated_mapping}
\begin{tabularx}{\textwidth}{P{3.0cm}L L}
\toprule
Etapa & Trabajo doctoral & Evidencia y entregables\\
\midrule
Diseño & L1/L2/L4; S4.7a, referencia, G0--G4 y Esc. E0--E1. & Especificación geométrica,
protocolo, comparaciones estáticas y caracterización de repetibilidad/sensibilidad.\\
Implementación & L3/L4/L5; Esc. E2--E3 y operadores multirradio. & Fantoma, referencia
sincronizada y estudio humano autorizado; datos con procedencia y particiones.\\
Integración & L5/L6; Esc. E4--E5, inversión y diseño de red. & Límites de separabilidad,
transferencia, incertidumbre y preparación del piloto según madurez demostrada.\\
\bottomrule
\end{tabularx}
\end{table}

Las etapas no implican que toda actividad de una fila deba concluir en los primeros
diez meses. El cronograma del \cref{ch:wbs_stage1} distingue preparación, prueba
preliminar y validación posterior. En particular, varios usuarios y transferencia
externa se condicionarán por la evidencia acumulada y los acuerdos necesarios.

\section{Productos científicos y criterios de revisión}

El primer producto será un banco comparativo que conserve resultados brutos, relativos
y diagnósticos oráculo separados, con particiones espaciales reproducibles. El segundo
será un protocolo metrológico que registre escena, antenas, relojes y desplazamiento
local, y evalúe G0--G4 sin confundir geometría con ajuste material. El tercero será
un conjunto preliminar de mediciones con reflector y fantoma, seguido de datos humanos
cuando se cumplan las condiciones técnicas y éticas. El cuarto será la comparación
de ramas y métodos de referencia bajo presupuestos de información explícitos.

Cada producto deberá contener la pregunta contrastada, configuración, artefactos de
reproducción, distribución de errores, incertidumbre y limitaciones. La disponibilidad
de una malla, un modelo entrenado o una demostración visual no basta para declarar
validado el producto. El criterio científico de revisión será que otro grupo pueda
reconstruir la correspondencia entre estado físico, señal adquirida, predicción y
resultado evaluado.

\section{Viabilidad y decisiones ante evidencia adversa}

La viabilidad se apoya en reutilizar el equipo reportado, priorizar referencia RF y
movimiento medido, y ejecutar S4.7a y el diseño del banco en paralelo. La inversión
en reconstrucción LiDAR se justificará por la comparación geométrica y la utilidad
multiescena; su adquisición no condiciona iniciar E0 ni una prueba local de E1 con
geometría simple verificada. La selección comercial se decidirá mediante especificación,
cotización y compatibilidad dentro de los rubros autorizados.

Una invariancia demasiado amplia se restringirá o se reservará a la rama de contexto.
Una referencia inestable motivará recalibración instrumental. Si el gemelo no reproduce
la sensibilidad mecánica, se investigarán geometría local, mecanismos de propagación
y observación antes de usarlo para inversión humana. La ausencia de separabilidad
motivará abstención o rediseño de enlaces. Estas decisiones proporcionan resultados
científicos útiles aun cuando una hipótesis de mejora no se confirme.

La colaboración prospectiva con UPNA y el Dr. Francisco Falcone se orientará a repetir
el protocolo canónico en otras bandas, sujeto a acuerdos, acceso y financiación.
La vinculación con Colmenas Guadalajara conserva su función de caracterización de sitio
y divulgación. Ninguna de estas actividades se presenta como ya ejecutada o como
sustituto de una validación experimental independiente.

\chapter{Plan de trabajo de la Etapa 1}
\label{ch:wbs_stage1}

\section{Alcance y criterio rector de la Etapa 1}

El plan de esta etapa se conserva como marco de actividades y entregables. Su avance científico actual se documenta en el \cref{ch:avances}: el banco dc01, B0--B2 y S4.1--S4.6d.2 desarrollan LT1 y el marco de evaluación de LT2, incluida la validación regional del predictor de retardo. Estos avances no implican que se hayan completado la instrumentación multimodal, el estudio con participantes o los hitos de vinculación comunitaria.

La Etapa 1 comprende los primeros diez meses del proyecto y corresponde a la fase de
\emph{modelado, diseño experimental y vinculación comunitaria}. Su propósito es
establecer las bases conceptuales, experimentales, metodológicas y éticas necesarias para
estudiar la recuperación de información fisiológica y de comportamiento a partir de
señales inalámbricas en escenarios reales. En consecuencia, la etapa articula seis tareas:
(i) simular escenarios Wi-Fi y LoRa; (ii) caracterizar las limitaciones físicas e
instrumentales; (iii) definir el conjunto de datos multimodal; (iv) formalizar el protocolo
de adquisición; (v) establecer métodos de referencia con simulación, datos abiertos y
pruebas preliminares; y (vi) seleccionar y caracterizar el sitio piloto mediante actividades
de vinculación.

Antes del despliegue comunitario se realizará una campaña experimental controlada en
laboratorio. Esta campaña constituye el vínculo empírico entre la simulación y el piloto:
permite verificar la instrumentación, la sincronización, la repetibilidad, la observabilidad
y los procedimientos de calidad. No sustituye la adquisición formal prevista para la
Etapa 2 ni pretende anticipar sus resultados; reduce la incertidumbre técnica antes de
trasladar el protocolo a una Colmena u otro espacio comunitario de Guadalajara.

El éxito de la Etapa 1 no se juzgará por el volumen de datos ni por el número de
participantes, sino por la reducción verificable de la incertidumbre metodológica. Al
finalizar el mes 10 deberán estar disponibles un modelo físico reproducible, una plataforma
de adquisición estable, un estudio preliminar de laboratorio documentado, un protocolo
validado y un sitio comunitario caracterizado para la campaña formal de la Etapa 2. Esta
secuencia mantiene la continuidad con la metodología general: primero se determina qué
debería ser observable, después qué puede medirse y, por último, bajo qué condiciones es
justificado avanzar al entorno comunitario.

\section{Organización de las líneas de trabajo}

La Etapa 1 se organiza en seis líneas de trabajo (LT). La organización conserva los
objetivos E1-1--E1-6 del proyecto y los agrupa de acuerdo con sus dependencias científicas
y operativas. Las líneas no son compartimentos independientes: el modelado fundamenta
el protocolo; el protocolo determina la instrumentación; la instrumentación hace posible
el estudio preliminar; y la evidencia reunida condiciona la transición a la Etapa 2.

\begin{table}[H]
\centering
\caption{Líneas de trabajo de la Etapa 1 y correspondencia con los objetivos específicos.}
\label{tab:wbs_stage1_summary}
\small
\begin{tabularx}{\textwidth}{P{1.15cm}P{3.25cm}P{1.9cm}L}
\toprule
\textbf{Línea} & \textbf{Denominación} & \textbf{Objetivos} & \textbf{Resultado esperado}\\
\midrule
LT1 & Modelado físico y gemelo inalámbrico inicial
& E1-1, E1-2
& Entorno reproducible de simulación Wi-Fi/LoRa, comparación con mediciones reales y métricas de propagación, fase y observabilidad.\\

LT2 & Marco de datos, protocolo y métodos de referencia
& E1-3, E1-4, E1-5
& Especificación multimodal, protocolo v1.0, esquema de metadatos, control de calidad y métodos reproducibles de referencia.\\

LT3 & Plataforma experimental multimodal de laboratorio
& E1-2, E1-4
& Banco experimental integrado y caracterizado: Wi-Fi, LoRa, referencias mmWave/visión/fisiología y sincronización.\\

LT4 & Estudio experimental preliminar y controlado
& E1-2, E1-5
& Conjunto preliminar de laboratorio y validación de presencia, movimiento, micromovimiento y respiración en condiciones controladas.\\

LT5 & Vinculación comunitaria, Colmenas y taller de IoT
& E1-6
& Sitio piloto caracterizado, caso de uso comunitario seleccionado y taller práctico ejecutado en Guadalajara.\\

LT6 & Integración, evaluación y transición a la Etapa 2
& E1-3--E1-6
& Informe técnico de Etapa 1, paquete de transición, matriz de riesgos y protocolo para despliegue formal.\\
\bottomrule
\end{tabularx}
\end{table}

\subsection{LT1 --- Modelado físico y gemelo inalámbrico inicial}
\label{sec:wp1_stage1}

\textbf{Periodo:} M1--M6; cierre inicial S4.7a en M1--M2 y ampliación geométrica en M4--M6.

\textbf{Avance documentado:} S1 y B0--B2 cuantifican la discrepancia estática; S4.1--S4.6d.2 analizan fase, estructura residual y transferencia del retardo al canal. S4.6 caracteriza la fase común, recupera el tiempo de adquisición y evalúa continuidad local. El trabajo siguiente S4.7a estudiará fidelidad relativa; S4.7b abordará incertidumbre por trayectoria. Esto aporta evidencia a P1.2 y P1.3. P1.1 y la revisión R1 requieren vincular los registros reproducibles indicados en la \cref{sec:trazabilidad_avances}.

\textbf{Pregunta rectora:} ¿qué información sobre entorno, movimiento y micromovimiento
puede preservar el canal Wi-Fi/LoRa bajo geometrías y configuraciones de radio
controladas?

\paragraph{Actividades.}
\begin{enumerate}[label=\textbf{LT1.\arabic*:},leftmargin=3.8em]
    \item Construcción y validación de casos de referencia: espacio libre, reflexión,
    multitrayectoria y perturbaciones geométricas controladas.
    \item Reproducción del banco de pruebas DICHASUS dc01 y generación del conjunto
    $(\vect p_i,H_i^{\rm real},\mathcal P_i^{\rm RT})$.
    \item Evaluación del trazado de rayos nominal y calibración progresiva de materiales,
    dispersión y errores locales de fase.
    \item Implementación del modelo de observación Wi-Fi OFDM y del modelo inicial
    LoRa CSS/IQ.
    \item Definición de métricas de fidelidad de campo y sensibilidad diferencial,
    incluyendo NMSE complejo, error de fase, PDP, dispersión cuadrática media del retardo y
    $\partial H/\partial q$.
    \item Construcción de mapas preliminares de cobertura y observabilidad en
    escenarios de referencia en interiores y exteriores.
\end{enumerate}

\paragraph{Productos y criterio de aceptación.}
\begin{itemize}
    \item \textbf{P1.1} Repositorio reproducible del modelo directo.
    \item \textbf{P1.2} Comparación DICHASUS entre mediciones y simulación.
    \item \textbf{P1.3} Documento técnico de limitaciones físicas e instrumentales.
    \item \textbf{Revisión R1:} las simulaciones son reproducibles y la discrepancia
    entre simulación y medición se encuentra cuantificada mediante métricas definidas de antemano.
\end{itemize}

\subsection{LT2 --- Marco de datos, protocolo y métodos de referencia}
\label{sec:wp2_stage1}

\textbf{Periodo:} M1--M5.

\textbf{Pregunta rectora:} ¿qué información debe conservarse y bajo qué reglas debe
adquirirse para que las mediciones sean comparables, trazables y reutilizables?

\paragraph{Actividades.}
\begin{enumerate}[label=\textbf{LT2.\arabic*:},leftmargin=3.8em]
    \item Revisión y reproducción de flujos de procesamiento sobre conjuntos de datos públicos de canal,
    actividad y fisiología.
    \item Definición del modelo de datos: señal de radiofrecuencia, valores de referencia, estado del participante,
    contexto, geometría, configuración de radio y metadatos.
    \item Especificación de nomenclatura, control de versiones, marcas temporales, estructura de
    archivos, trazabilidad y criterios de inclusión/exclusión.
    \item Diseño del protocolo multimodal para configuraciones sin participantes y
    con uno o varios participantes.
    \item Definición de procedimientos de sincronización y control de calidad.
    \item Formalización de consentimiento, anonimización, gestión de video,
    acceso y gobernanza de datos.
    \item Implementación de métodos de referencia sin selección oráculo de canales,
    subportadoras o componentes.
\end{enumerate}

\paragraph{Productos y criterio de aceptación.}
\begin{itemize}
    \item \textbf{P2.1} Diccionario y esquema del conjunto de datos multimodal.
    \item \textbf{P2.2} Protocolo de adquisición v1.0.
    \item \textbf{P2.3} Lineamientos de ética, privacidad y gobernanza.
    \item \textbf{P2.4} Métodos reproducibles de referencia sobre datos públicos.
    \item \textbf{Revisión R2:} no se iniciará una campaña formal con participantes sin
    protocolo, trazabilidad, control de calidad y autorización ética aplicable.
\end{itemize}

\subsection{LT3 --- Plataforma experimental multimodal de laboratorio}
\label{sec:wp3_stage1}

\textbf{Periodo:} M1--M6; diseño y adquisición M1--M3, caracterización E0 en M2--M4.

\textbf{Objetivo:} implementar una plataforma común que permita observar el mismo
fenómeno físico mediante radios de interés y modalidades de referencia.

La arquitectura experimental se define como

\begin{equation}
\begin{aligned}
\mathcal A_{\rm lab}=\{&Y_{\rm WiFi},Y_{\rm LoRa},Y_{\rm mmWave},Y_{\rm ref},\\
&V_{\rm RGB(D)},x_{\rm stage},x_{\rm resp},x_{\rm ECG/RR},x_{\rm IMU}\}.
\end{aligned}
\label{eq:lab_architecture}
\end{equation}

Esta arquitectura es un catálogo de modalidades posibles; cada captura declarará el
subconjunto realmente disponible. $Y_{\rm ref}$ y $x_{\rm stage}$ corresponden a la
referencia RF y al desplazamiento mecánico medido. ECG e inerciales son auxiliares
según tarea y no requisitos para cerrar el banco mecánico.

\paragraph{Bloques instrumentales.}
\begin{description}[style=nextline]
    \item[\textbf{Referencia RF y mecánica.}]
    Adquisición coherente caracterizada, etapa con medición independiente y soportes
    registrados. La incertidumbre del banco determina los pasos utilizables en E1.
    \item[\textbf{Wi-Fi.}]
    Enlaces configurables para CSI, con control de canal, ancho de banda,
    geometría Tx--Rx, orientación y tasa de paquetes.

    \item[\textbf{LoRa/LoRaWAN.}]
    Nodos y puerta de enlace de laboratorio; cuando sea posible, observación I/Q mediante
    SDR para preservar información anterior a la reducción a RSSI/SNR.

    \item[\textbf{Radar de ondas milimétricas.}]
    Modalidad de radiofrecuencia de referencia para movimiento y micromovimiento.

    \item[\textbf{RGB/RGB-D.}]
    Referencia geométrica para posición, orientación, trayectoria y actividad.

    \item[\textbf{Fisiología e inerciales.}]
    Banda de expansión respiratoria, ECG o banda torácica con intervalos RR,
    y una o más IMU para movimiento corporal.
\end{description}

\paragraph{Actividades de caracterización.}
Se medirán estabilidad temporal, piso de ruido, pérdida de paquetes, CFO/SFO/CPO,
latencia entre modalidades y deriva temporal. La sincronización se expresará como

\begin{equation}
a_m t_i^{(m)}+b_m=t_i^{\rm global}+\epsilon_i^{(m)},
\end{equation}

donde $a_m,b_m$ corrigen escala y origen del reloj. El error residual
$\epsilon_i^{(m)}$ deberá estimarse y reportarse, no asumirse nulo; se segmentarán
las sesiones si los reinicios o la deriva invalidan un único mapa afín.

\paragraph{Productos y criterio de aceptación.}
\begin{itemize}
    \item \textbf{P3.1} Banco experimental inventariado y documentado.
    \item \textbf{P3.2} Procedimiento de calibración y sincronización.
    \item \textbf{P3.3} Informe de estabilidad y repetibilidad.
    \item \textbf{Revisión R3:} las modalidades del banco mínimo E0--E1 producen
    registros trazables, con referencia y sincronización caracterizadas. La integración
    adicional continuará según el cronograma y se validará antes de utilizarla.
\end{itemize}

\subsection{LT4 --- Estudio experimental preliminar y controlado}
\label{sec:wp4_stage1}

\textbf{Periodo:} M3--M9; reflector M3--M5, fantoma M5--M7 y una persona M7--M9.

El estudio preliminar se ejecutará de forma acumulativa. Cada nivel habilita el siguiente; no se
introducirán participantes antes de validar primero el comportamiento instrumental y
geométrico.

\begin{table}[H]
\centering
\caption{Matriz experimental mínima del estudio preliminar de laboratorio.}
\label{tab:lab_prepilot}
\small
\begin{tabularx}{\textwidth}{P{1.1cm}P{3.0cm}L L}
\toprule
\textbf{Nivel} & \textbf{Condición} & \textbf{Variable controlada} & \textbf{Resultado buscado}\\
\midrule
N0 & Habitación vacía & Tiempo, canal, dispositivo & Piso de ruido, estabilidad y perturbaciones instrumentales.\\
N1 & Reflector estático & Posición y material & Sensibilidad espacial y repetibilidad.\\
N2 & Reflector móvil & Desplazamiento conocido & Fidelidad diferencial y $\partial H/\partial x$.\\
N3 & Fantoma respiratorio & Amplitud, frecuencia y movimiento lento & Forma de onda mecánica y sensibilidad reproducibles.\\
N4 & Una persona & Postura, actividad y respiración & Movimiento y forma de onda frente a referencias sincronizadas.\\
N5 & Wi-Fi y LoRa pareados & Misma perturbación física & Comparación de sensibilidad y observabilidad entre tecnologías.\\
N6 & Dos personas, exploratorio & Separación y actividad & Viabilidad y condicionamiento antes de la Etapa 2.\\
\bottomrule
\end{tabularx}
\end{table}

El conjunto de salida se define como

\begin{equation}
\begin{aligned}
\mathcal D_{\rm lab}^{(0)}=\{&Y_{\rm WiFi},Y_{\rm LoRa},Y_{\rm mmWave},Y_{\rm ref},\\
&V_{\rm RGB(D)},x_{\rm stage},x_{\rm resp},x_{\rm ECG/RR},x_{\rm IMU},M\},
\end{aligned}
\end{equation}

donde $M$ contiene metadatos y estado de control de calidad.

\paragraph{Revisión R4.}
En M7 se revisará la evidencia de E0--E2: estabilidad instrumental, respuesta a
movimiento medido y recuperación del fantoma con incertidumbre. La adquisición de
una persona se prevé en M7--M9 y requiere además aprobación ética. El cumplimiento
del hito mecánico no equivale a validación fisiológica ni habilita automáticamente
un despliegue humano comunitario.

\subsection{LT5 --- Vinculación comunitaria, Colmenas Guadalajara y taller de IoT}
\label{sec:wp5_stage1}

\textbf{Periodo:} M6--M10; preparación institucional desde M6 y taller previsto en M9--M10.

\textbf{Objetivo:} seleccionar y caracterizar un sitio comunitario en Guadalajara y
establecer una relación de trabajo que combine preparación del piloto con apropiación
tecnológica.

La vinculación seguirá cuatro fases:
\begin{enumerate}[label=\textbf{V\arabic*:},leftmargin=3em]
    \item \textbf{Vinculación institucional:} identificación de contraparte local,
    condiciones de acceso, responsables y logística.
    \item \textbf{Caracterización técnica:} geometría, materiales, energía,
    infraestructura Wi-Fi, cobertura LoRa, ruido RF, zonas de actividad y restricciones
    de instalación.
    \item \textbf{Codiseño comunitario:} selección de un único caso de uso ambiental
    para el nodo IoT de divulgación.
    \item \textbf{Taller y preparación del despliegue:} taller práctico, demostración de sensores,
    conectividad y energía autónoma, y pruebas técnicas sin convertir el taller en una
    campaña experimental de sujetos humanos.
\end{enumerate}

El resultado técnico de la visita será una ficha de sitio y un primer modelo
$G_E^{\rm Colmena}$ utilizable posteriormente para simulación y planificación de la
Etapa 2.

\paragraph{Revisión R5.}
El sitio deberá contar con contraparte identificada, caracterización técnica, restricciones
operativas documentadas y un caso de uso comunitario definido.

\subsection{LT6 --- Integración y transición a la Etapa 2}
\label{sec:wp6_stage1}

\textbf{Periodo:} M8--M10.

La línea final consolida los resultados de las cinco líneas anteriores y establece la primera
versión operativa del sistema.

\paragraph{Productos de cierre.}
\begin{enumerate}[label=\textbf{P6.\arabic*:},leftmargin=3em]
    \item Entorno de simulación reproducible Wi-Fi/LoRa.
    \item Plataforma multimodal de laboratorio caracterizada.
    \item Protocolo de adquisición, metadatos y control de calidad v1.0.
    \item Conjunto preliminar $\mathcal D_{\rm lab}^{(0)}$.
    \item Métodos reproducibles de referencia.
    \item Sitio comunitario caracterizado.
    \item Taller y materiales de vinculación ejecutados.
    \item Matriz de riesgos y pendientes.
    \item Plan cerrado de despliegue y adquisición para la Etapa 2.
\end{enumerate}

La transición se formaliza mediante

\begin{equation}
R_{\rm E2}=
R_{\rm RF}\land
R_{\rm sinc}\land
R_{\rm protocolo}\land
R_{\rm etica}\land
R_{\rm sitio}\land
R_{\rm datos}.
\label{eq:stage2_gate}
\end{equation}

La Etapa 2 deberá iniciar su campaña formal cuando estas condiciones hayan sido
verificadas, y no únicamente por cumplimiento cronológico del mes 10.

\section{Secuencia y dependencias}

Las dependencias principales siguen una secuencia de evidencia. LT1 establece las
predicciones físicas que deben contrastarse; LT2 determina cómo adquirir y conservar las
mediciones; LT3 materializa esas decisiones en una plataforma común; y LT4 contrasta el
modelo con perturbaciones controladas. LT5 se inicia cuando existe suficiente madurez
técnica para caracterizar el sitio comunitario, mientras que LT6 integra la evidencia y
documenta los supuestos pendientes. El solapamiento temporal mostrado en la
\cref{tab:cronograma_etapa1} expresa trabajo concurrente, no independencia entre líneas.

\section{Cronograma de la Etapa 1}

\subsection{Paquete técnico de las primeras ocho semanas}
\label{sec:eight_week_package}

Dentro de M1--M2 se propone un paquete corto que convierte la revisión de literatura
en artefactos verificables antes de cerrar especificaciones de compra:
\begin{table}[htbp]
\centering\small
\caption{Secuencia técnica inmediata; las semanas son relativas al inicio del paquete.}
\label{tab:eight_week_package}
\begin{tabularx}{\textwidth}{P{1.5cm}L L}
\toprule
Semana & Actividad & Evidencia de cierre\\
\midrule
1 & Congelar versiones, manifiestos, particiones y artefactos S1--S4.6. &
Correspondencia reproducible y sumas de verificación.\\
2 & Implementar $\mathsf V_0$--$\mathsf V_7$ y $\mathsf T_1$--$\mathsf T_7$. &
Pruebas algebraicas y matriz de respuesta a perturbaciones.\\
3 & Ejecutar fidelidad estática de dc01. & Resultados densos y por bloques con
normalización ajustada sólo en entrenamiento.\\
4 & Auditar d010 y realizar pruebas de sensibilidad si sus metadatos lo permiten. &
Pareto entre rechazo de perturbaciones y cambio geométrico del transmisor.\\
5 & Implementar S4.7b. & Baselines circulares, NLL, calibración y abstención.\\
6 & Integrar perturbaciones seleccionadas de RadioRange. & Vulnerabilidades por
observable bajo rangos trazables.\\
7 & Ejecutar un caso mínimo de Indoor Radio Mapping y WiSegRT. & Flujo de nube/malla
a RT; RSSI y canal simulado con límites declarados.\\
8 & Auditar subconjuntos CAEZ y OpenLoRa. & Adaptadores iniciales y requisitos
medidos para el banco propio.\\
\bottomrule
\end{tabularx}
\end{table}

El paquete no condiciona toda adquisición a obtener un método ganador. Debe responder
qué fase necesita referencia, qué ancho de banda y coherencia preservan sensibilidad,
qué precisión geométrica resulta relevante y qué salidas proporciona cada radio.
Las compras de baja ambigüedad necesarias para mecánica, montaje y trazabilidad pueden
prepararse en paralelo; las especificaciones de recepción coherente se cerrarán con
la evidencia disponible y las cotizaciones compatibles.

\begin{table}[htbp]
\centering\small
\caption{Cronograma revisado: meses relativos al inicio operativo de la etapa; ventanas propuestas sujetas a evidencia.}
\label{tab:cronograma_etapa1}
\begin{tabularx}{\textwidth}{P{1.8cm}P{2.2cm}L L}
\toprule
Meses & Líneas & Actividad & Producto revisable\\
\midrule
M1--M2 & LT1/LT2 & S4.7a en dc01. & Comparación $\mathsf V_i$ y controles espaciales.\\
M1--M3 & LT2/LT3 & Diseño del banco y compras prioritarias. & Referencia RF, metrología y captura especificadas.\\
M2--M4 & LT3 & Esc. E0. & Deriva y repetibilidad documentadas.\\
M3--M5 & LT3/LT4 & Esc. E1, reflector. & $\Delta H$ y Jacobiano con desplazamiento medido.\\
M4--M6 & LT1/LT3 & Reconstrucción RGB-D/LiDAR. & G0--G4, registro y controles independientes.\\
M5--M7 & LT4 & Esc. E2, fantoma. & Fidelidad diferencial y forma de onda.\\
M6--M8 & LT1/LT4 & Wi-Fi, LoRa y radar comparativos. & Operadores y referencias bajo perturbación común.\\
M7--M9 & LT4 & Esc. E3 con aprobación ética. & Datos humanos preliminares y límites.\\
M8--M10 & LT6 & Integración y transferencia inicial E5. & Protocolo y conjunto v1; pendientes de Etapa 2.\\
M6--M10 & LT5 & Vinculación; taller en M9--M10. & Sitio caracterizado y demostración de divulgación.\\
\bottomrule
\end{tabularx}
\end{table}

El calendario expresa dependencias técnicas y no acredita ejecución por transcurso
del tiempo. Esc. E4 y transferencia externa amplia se incorporarán según madurez y
recursos; no se prometen como resultados humanos completos dentro de esta ventana.
Mientras se integra el equipo se puede avanzar con CAEZ, el flujo geometría--RSSI,
OpenLoRa, reconstrucción RGB-D, marcas temporales y diseño de referencia. Cada conjunto
adicional se someterá primero a una auditoría de observables y procedencia.

\section{Movilidad científica}

El presupuesto aprobado considera una bolsa de \$68,000 MXN para viáticos y
\$30,000 MXN para pasajes nacionales, originalmente vinculados a actividades de
divulgación y vinculación. La planeación operativa propone concentrar estas movilidades
en una ventana M9--M10, de modo que una misma visita contribuya a: (i) revisión
experimental; (ii) validación Wi-Fi/LoRa; (iii) caracterización de la Colmena;
(iv) impartición del taller; y (v) preparación de la Etapa 2.

La solicitud aprobada refiere hospedaje y alimentación para dos personas por 12 noches y
tres vuelos nacionales redondos. Por tanto, una estancia presencial completa de un mes
para el especialista deberá considerarse una extensión deseable, pero su ejecución requerirá
validar previamente con la administración del proyecto la posibilidad de redistribución del
rubro o complementarla con recursos institucionales. El documento operativo no debe
suponer dicha ampliación como autorizada.

La movilidad prospectiva a UPNA constituye una ampliación científica distinta de los
pasajes nacionales presupuestados. Requerirá acuerdo con el grupo anfitrión, disponibilidad
de instrumentación y financiación específica; su objetivo será repetir E1/E2 con el
protocolo ya validado localmente. Una propuesta de colaboración no se registra como
estancia concedida ni como validación realizada.

\section{Plan de adquisiciones}

El presupuesto aprobado contiene tres bolsas materiales diferenciadas:
\$40,000 MXN para dispositivos electrónicos de radio configurables y accesorios;
\$40,000 MXN para sensores y suministros de laboratorio de referencia; y
\$70,000 MXN para 25 kits electrónicos de divulgación.


\subsection{Inventario reportado y verificación funcional}

La planeación considera como inventario reportado plataformas Wi-Fi/ESP32, nodos
LoRa STM32WL, equipos Notecard y RAK2287, un ADALM-Pluto, tres IWRL6432BOOST y
cámaras RGB, RGB-D o estéreo. Antes de comprar se registrarán modelos exactos,
interfaces y estado funcional. La denominación Notecard no basta para identificar
tecnología o acceso a señal; se comprobará el SKU. Tampoco se asumirá que un
concentrador LoRaWAN entrega I/Q coherente.

Los IWRL6432BOOST se basan en la familia radar de 57--64 GHz
\citep{TIIWRL6432BOOST}. Su disponibilidad reportada permite plantear vistas complementarias;
se verificarán firmware, acceso a datos, calibración y sincronización antes de fijar
el protocolo. Su precisión de desplazamiento deberá medirse y no se adoptará como
referencia submilimétrica independiente. Se reutilizarán estas unidades antes de
priorizar otro radar equivalente.

\subsection{Plataforma experimental de radiofrecuencia --- bolsa autorizada: \$40,000 MXN}

\begin{table}[htbp]
\centering\small
\caption{Prioridades RF revisadas a partir del inventario reportado.}
\label{tab:rf_procurement}
\begin{tabularx}{\textwidth}{P{3.4cm}L P{2.9cm}}
\toprule
Bloque & Función y condición & Prioridad\\
\midrule
Referencia coherente & Resolver recepción simultánea o referencia verificable; comprobar
compatibilidad con SDR y reloj antes de seleccionar equipo. & Muy alta\\
Coaxiales, divisores, atenuadores y adaptadores & Caracterizar pérdidas/fase y estabilizar
cadenas; cantidades según arquitectura, sin asumir calibración de fábrica suficiente. & Muy alta\\
Antenas y soportes rígidos & Registro geométrico, polarización y repetibilidad del centro
de montaje; cables inmovilizados. & Alta\\
Sincronización y automatización & Marcas temporales, disparos disponibles y control
mecánico; medición de deriva entre modalidades. & Alta\\
Ampliación Wi-Fi/LoRa/SDR & Cubrir una carencia demostrada de CSI, I/Q o enlaces;
reutilizar primero plataformas existentes. & Condicionada\\
\bottomrule
\end{tabularx}
\end{table}

La referencia es un requisito funcional, no una compra cerrada: el Pluto documentado
no ofrece por sí solo dos receptores simultáneos. Se compararán soluciones compatibles
y su coste total dentro del rubro. Compartir frecuencia o PPS no exime de verificar
fase relativa tras reinicios.

\subsection{Metrología y referencias --- bolsa autorizada: \$40,000 MXN}

\begin{table}[htbp]
\centering\small
\caption{Prioridades para metrología y referencias del banco actualizado.}
\label{tab:reference_procurement}
\begin{tabularx}{\textwidth}{P{3.4cm}L P{2.9cm}}
\toprule
Instrumento & Uso científico & Prioridad\\
\midrule
Etapa lineal y referencia micrométrica & Desplazamiento medido, control de holgura,
reflector y fantoma; incertidumbre compatible con E1. & Muy alta\\
Fijaciones y marcadores & Registro reproducible de escena, sensores y antenas;
AprilTags u otros fiduciales con dimensiones verificadas. & Muy alta\\
Referencia respiratoria & Forma de onda sincronizada para E3; no sustituible por una
etiqueta de BPM o ECG aislado. & Alta\\
LiDAR 3D de gama media & Geometría global y repetibilidad multiescena; comparar G1--G3
frente a RGB-D existente. & Alta; cotizar\\
Distanciómetro o control métrico & Comprobación independiente de dimensiones; Bosch GLM
o equivalente es opcional, no requisito de marca. & Complementaria\\
Radar, RGB-D, IMU y ECG & Reutilizar disponibilidad; ampliar sólo ante una necesidad
medida. Cardiología queda fuera del núcleo E0--E2. & Condicionada\\
\bottomrule
\end{tabularx}
\end{table}

Para LiDAR se propone evaluar una clase Livox Mid-360 o equivalente, sin fijar compra,
precio ni exactitud no comprobados. La especificación exigirá exportación de nube de
puntos, marcas temporales accesibles, integración mediante SDK/ROS2 o interfaz compatible,
registro métrico y documentación de incertidumbre. No se requiere un equipo topográfico
ni se presupone que un escáner de mayor coste resolverá la fase del canal.

La selección monetaria se realizará contra cotizaciones y compatibilidad real. La
combinación completa de instrumentos no se da por financiada dentro de \$40,000 MXN.
Si excede el rubro, se priorizarán mecánica, referencia y sincronización, usando RGB-D
para iniciar el registro y gestionando préstamo o recursos complementarios para LiDAR.
No se trasladarán recursos entre bolsas ni se imputará movilidad internacional a pasajes
nacionales mediante esta actualización metodológica.

\subsection{Kits de divulgación --- bolsa autorizada: \$70,000 MXN}

El presupuesto aprobado corresponde a 25 kits, equivalente a un promedio de

\begin{equation}
\bar C_{\rm kit}=
\frac{70,000}{25}=2,800\ {\rm MXN/kit}.
\end{equation}

Cada kit se diseñará como plataforma modular con microcontrolador, conectividad,
sensores, alimentación y material de prototipado. El sensor específico del caso de uso
ambiental deberá seleccionarse después de la caracterización de la Colmena, evitando
adquirir anticipadamente 25 sensores especializados para un problema todavía no
seleccionado.

\section{Presupuesto consolidado de la Etapa 1}

\begin{table}[H]
\centering
\caption{Presupuesto aprobado y asignación por líneas de trabajo durante la Etapa 1.}
\label{tab:stage1_budget_formal}
\small
\begin{tabularx}{\textwidth}{L r P{2.1cm}}
\toprule
\textbf{Rubro} & \textbf{Monto MXN} & \textbf{Líneas principales}\\
\midrule
Apoyo a la investigación científica
& \$240,000 & LT1--LT6\\
Honorarios por servicios profesionales
& \$240,000 & LT1--LT6\\
Material eléctrico/electrónico para kits del taller
& \$70,000 & LT5\\
Viáticos
& \$68,000 & LT4--LT5\\
Pasajes nacionales
& \$30,000 & LT4--LT5\\
Material eléctrico/electrónico para pilotos
& \$40,000 & LT3--LT4\\
Sensores y suministros de laboratorio
& \$40,000 & LT3--LT4\\
\midrule
\textbf{Total Etapa 1}
& \textbf{\$728,000}
& \\
\bottomrule
\end{tabularx}
\end{table}

El cuadro anterior conserva íntegramente los rubros y montos aprobados. Las
subasignaciones internas de equipo se consideran planeación técnica y deberán cerrarse
contra cotizaciones y reglas administrativas antes de comprometer el gasto.

\section{Matriz de hitos y decisiones}

\begin{table}[H]
\centering
\caption{Hitos de decisión de la Etapa 1.}
\label{tab:stage1_milestones}
\small
\begin{tabularx}{\textwidth}{P{1.7cm}P{1.4cm}L L}
\toprule
\textbf{Revisión} & \textbf{Mes} & \textbf{Evidencia requerida} & \textbf{Decisión resultante}\\
\midrule
R1 & M4
& Gemelo estático y S4.7a reproducibles; discrepancia y acceso a fase declarados.
& Establecer el modelo físico de referencia.\\
R2 & M5
& Protocolo, metadatos, control de calidad y ética formalizados.
& Autorizar el estudio preliminar con participantes cuando corresponda.\\
R3 & M5
& Banco mínimo E0--E1 estable, referencia y sincronización caracterizadas.
& Continuar fantoma e integración multimodal, con control de deriva.\\
R4 & M7
& E0--E2: repetibilidad, sensibilidad mecánica y fantoma con incertidumbre.
& Iniciar E3 sólo con ética y referencias verificadas; preparar transferencia.\\
R5 & M9
& Colmena caracterizada, contraparte y caso de uso definidos.
& Preparar la instalación de la Etapa 2.\\
$R_{\rm E2}$ & M10
& Cumplimiento conjunto de radiofrecuencia, sincronización, protocolo, ética, sitio y datos.
& Iniciar campaña formal de Etapa 2.\\
\bottomrule
\end{tabularx}
\end{table}

\section{Síntesis}

La Etapa 1 constituye una fase de reducción de incertidumbre científica, no sólo de
cumplimiento administrativo. El modelado establece qué debería observarse; la
instrumentación verifica qué puede medirse; el estudio preliminar determina qué
información se conserva en condiciones reales; y la vinculación comunitaria procura que
el despliegue posterior ocurra en un sitio técnicamente caracterizado y con una contraparte
local. En conjunto, estas actividades convierten los fundamentos desarrollados en los
capítulos anteriores en una secuencia experimental auditable y preparan, sin anticiparla,
la validación de la Etapa 2.

\chapter{Conclusiones provisionales}
\label{ch:conclusions}

\textbf{La fase se mejoró mediante una corrección espectral parcialmente
predecible, aunque no se reconstruyó de forma completamente autónoma.}
El resultado central de S4 es distinguir una componente de retardo compartida
y transferible, una dirección común al arreglo condicionada a la captura y un
residual relativo con estructura espectral. Esta distinción permite interpretar
qué se ha logrado y elegir el tratamiento adecuado de lo que permanece abierto.


La primera etapa de investigación establece un marco reproducible de comparación entre
un canal medido y su representación electromagnética, y distingue fidelidad macroscópica,
fidelidad de magnitud, concordancia de fase y respuesta a perturbaciones físicas. La
cadena desde escena hasta observable permite separar calibración del entorno, tratamiento
instrumental y aprendizaje del problema inverso.

Sobre DICHASUS dc01, la calibración B0--B2 mejora potencia, magnitud y dispersión de
retardos, sin una mejora correspondiente de fase bruta. La comparación se sitúa en la
línea de calibración diferenciable de materiales y no identifica por sí sola parámetros
constitutivos verdaderos. Los resultados justifican evaluar expresamente qué propiedades
no quedan determinadas por el objetivo macroscópico de entrenamiento.

S4.2 caracteriza una corrección afín de fase formada por retardo efectivo e interceptos
por enlace. El retardo es relativo al modelo y a la referencia de procesamiento; no se
identifica automáticamente con STO físico. La transferencia entre antenas y frecuencias
muestra estructura compartida, con diferencias claras entre interpolación espectral y
extrapolación entre medias bandas. La alineación oráculo conserva su función diagnóstica
y no se interpreta como predicción autónoma.

S4.3 documenta estructura espectral residual y asociaciones con la distribución energética
de las trayectorias. El análisis entre posiciones, dentro de posición y mediante
remuestreo agrupado evidencia que el nivel de agregación modifica la interpretación.
La asociación observada no determina una causa única ni reemplaza un contraste explícito
de independencia que conserve la distribución marginal y el procedimiento de ajuste.

S4.4 muestra que la partición original favorece interpolación densa: la distancia mediana
al vecino de entrenamiento es 0.0483 m. La evaluación posterior mediante cinco bloques
espaciales y exclusión de 0.5 m produce un MAE medio de retardo de 27.56 ns con posición,
19.29 ns con descriptores RT y 18.69 ns con ambos. La información del trazado conserva
utilidad en regiones reservadas de dc01, aunque la mejora depende de la región y no
constituye todavía transferencia entre escenas.

S4.5 vincula el retardo predicho con métricas del canal. En interpolación densa,
el error angular mediano disminuye de 81.42 a 45.59 grados y la coherencia
aumenta de 0.156 a 0.634. Bajo evaluación por bloques,
el predictor combinado alcanza error angular mediano de 64.08 grados y coherencia
mediana de 0.448, frente a 81.49 grados y 0.157 con corrección exclusiva de CPO.
El CPO continúa ajustándose con la medición evaluada: la contribución demostrada es
la transferencia del retardo y su utilidad condicionada, no la recuperación completa
de fase sin acceso al canal. La degradación respecto de interpolación densa cuantifica
la importancia del protocolo espacial.

S4.6 amplía el diagnóstico a 640\,000 fases residuales en 10\,000 capturas.
La resultante global de 0.001088 y las medianas por antena de 0.01024 y 0.00630
no respaldan una fase absoluta fija fuertemente concentrada. La dirección común
estimada con las 32 antenas reduce el error mediano de 88.99 a 45.97 grados;
esa reducción es descriptiva y requiere controles de ajuste finito para
cuantificar su significación. La concentración marginal baja no excluye por
sí sola sesgos relativos estables entre antenas.

La recuperación temporal produjo correspondencia XYZ exacta y un manifiesto
sin tiempos faltantes. Con esa referencia, copiar la fase común del vecino
espacial alcanza errores medianos de 89.75 y 88.85 grados; copiarla del vecino
temporal produce 90.34 y 90.56 grados. La falta de una ventaja clara se observa
también en los intervalos temporales examinados. El resultado delimita el
alcance de esas reglas locales, sin demostrar independencia, equivalencia
estadística o una causa instrumental única.

La etapa siguiente S4.7a evaluará concordancia relativa respecto de fase común.
Los productos espaciales $G$ eliminan además el retardo común; por ello su
contribución predictiva debe medirse con productos entre frecuencias $K$ o con
la distancia del canal apilado respecto de una sola rotación por arreglo.
S4.7b comparará incertidumbre circular y por trayectoria cuando el residual lo justifique.
S5 evaluará representaciones con entradas, invariancias y costes declarados;
S6 contrastará tanto $\Delta H$ como la sensibilidad del observable invariante
mediante desplazamientos conocidos y calibración basal fija; S7 validará
operadores Wi-Fi y LoRa específicos. La recuperación respiratoria y la
incertidumbre de la tarea permanecen hipótesis por validar con mediciones
independientes. Una invariancia que mejora fidelidad estática puede eliminar
parte de la sensibilidad física y, por tanto, debe evaluarse frente a la tarea.

El criterio doctoral de éxito consiste en establecer qué representación y qué tratamiento
de la discrepancia conservan información útil para sensado, bajo qué condiciones esa
información se transfiere y cuándo el sistema debe reconocer que la evidencia disponible
es insuficiente. La comparación entre error de reconstrucción y error diferencial
proporciona el vínculo matemático y experimental entre los avances de canal y ese objetivo.

La propuesta actualizada convierte ese criterio en dos ramas complementarias y una
secuencia de intervenciones físicas. La rama robusta buscará estabilidad frente a
transformaciones instrumentales especificadas; la coherente conservará micromovimiento
cuando exista una referencia verificable. Su separación constituye una hipótesis a
contrastar, porque una transformación común puede contener simultáneamente deriva y
señal física. Por ello la invariancia estática y la sensibilidad diferencial se medirán
juntas, sin convertir una mejora de normalización en evidencia automática de sensado.

La geometría se tratará en tres escalas: recinto, registro de antenas y desplazamiento
local. G0--G4 permitirá medir el aporte de reconstrucción y calibración; E0--E2 establecerá
el piso instrumental y la respuesta mecánica antes de E3--E4 con personas. E5 evaluará
transferencia real entre condiciones independientes. El cronograma y la propuesta de
Ciencia de Frontera priorizan referencia RF, mecánica y trazabilidad, reutilizando el
equipo reportado. Los resultados actuales terminan en S4.6d.2; el programa restante
expresa compromisos de validación y decisiones experimentales, no logros anticipados.

\clearpage
\bibliographystyle{ieeetr}
\bibliography{references}
\end{document}
